% !TeX program = xelatex
% !TeX encoding = UTF-8
\documentclass[UTF8,10pt,a4paper,fontset=fandol]{ctexart}
\usepackage{amsmath,amssymb,graphicx,tabularx,array,multirow,enumitem,titlesec,xcolor,ragged2e,etoolbox}
\usepackage[a4paper,top=1.6cm,bottom=1.6cm,left=1.6cm,right=1.6cm]{geometry}
\setlength{\parindent}{2em}
\setlength{\parskip}{0.15em}
\setlength{\tabcolsep}{4pt}
\renewcommand{\arraystretch}{1.15}
\setlist[enumerate]{leftmargin=2.6em,itemsep=0.1em,topsep=0.2em}
\titleformat{\section}{\large\bfseries}{\thesection}{0.6em}{}
\titleformat{\subsection}{\normalsize\bfseries}{\thesubsection}{0.6em}{}
\titleformat{\subsubsection}{\normalsize\bfseries}{\thesubsubsection}{0.6em}{}
\titlespacing*{\section}{0pt}{0.7em}{0.25em}
\titlespacing*{\subsection}{0pt}{0.55em}{0.2em}
\titlespacing*{\subsubsection}{0pt}{0.4em}{0.15em}
\raggedbottom
\BeforeBeginEnvironment{tabularx}{\par\noindent}
\newcommand{\fieldvalue}[1]{#1}
\newcommand{\handwritten}[1]{#1}
\newcommand{\TODO}[1]{#1}
\newcommand{\checkboxfield}[1]{\ifstrequal{#1}{checked}{☑}{\ifstrequal{#1}{unclear}{?}{☐}}}
\begin{document}

% TODO: The supplied page image is upside down and several handwritten entries are
% difficult to read. The following transcription preserves the visible form layout.

\begin{center}
{\Large\textbf{普通高等学校毕业生就业推荐表}}
\end{center}

\begin{center}
\small
学校名称：\underline{\hspace{7em}}
\qquad
院（系）：\underline{\hspace{7em}}
\qquad
专业：\underline{\hspace{7em}}
\end{center}

\vspace{0.4em}

\begin{tabularx}{\textwidth}{|p{3.5em}|X|p{4.5em}|X|}
\hline
姓名 & 
% #VALUE_ID: LEX-P0001-V0001
% #FIELD_VALUE: 姓名
% TODO #HANDWRITTEN: 难以辨认; 蓝色手写姓名
\fieldvalue{\handwritten{（难以辨认）}} &
性别 & \rule{0pt}{1.8em} \\ \hline
出生年月 &
% #VALUE_ID: LEX-P0001-V0002
% #FIELD_VALUE: 出生年月
% TODO #HANDWRITTEN: 2025年06月18日（疑似）; 字迹及字段对应关系不清
\fieldvalue{\handwritten{2025年06月18日（疑似）}} &
政治面貌 & \\ \hline
培养方式 & 
% #VALUE_ID: LEX-P0001-V0003
% #FIELD_VALUE: 培养方式
% TODO #HANDWRITTEN: 难以辨认; 蓝色手写内容
\fieldvalue{\handwritten{（难以辨认）}} &
健康状况 & \\ \hline
生源地 & & 学历 & \\ \hline
学制 & & 专业 & \\ \hline
毕业时间 &
% #VALUE_ID: LEX-P0001-V0004
% #FIELD_VALUE: 毕业时间
% TODO #HANDWRITTEN: 2025年06月18日（疑似）; 蓝色手写日期
\fieldvalue{\handwritten{2025年06月18日（疑似）}} &
外语水平 & \\ \hline
计算机水平 & & 联系方式 &
% #VALUE_ID: LEX-P0001-V0005
% #FIELD_VALUE: 联系方式
% TODO #HANDWRITTEN: 难以辨认; 蓝色手写数字
\fieldvalue{\handwritten{（难以辨认）}} \\ \hline
\end{tabularx}

\vspace{0.5em}

\begin{tabularx}{\textwidth}{|p{7em}|X|}
\hline
奖惩情况 &
\rule{0pt}{3.2em} \\ \hline
\end{tabularx}

\vspace{0.3em}

\begin{tabularx}{\textwidth}{|p{7em}|X|}
\hline
\multirow{4}{*}{学习成绩}
&
{\begin{tabularx}{\linewidth}{@{}p{11em}|X@{}}
课程成绩（或平均成绩） & \\
\cline{1-2}
必修课平均成绩 &
% #VALUE_ID: LEX-P0001-V0006
% #FIELD_VALUE: 必修课平均成绩
% TODO #HANDWRITTEN: 难以辨认; 蓝色手写数字
\fieldvalue{\handwritten{（难以辨认）}} \\
\cline{1-2}
选修课平均成绩 &
% #VALUE_ID: LEX-P0001-V0007
% #FIELD_VALUE: 选修课平均成绩
% TODO #HANDWRITTEN: 难以辨认; 蓝色手写数字
\fieldvalue{\handwritten{（难以辨认）}} \\
\cline{1-2}
综合测评成绩 &
% #VALUE_ID: LEX-P0001-V0008
% #FIELD_VALUE: 综合测评成绩
% TODO #HANDWRITTEN: 难以辨认; 蓝色手写数字
\fieldvalue{\handwritten{（难以辨认）}} \\
\end{tabularx}}
\\ \hline
\end{tabularx}

\vspace{0.3em}

\begin{tabularx}{\textwidth}{|p{7em}|X|}
\hline
自我鉴定 &
\rule{0pt}{5em} \\ \hline
学校意见 &
\rule{0pt}{4em} \\ \hline
\end{tabularx}

\vspace{0.35em}

\noindent
联系人：\underline{\hspace{8em}}
\qquad
联系电话：
% #VALUE_ID: LEX-P0001-V0009
% #FIELD_VALUE: 联系电话
% TODO #HANDWRITTEN: 难以辨认; 页面底部蓝色手写电话号码
\fieldvalue{\handwritten{（难以辨认）}}
\qquad
日期：\underline{\hspace{7em}}

% TODO: A red official seal and a printed page number are visible on the page.
% They are fixed boilerplate/seal content and are not treated as editable values.

% LEXOID_PAGE_COMPLETED: 1/78

\newpage

% TODO: The page is scanned upside down and several Chinese characters are low-resolution; transcriptions marked as uncertain should be reviewed against the source PDF.

\begin{center}
\textbf{工程质量验收相关记录}
\end{center}

\begin{tabularx}{\textwidth}{|>{\RaggedRight\arraybackslash}p{0.20\textwidth}|X|>{\RaggedRight\arraybackslash}p{0.20\textwidth}|X|}
\hline
工程名称 & 
% #VALUE_ID: LEX-P0002-V0001
% #FIELD_VALUE: 工程名称
% #TODO #HANDWRITTEN: （疑似工程名称）; 字迹及扫描方向不清
\fieldvalue{\handwritten{（疑似工程名称）}} &
单位工程名称 & \\ \hline
施工单位 & & 监理单位 & \\ \hline
项目负责人 & & 记录编号 & \\ \hline
\end{tabularx}

\vspace{1em}

\section{8.1 施工单位工程质量检查评定记录}

\begin{tabularx}{\textwidth}{|>{\RaggedRight\arraybackslash}p{0.58\textwidth}|>{\RaggedRight\arraybackslash}p{0.22\textwidth}|X|}
\hline
检查项目 & 自评 & 核查结果 \\ \hline
工程资料及质量保证资料检查 & 
% #VALUE_ID: LEX-P0002-V0002
% #FIELD_VALUE: 工程资料及质量保证资料检查
% #TODO #HANDWRITTEN: 合格; 蓝色手写内容部分难以辨认
\fieldvalue{\handwritten{合格}} & \\ \hline
建筑工程施工质量验收统一标准及相关规范执行情况 & 
% #VALUE_ID: LEX-P0002-V0003
% #FIELD_VALUE: 建筑工程施工质量验收统一标准及相关规范执行情况
% #HANDWRITTEN: 合格
\fieldvalue{\handwritten{合格}} & \\ \hline
分部、分项工程质量验收情况 & 
% #VALUE_ID: LEX-P0002-V0004
% #FIELD_VALUE: 分部、分项工程质量验收情况
% #HANDWRITTEN: 合格
\fieldvalue{\handwritten{合格}} & \\ \hline
观感质量检查情况 & 
% #VALUE_ID: LEX-P0002-V0005
% #FIELD_VALUE: 观感质量检查情况
% #HANDWRITTEN: 合格
\fieldvalue{\handwritten{合格}} & \\ \hline
主要使用功能抽查情况 & 
% #VALUE_ID: LEX-P0002-V0006
% #FIELD_VALUE: 主要使用功能抽查情况
% #HANDWRITTEN: 合格
\fieldvalue{\handwritten{合格}} & \\ \hline
\end{tabularx}

\vspace{0.5em}

\begin{tabularx}{\textwidth}{|>{\RaggedRight\arraybackslash}p{0.58\textwidth}|>{\RaggedRight\arraybackslash}p{0.22\textwidth}|X|}
\hline
施工单位自评结论 & 
% #VALUE_ID: LEX-P0002-V0007
% #FIELD_VALUE: 施工单位自评结论
% #HANDWRITTEN: 合格
\fieldvalue{\handwritten{合格}} & \\ \hline
项目负责人签字 & 
% #VALUE_ID: LEX-P0002-V0008
% #FIELD_VALUE: 项目负责人签字
% #TODO #HANDWRITTEN: （签名）; 手写签名无法辨认
\fieldvalue{\handwritten{（签名）}} & \\ \hline
施工单位技术负责人签字 & 
% #VALUE_ID: LEX-P0002-V0009
% #FIELD_VALUE: 施工单位技术负责人签字
% #TODO #HANDWRITTEN: （签名）; 手写签名无法辨认
\fieldvalue{\handwritten{（签名）}} & \\ \hline
日期 & 
% #VALUE_ID: LEX-P0002-V0010
% #FIELD_VALUE: 日期
% #HANDWRITTEN: 2025.05.16
\fieldvalue{\handwritten{2025.05.16}} & \\ \hline
\end{tabularx}

\vspace{1em}

\section{8.2 监理单位工程质量评估报告}

\begin{tabularx}{\textwidth}{|>{\RaggedRight\arraybackslash}p{0.58\textwidth}|>{\RaggedRight\arraybackslash}p{0.22\textwidth}|X|}
\hline
监理单位评估意见 & 
% #VALUE_ID: LEX-P0002-V0011
% #FIELD_VALUE: 监理单位评估意见
% #TODO #HANDWRITTEN: 合格; 蓝色手写内容部分难以辨认
\fieldvalue{\handwritten{合格}} & \\ \hline
总监理工程师签字 & 
% #VALUE_ID: LEX-P0002-V0012
% #FIELD_VALUE: 总监理工程师签字
% #TODO #HANDWRITTEN: （签名）; 手写签名无法辨认
\fieldvalue{\handwritten{（签名）}} & \\ \hline
日期 & 
% #VALUE_ID: LEX-P0002-V0013
% #FIELD_VALUE: 日期
% #HANDWRITTEN: 2025.05.18
\fieldvalue{\handwritten{2025.05.18}} & \\ \hline
\end{tabularx}

\vspace{1em}

\begin{tabularx}{\textwidth}{|>{\RaggedRight\arraybackslash}p{0.20\textwidth}|X|>{\RaggedRight\arraybackslash}p{0.20\textwidth}|X|}
\hline
验收结论 & 
% #VALUE_ID: LEX-P0002-V0014
% #FIELD_VALUE: 验收结论
% #TODO #HANDWRITTEN: 合格; 蓝色手写内容部分难以辨认
\fieldvalue{\handwritten{合格}} &
验收日期 & \\ \hline
\end{tabularx}

% TODO: The lower portion of the page contains additional tabular/header material visible only partially in the scan; continuation and exact printed wording require review.
% LEXOID_PAGE_COMPLETED: 2/78

\newpage

\begin{center}
\textbf{药品及相关记录表}
\end{center}

\noindent
\textbf{登记日期：}
% #VALUE_ID: LEX-P0003-V0001
% #FIELD_VALUE: 登记日期
\fieldvalue{2020-01-02}
\hfill
\textbf{编号：}
% #VALUE_ID: LEX-P0003-V0002
% #FIELD_VALUE: 编号
\fieldvalue{200--408}

\begin{tabularx}{\textwidth}{|c|X|X|X|X|X|c|}
\hline
序号 & 药品名称 & 规格 & 剂型 & 批号 & 生产厂家 & 有效期 \\
\hline
2 &
% #VALUE_ID: LEX-P0003-V0003
% #FIELD_VALUE: 药品名称（第2行）
\fieldvalue{维生素B1} &
% #VALUE_ID: LEX-P0003-V0004
% #FIELD_VALUE: 规格（第2行）
\fieldvalue{100\,mg} &
% #VALUE_ID: LEX-P0003-V0005
% #FIELD_VALUE: 剂型（第2行）
\fieldvalue{片剂} &
% #VALUE_ID: LEX-P0003-V0006
% #FIELD_VALUE: 批号（第2行）
\fieldvalue{--} &
% #VALUE_ID: LEX-P0003-V0007
% #FIELD_VALUE: 生产厂家（第2行）
\fieldvalue{Bioscience CRA} &
% #VALUE_ID: LEX-P0003-V0008
% #FIELD_VALUE: 有效期（第2行）
\fieldvalue{11/2021} \\
\hline
3 &
% #VALUE_ID: LEX-P0003-V0009
% #FIELD_VALUE: 药品名称（第3行）
\fieldvalue{复方氨基酸注射液} &
% #VALUE_ID: LEX-P0003-V0010
% #FIELD_VALUE: 规格（第3行）
\fieldvalue{250\,mL} &
% #VALUE_ID: LEX-P0003-V0011
% #FIELD_VALUE: 剂型（第3行）
\fieldvalue{注射液} &
% #VALUE_ID: LEX-P0003-V0012
% #FIELD_VALUE: 批号（第3行）
\fieldvalue{--} &
% #VALUE_ID: LEX-P0003-V0013
% #FIELD_VALUE: 生产厂家（第3行）
\fieldvalue{Gentech Exa Pro} &
% #VALUE_ID: LEX-P0003-V0014
% #FIELD_VALUE: 有效期（第3行）
\fieldvalue{11/2021} \\
\hline
4 &
% #VALUE_ID: LEX-P0003-V0015
% #FIELD_VALUE: 药品名称（第4行）
\fieldvalue{复方氨基酸注射液} &
% #VALUE_ID: LEX-P0003-V0016
% #FIELD_VALUE: 规格（第4行）
\fieldvalue{250\,mL} &
% #VALUE_ID: LEX-P0003-V0017
% #FIELD_VALUE: 剂型（第4行）
\fieldvalue{注射液} &
% #VALUE_ID: LEX-P0003-V0018
% #FIELD_VALUE: 批号（第4行）
\fieldvalue{--} &
% #VALUE_ID: LEX-P0003-V0019
% #FIELD_VALUE: 生产厂家（第4行）
\fieldvalue{Gentech Exa Pro} &
% #VALUE_ID: LEX-P0003-V0020
% #FIELD_VALUE: 有效期（第4行）
\fieldvalue{11/2021} \\
\hline
5 &
% #VALUE_ID: LEX-P0003-V0021
% #FIELD_VALUE: 药品名称（第5行）
\fieldvalue{复方氨基酸注射液} &
% #VALUE_ID: LEX-P0003-V0022
% #FIELD_VALUE: 规格（第5行）
\fieldvalue{250\,mL} &
% #VALUE_ID: LEX-P0003-V0023
% #FIELD_VALUE: 剂型（第5行）
\fieldvalue{注射液} &
% #VALUE_ID: LEX-P0003-V0024
% #FIELD_VALUE: 批号（第5行）
\fieldvalue{--} &
% #VALUE_ID: LEX-P0003-V0025
% #FIELD_VALUE: 生产厂家（第5行）
\fieldvalue{Gentech Exa Pro} &
% #VALUE_ID: LEX-P0003-V0026
% #FIELD_VALUE: 有效期（第5行）
\fieldvalue{11/2021} \\
\hline
6 &
% #VALUE_ID: LEX-P0003-V0027
% #FIELD_VALUE: 药品名称（第6行）
\fieldvalue{复方氨基酸注射液} &
% #VALUE_ID: LEX-P0003-V0028
% #FIELD_VALUE: 规格（第6行）
\fieldvalue{250\,mL} &
% #VALUE_ID: LEX-P0003-V0029
% #FIELD_VALUE: 剂型（第6行）
\fieldvalue{注射液} &
% #VALUE_ID: LEX-P0003-V0030
% #FIELD_VALUE: 批号（第6行）
\fieldvalue{--} &
% #VALUE_ID: LEX-P0003-V0031
% #FIELD_VALUE: 生产厂家（第6行）
\fieldvalue{Gentech Exa Pro} &
% #VALUE_ID: LEX-P0003-V0032
% #FIELD_VALUE: 有效期（第6行）
\fieldvalue{11/2021} \\
\hline
7 &
% #VALUE_ID: LEX-P0003-V0033
% #FIELD_VALUE: 药品名称（第7行）
\fieldvalue{复方氨基酸注射液} &
% #VALUE_ID: LEX-P0003-V0034
% #FIELD_VALUE: 规格（第7行）
\fieldvalue{250\,mL} &
% #VALUE_ID: LEX-P0003-V0035
% #FIELD_VALUE: 剂型（第7行）
\fieldvalue{注射液} &
% #VALUE_ID: LEX-P0003-V0036
% #FIELD_VALUE: 批号（第7行）
\fieldvalue{--} &
% #VALUE_ID: LEX-P0003-V0037
% #FIELD_VALUE: 生产厂家（第7行）
\fieldvalue{Gentech Exa Pro} &
% #VALUE_ID: LEX-P0003-V0038
% #FIELD_VALUE: 有效期（第7行）
\fieldvalue{11/2021} \\
\hline
\end{tabularx}

\vspace{1em}

\noindent
\textbf{备注：}
% #VALUE_ID: LEX-P0003-V0039
% #FIELD_VALUE: 备注
% #TODO #HANDWRITTEN: 字迹不清；蓝色手写备注
\fieldvalue{\handwritten{字迹不清}}

\begin{tabularx}{\textwidth}{|c|X|X|X|c|}
\hline
序号 & 项目名称 & 规格／型号 & 生产厂家 & 备注 \\
\hline
1 &
% #VALUE_ID: LEX-P0003-V0040
% #FIELD_VALUE: 项目名称（第1行）
\fieldvalue{一次性使用输液器} &
% #VALUE_ID: LEX-P0003-V0041
% #FIELD_VALUE: 规格／型号（第1行）
\fieldvalue{--} &
% #VALUE_ID: LEX-P0003-V0042
% #FIELD_VALUE: 生产厂家（第1行）
\fieldvalue{--} &
% #VALUE_ID: LEX-P0003-V0043
% #FIELD_VALUE: 备注（第1行）
% #TODO #HANDWRITTEN: 2；蓝色手写数字
\fieldvalue{\handwritten{2}} \\
\hline
2 &
% #VALUE_ID: LEX-P0003-V0044
% #FIELD_VALUE: 项目名称（第2行）
\fieldvalue{一次性使用无菌注射器} &
% #VALUE_ID: LEX-P0003-V0045
% #FIELD_VALUE: 规格／型号（第2行）
\fieldvalue{--} &
% #VALUE_ID: LEX-P0003-V0046
% #FIELD_VALUE: 生产厂家（第2行）
\fieldvalue{Bioscience CRA} &
% #VALUE_ID: LEX-P0003-V0047
% #FIELD_VALUE: 备注（第2行）
% #TODO #HANDWRITTEN: 11/04；蓝色手写日期
\fieldvalue{\handwritten{11/04}} \\
\hline
3 &
% #VALUE_ID: LEX-P0003-V0048
% #FIELD_VALUE: 项目名称（第3行）
\fieldvalue{一次性使用无菌注射器} &
% #VALUE_ID: LEX-P0003-V0049
% #FIELD_VALUE: 规格／型号（第3行）
\fieldvalue{--} &
% #VALUE_ID: LEX-P0003-V0050
% #FIELD_VALUE: 生产厂家（第3行）
\fieldvalue{Gentech Exa Pro} &
% #VALUE_ID: LEX-P0003-V0051
% #FIELD_VALUE: 备注（第3行）
% #TODO #HANDWRITTEN: 11/05；蓝色手写日期
\fieldvalue{\handwritten{11/05}} \\
\hline
4 &
% #VALUE_ID: LEX-P0003-V0052
% #FIELD_VALUE: 项目名称（第4行）
\fieldvalue{一次性使用无菌注射器} &
% #VALUE_ID: LEX-P0003-V0053
% #FIELD_VALUE: 规格／型号（第4行）
\fieldvalue{--} &
% #VALUE_ID: LEX-P0003-V0054
% #FIELD_VALUE: 生产厂家（第4行）
\fieldvalue{Gentech Exa Pro} &
% #VALUE_ID: LEX-P0003-V0055
% #FIELD_VALUE: 备注（第4行）
% #TODO #HANDWRITTEN: 11/06；蓝色手写日期
\fieldvalue{\handwritten{11/06}} \\
\hline
5 &
% #VALUE_ID: LEX-P0003-V0056
% #FIELD_VALUE: 项目名称（第5行）
\fieldvalue{一次性使用无菌注射器} &
% #VALUE_ID: LEX-P0003-V0057
% #FIELD_VALUE: 规格／型号（第5行）
\fieldvalue{--} &
% #VALUE_ID: LEX-P0003-V0058
% #FIELD_VALUE: 生产厂家（第5行）
\fieldvalue{Gentech Exa Pro} &
% #VALUE_ID: LEX-P0003-V0059
% #FIELD_VALUE: 备注（第5行）
% #TODO #HANDWRITTEN: 11/07；蓝色手写日期
\fieldvalue{\handwritten{11/07}} \\
\hline
6 &
% #VALUE_ID: LEX-P0003-V0060
% #FIELD_VALUE: 项目名称（第6行）
\fieldvalue{一次性使用无菌注射器} &
% #VALUE_ID: LEX-P0003-V0061
% #FIELD_VALUE: 规格／型号（第6行）
\fieldvalue{--} &
% #VALUE_ID: LEX-P0003-V0062
% #FIELD_VALUE: 生产厂家（第6行）
\fieldvalue{Gentech Exa Pro} &
% #VALUE_ID: LEX-P0003-V0063
% #FIELD_VALUE: 备注（第6行）
% #TODO #HANDWRITTEN: 11/08；蓝色手写日期
\fieldvalue{\handwritten{11/08}} \\
\hline
7 &
% #VALUE_ID: LEX-P0003-V0064
% #FIELD_VALUE: 项目名称（第7行）
\fieldvalue{一次性使用无菌注射器} &
% #VALUE_ID: LEX-P0003-V0065
% #FIELD_VALUE: 规格／型号（第7行）
\fieldvalue{--} &
% #VALUE_ID: LEX-P0003-V0066
% #FIELD_VALUE: 生产厂家（第7行）
\fieldvalue{Gentech Exa Pro} &
% #VALUE_ID: LEX-P0003-V0067
% #FIELD_VALUE: 备注（第7行）
% #TODO #HANDWRITTEN: 11/09；蓝色手写日期
\fieldvalue{\handwritten{11/09}} \\
\hline
\end{tabularx}

% TODO: The page contains additional rotated tabular content and handwritten annotations; exact text is illegible at the supplied resolution.
% LEXOID_PAGE_COMPLETED: 3/78

\newpage

\section*{职业病危害因素检测结果}

\noindent
报告编号：WZ-...；检测日期：2015-04-18

% TODO: 部分页眉文字及编号因页面扫描方向和分辨率较低，无法完全辨认。

\begin{tabularx}{\textwidth}{|c|X|c|c|}
\hline
\textbf{检测点编号} & \textbf{检测项目} & \textbf{检测结果} & \textbf{单位} \\
\hline
1 & 粉尘（呼吸性粉尘） & 
% #VALUE_ID: LEX-P0004-V0001
% #FIELD_VALUE: 检测结果
% #HANDWRITTEN: 0.026
\fieldvalue{\handwritten{0.026}} &
% #VALUE_ID: LEX-P0004-V0002
% #FIELD_VALUE: 单位
% #HANDWRITTEN: 2
\fieldvalue{\handwritten{2}} \\
\hline
2 & 粉尘（总尘） & 
% #VALUE_ID: LEX-P0004-V0003
% #FIELD_VALUE: 检测结果
% #HANDWRITTEN: 0.051
\fieldvalue{\handwritten{0.051}} &
% #VALUE_ID: LEX-P0004-V0004
% #FIELD_VALUE: 单位
% #HANDWRITTEN: 2
\fieldvalue{\handwritten{2}} \\
\hline
3 & 粉尘（呼吸性粉尘） & 
% #VALUE_ID: LEX-P0004-V0005
% #FIELD_VALUE: 检测结果
% #HANDWRITTEN: 0.016
\fieldvalue{\handwritten{0.016}} &
% #VALUE_ID: LEX-P0004-V0006
% #FIELD_VALUE: 单位
% #HANDWRITTEN: 2
\fieldvalue{\handwritten{2}} \\
\hline
4 & & & \\
\hline
5 & & & \\
\hline
\end{tabularx}

\vspace{1em}

% #VALUE_ID: LEX-P0004-V0007
% #FIELD_VALUE: 采样日期
% #HANDWRITTEN: 2015.04.18
采样日期：\fieldvalue{\handwritten{2015.04.18}}

\hspace{3em}
% #VALUE_ID: LEX-P0004-V0008
% #FIELD_VALUE: 检测日期
% #HANDWRITTEN: 2015.04.18
检测日期：\fieldvalue{\handwritten{2015.04.18}}

\section*{检测结果及判定}

\begin{tabularx}{\textwidth}{|c|X|c|c|}
\hline
\textbf{序号} & \textbf{检测项目} & \multicolumn{2}{c|}{\textbf{判定}} \\
\cline{3-4}
 & & \textbf{是} & \textbf{否} \\
\hline
(1) & 粉尘中游离二氧化硅含量 & 
% #VALUE_ID: LEX-P0004-V0009
% #FIELD_VALUE: （1）判定
% #HANDWRITTEN: ✓
\fieldvalue{\handwritten{\checkmark}} & $\square$ \\
\hline
(2) & 粉尘分散度 & 
% #VALUE_ID: LEX-P0004-V0010
% #FIELD_VALUE: （2）判定
% #HANDWRITTEN: ✓
\fieldvalue{\handwritten{\checkmark}} & $\square$ \\
\hline
(3) & 粉尘检测结果是否符合要求 & 
% #VALUE_ID: LEX-P0004-V0011
% #FIELD_VALUE: （3）判定
% #HANDWRITTEN: ✓
\fieldvalue{\handwritten{\checkmark}} & $\square$ \\
\hline
(4) & 作业场所空气中有害物质浓度 & 
% #VALUE_ID: LEX-P0004-V0012
% #FIELD_VALUE: （4）判定
% #HANDWRITTEN: ✓
\fieldvalue{\handwritten{\checkmark}} & $\square$ \\
\hline
(5) & 其他检测项目 & 
% #VALUE_ID: LEX-P0004-V0013
% #FIELD_VALUE: （5）判定
% #HANDWRITTEN: ✓
\fieldvalue{\handwritten{\checkmark}} & $\square$ \\
\hline
(6) & 检测结果评价 & 
% #VALUE_ID: LEX-P0004-V0014
% #FIELD_VALUE: （6）判定
% #HANDWRITTEN: ✓
\fieldvalue{\handwritten{\checkmark}} & $\square$ \\
\hline
(7) & 检测结论 & 
% #VALUE_ID: LEX-P0004-V0015
% #FIELD_VALUE: （7）判定
% #HANDWRITTEN: ✓
\fieldvalue{\handwritten{\checkmark}} & $\square$ \\
\hline
\end{tabularx}

\vspace{1em}

\noindent
备注：

% #VALUE_ID: LEX-P0004-V0016
% #FIELD_VALUE: 备注
% #HANDWRITTEN: 1625060
\fieldvalue{\handwritten{1625060}}

% TODO: 页面中部分手写数字及检测项目名称无法可靠辨认，以上以可识别内容转录；表格在下一页可能继续。
% LEXOID_PAGE_COMPLETED: 4/78

\newpage

\begin{flushright}
\textbf{令和7年度　確認票}
\end{flushright}

\begin{flushright}
% #VALUE_ID: LEX-P0005-V0001
% #FIELD_VALUE: 日付
% #TODO #HANDWRITTEN: 令和7年?月?日; 赤字の手書き文字が不鮮明
\fieldvalue{\handwritten{令和7年?月?日}}
\end{flushright}

\begin{tabularx}{\textwidth}{|c|c|X|}
\hline
\textbf{番号} & \textbf{回答} & \textbf{確認事項} \\
\hline
 & 
% #VALUE_ID: LEX-P0005-V0002
% #FIELD_VALUE: 管理番号
% #TODO #HANDWRITTEN: 232?25?006; 青字の手書き番号が不鮮明
\fieldvalue{\handwritten{232?25?006}}
&
\textbf{基本事項} \\
\hline
(8) &
% #VALUE_ID: LEX-P0005-V0003
% #FIELD_VALUE: 回答（有／無）
% #HANDWRITTEN: ✓
\fieldvalue{\handwritten{✓}}
&
（確認事項の本文は判読困難） \\
\hline
(9) &
% #VALUE_ID: LEX-P0005-V0004
% #FIELD_VALUE: 回答（有／無）
% #HANDWRITTEN: ✓
\fieldvalue{\handwritten{✓}}
&
（確認事項の本文は判読困難） \\
\hline
(10) &
% #VALUE_ID: LEX-P0005-V0005
% #FIELD_VALUE: 回答（数値）
% #TODO #HANDWRITTEN: 16 / 20; 青字の手書き数値
\fieldvalue{\handwritten{16 / 20}}
&
（確認事項の本文は判読困難）
\\[-1mm]
&
% #VALUE_ID: LEX-P0005-V0006
% #FIELD_VALUE: 回答（有／無）
% #HANDWRITTEN: ✓
\fieldvalue{\handwritten{✓}}
&
\\
\hline
(11) &
% #VALUE_ID: LEX-P0005-V0007
% #FIELD_VALUE: 回答（有／無）
% #HANDWRITTEN: ✓
\fieldvalue{\handwritten{✓}}
&
（確認事項の本文は判読困難） \\
\hline
(12) &
% #VALUE_ID: LEX-P0005-V0008
% #FIELD_VALUE: 回答（有／無）
% #HANDWRITTEN: ✓
\fieldvalue{\handwritten{✓}}
&
（確認事項の本文は判読困難） \\
\hline
(13) &
% #VALUE_ID: LEX-P0005-V0009
% #FIELD_VALUE: 回答（有／無）
% #HANDWRITTEN: ✓
\fieldvalue{\handwritten{✓}}
&
（確認事項の本文は判読困難） \\
\hline
(14) &
% #VALUE_ID: LEX-P0005-V0010
% #FIELD_VALUE: 回答（有／無）
% #HANDWRITTEN: ✓
\fieldvalue{\handwritten{✓}}
&
（確認事項の本文は判読困難） \\
\hline
(15) &
% #VALUE_ID: LEX-P0005-V0011
% #FIELD_VALUE: 回答（有／無）
% #HANDWRITTEN: ✓
\fieldvalue{\handwritten{✓}}
&
（確認事項の本文は判読困難） \\
\hline
(16) &
% #VALUE_ID: LEX-P0005-V0012
% #FIELD_VALUE: 回答（有／無）
% #HANDWRITTEN: ✓
\fieldvalue{\handwritten{✓}}
&
（確認事項の本文は判読困難） \\
\hline
(17) &
% #VALUE_ID: LEX-P0005-V0013
% #FIELD_VALUE: 回答（有／無）
% #HANDWRITTEN: ✓
\fieldvalue{\handwritten{✓}}
&
（確認事項の本文は判読困難） \\
\hline
\end{tabularx}

% TODO: The page is a rotated scan. Several Japanese question texts and the red date entry are not sufficiently legible for reliable transcription.
% LEXOID_PAGE_COMPLETED: 5/78

\newpage

\begin{center}
\textbf{医疗机构制剂调剂使用申请表}
\end{center}

\noindent
\begin{tabularx}{\textwidth}{|p{0.18\textwidth}|X|p{0.16\textwidth}|}
\hline
\textbf{类别} & \textbf{品名及规格} & \textbf{数量} \\
\hline
\multicolumn{3}{|l|}{\textbf{申请调剂使用的制剂信息}} \\
\hline
\end{tabularx}

\vspace{2mm}

\small
\begin{tabularx}{\textwidth}{|p{0.14\textwidth}|X|p{0.13\textwidth}|}
\hline
\textbf{序号} & \textbf{申请理由及相关资料} & \textbf{是否符合} \\
\hline
1 &
% TODO: The printed text in this row is partially illegible in the scan.
申请调剂使用的制剂及相关资料。 &
% #VALUE_ID: LEX-P0006-V0001
% #FIELD_VALUE: 第1项是否符合
% #HANDWRITTEN: √
\fieldvalue{\handwritten{√}} \\
\hline
2 &
% TODO: The printed text in this row is partially illegible in the scan.
制剂质量标准、检验报告及相关证明材料。 &
% #VALUE_ID: LEX-P0006-V0002
% #FIELD_VALUE: 第2项是否符合
% #HANDWRITTEN: √
\fieldvalue{\handwritten{√}} \\
\hline
3 &
% TODO: The printed text in this row is partially illegible in the scan.
医疗机构制剂注册、备案及调剂使用相关资料。 &
% #VALUE_ID: LEX-P0006-V0003
% #FIELD_VALUE: 第3项是否符合
% #HANDWRITTEN: √
\fieldvalue{\handwritten{√}} \\
\hline
4 &
% TODO: The printed text in this row is partially illegible in the scan.
制剂说明书及包装、标签等资料。 &
% #VALUE_ID: LEX-P0006-V0004
% #FIELD_VALUE: 第4项是否符合
% #HANDWRITTEN: √
\fieldvalue{\handwritten{√}} \\
\hline
5 &
% TODO: The printed text in this row is partially illegible in the scan.
制剂调剂使用的相关要求。 &
未见填写 \\
\hline
6 &
% TODO: The printed text in this row is partially illegible in the scan.
其他需要说明的事项。 &
% #VALUE_ID: LEX-P0006-V0005
% #FIELD_VALUE: 第6项是否符合
% #HANDWRITTEN: √
\fieldvalue{\handwritten{√}} \\
\hline
7 &
% TODO: The printed text in this row is partially illegible in the scan.
审核意见及结论。 &
% #VALUE_ID: LEX-P0006-V0006
% #FIELD_VALUE: 第7项是否符合
% #HANDWRITTEN: √
\fieldvalue{\handwritten{√}} \\
\hline
\end{tabularx}
\normalsize

\vspace{3mm}

\noindent
\begin{tabularx}{\textwidth}{|p{0.18\textwidth}|X|p{0.18\textwidth}|}
\hline
\textbf{备注} & & \textbf{日期} \\
\hline
\multicolumn{3}{|l|}{\textit{（本页表格内容在此处结束；如有续页，续页内容未显示。）}} \\
\hline
\end{tabularx}

% TODO: A red official stamp and a printed institutional emblem are visible near the top of the page; they are fixed marks rather than editable form values.
% LEXOID_PAGE_COMPLETED: 6/78

\newpage

\begin{center}
\textbf{設備規格及檢查資料}
\end{center}

\begin{tabularx}{\textwidth}{|l|X|l|l|}
\hline
品名 & 規格 & 數量 & 備註 \\
\hline
% #VALUE_ID: LEX-P0007-V0001
% #FIELD_VALUE: 備註
% #TODO #HANDWRITTEN: （字跡難以辨識）; 藍色手寫內容不清楚
\fieldvalue{\handwritten{（字跡難以辨識）}} &
\multicolumn{2}{l|}{200--480 請參閱相關規格} \\
\hline
\end{tabularx}

\vspace{1em}

\begin{tabularx}{\textwidth}{|l|X|c|}
\hline
項次 & 檢查項目 & 結果 \\
\hline
1 & 設備外觀及基本資料檢查 & \\
\hline
2 & 相關文件及資料是否齊全 & 
% #VALUE_ID: LEX-P0007-V0002
% #FIELD_VALUE: 相關文件及資料是否齊全
% #HANDWRITTEN: 勾
\fieldvalue{\handwritten{勾}} \\
\hline
3 & 設備操作及功能檢查 & \\
\hline
4 & 安全裝置及防護設施檢查 & \\
\hline
5 & 測試結果是否符合規定 & 
% #VALUE_ID: LEX-P0007-V0003
% #FIELD_VALUE: 測試結果是否符合規定
% #HANDWRITTEN: 勾
\fieldvalue{\handwritten{勾}} \\
\hline
6 & 其他應檢附之資料是否齊全 & 
% #VALUE_ID: LEX-P0007-V0004
% #FIELD_VALUE: 其他應檢附之資料是否齊全
% #TODO #HANDWRITTEN: （數字或符號難以辨識）; 藍色手寫內容不清楚
\fieldvalue{\handwritten{（難以辨識）}} \\
\hline
\end{tabularx}

\vspace{1em}

\begin{tabularx}{\textwidth}{|l|X|c|}
\hline
(2) & 相關設備資料是否已完成確認 & 
% #VALUE_ID: LEX-P0007-V0005
% #FIELD_VALUE: 相關設備資料是否已完成確認
% #HANDWRITTEN: 勾
\fieldvalue{\handwritten{勾}} \\
\hline
(5) & 是否符合相關規範及要求 & 
% #VALUE_ID: LEX-P0007-V0006
% #FIELD_VALUE: 是否符合相關規範及要求
% #HANDWRITTEN: 勾
\fieldvalue{\handwritten{勾}} \\
\hline
(6) & 其他備註事項 & 
% #VALUE_ID: LEX-P0007-V0007
% #FIELD_VALUE: 其他備註事項
% #HANDWRITTEN: 勾
\fieldvalue{\handwritten{勾}} \\
\hline
\end{tabularx}

\vspace{1em}

\begin{tabularx}{\textwidth}{|l|r|}
\hline
規格項目 & 數值 \\
\hline
尺寸 & 200 \\
\hline
長度 & 2200 \\
\hline
數量 & 4 \\
\hline
寬度 & 400 \\
\hline
高度 & 150 \\
\hline
運轉時間 & 180 \\
\hline
最高轉速 & 800 \\
\hline
最低轉速 & 300 \\
\hline
時間 & 60 \\
\hline
\end{tabularx}

% TODO: The page is rotated in the source image; several Chinese labels and the blue handwritten entries are partially illegible.
% LEXOID_PAGE_COMPLETED: 7/78

\newpage

\section*{建筑起重机械检验记录（表二）}

\noindent
\begin{tabularx}{\textwidth}{|l|X|l|X|}
\hline
设备名称 & 
% #VALUE_ID: LEX-P0008-V0001
% #FIELD_VALUE: 设备名称
% #TODO #HANDWRITTEN: [无法辨认]; 字迹及图像分辨率不足
\fieldvalue{\handwritten{[无法辨认]}} &
规格型号 &
% #VALUE_ID: LEX-P0008-V0002
% #FIELD_VALUE: 规格型号
% #TODO #HANDWRITTEN: [无法辨认]; 字迹及图像分辨率不足
\fieldvalue{\handwritten{[无法辨认]}} \\
\hline
制造厂家 &
% #VALUE_ID: LEX-P0008-V0003
% #FIELD_VALUE: 制造厂家
% #TODO #HANDWRITTEN: [无法辨认]; 字迹及图像分辨率不足
\fieldvalue{\handwritten{[无法辨认]}} &
出厂日期 &
% #VALUE_ID: LEX-P0008-V0004
% #FIELD_VALUE: 出厂日期
% #TODO #HANDWRITTEN: [无法辨认]; 字迹及图像分辨率不足
\fieldvalue{\handwritten{[无法辨认]}} \\
\hline
设备编号 &
% #VALUE_ID: LEX-P0008-V0005
% #FIELD_VALUE: 设备编号
% #TODO #HANDWRITTEN: [无法辨认]; 字迹及图像分辨率不足
\fieldvalue{\handwritten{[无法辨认]}} &
检验日期 &
% #VALUE_ID: LEX-P0008-V0006
% #FIELD_VALUE: 检验日期
% #TODO #HANDWRITTEN: [无法辨认]; 字迹及图像分辨率不足
\fieldvalue{\handwritten{[无法辨认]}} \\
\hline
\end{tabularx}

\vspace{0.5em}

\noindent
\begin{tabularx}{\textwidth}{|c|X|c|}
\hline
序号 & 检验项目 & 检验结果 \\
\hline
7 &
垂直度偏差应符合规定要求，基础及附着装置应满足使用要求。 &
% #VALUE_ID: LEX-P0008-V0007
% #FIELD_VALUE: 检验结果（第7项）
% #HANDWRITTEN: 合格
\fieldvalue{\handwritten{合格}} \\
\hline
8 &
各安全装置应齐全、灵敏、可靠；限位装置应符合有关规定。 &
% #VALUE_ID: LEX-P0008-V0008
% #FIELD_VALUE: 检验结果（第8项）
% #HANDWRITTEN: √
\fieldvalue{\handwritten{√}} \\
\hline
9 &
钢丝绳、滑轮及卷筒应符合安全技术要求，连接及固定应牢固。 &
% #VALUE_ID: LEX-P0008-V0009
% #FIELD_VALUE: 检验结果（第9项）
% #HANDWRITTEN: √
\fieldvalue{\handwritten{√}} \\
\hline
10 &
制动器应工作可靠，制动性能应符合要求。 &
% #VALUE_ID: LEX-P0008-V0010
% #FIELD_VALUE: 检验结果（第10项）
% #HANDWRITTEN: √
\fieldvalue{\handwritten{√}} \\
\hline
11 &
起重机与输电线路的安全距离应符合规定，接地及避雷装置应可靠。 &
% #VALUE_ID: LEX-P0008-V0011
% #FIELD_VALUE: 检验结果（第11项）
% #HANDWRITTEN: √
\fieldvalue{\handwritten{√}} \\
\hline
12 &
电气系统、控制系统及照明系统应安装正确、运行正常。 &
% #VALUE_ID: LEX-P0008-V0012
% #FIELD_VALUE: 检验结果（第12项）
% #HANDWRITTEN: √
\fieldvalue{\handwritten{√}} \\
\hline
13 &
空载、额定载荷及相应工况下的运行试验应符合要求。 &
% #VALUE_ID: LEX-P0008-V0013
% #FIELD_VALUE: 检验结果（第13项）
% #TODO #HANDWRITTEN: [疑似“合格”]; 手写内容模糊
\fieldvalue{\handwritten{[疑似“合格”]}} \\
\hline
\end{tabularx}

% TODO：本页表格内容在页面边界处可能与相邻页面连续；仅保留本页可见部分。
% LEXOID_PAGE_COMPLETED: 8/78

\newpage

\section*{8.2.2 供应商资格审查表}

\begin{tabularx}{\textwidth}{|c|X|c|}
\hline
版本 & 变更记录 & 日期 \\
\hline
% #VALUE_ID: LEX-P0009-V0001
% #FIELD_VALUE: 日期
% #HANDWRITTEN: 2023年
\fieldvalue{\handwritten{2023年}} &
% #VALUE_ID: LEX-P0009-V0002
% #FIELD_VALUE: 变更内容
% #TODO #HANDWRITTEN: 内容不清；蓝色手写文字难以辨认
\fieldvalue{\handwritten{内容不清}} &
% #VALUE_ID: LEX-P0009-V0003
% #FIELD_VALUE: 版本
% #TODO #HANDWRITTEN: 版本号不清；蓝色手写字符难以辨认
\fieldvalue{\handwritten{不清}} \\
\hline
\end{tabularx}

\vspace{0.5em}

\begin{tabularx}{\textwidth}{|c|X|c|}
\hline
序号 & 检查项目 & 审核结果 \\
\hline
1 &
营业执照及相关证明文件是否齐全、有效 &
% #VALUE_ID: LEX-P0009-V0004
% #FIELD_VALUE: 第1项审核结果
% #HANDWRITTEN: ✓
\fieldvalue{\handwritten{\checkmark}} \\
\hline
2 &
资质证书及其相关证明材料是否齐全、有效 &
% #VALUE_ID: LEX-P0009-V0005
% #FIELD_VALUE: 第2项审核结果
% #HANDWRITTEN: ✓
\fieldvalue{\handwritten{\checkmark}} \\
\hline
3 &
安全生产许可证等相关证照是否在有效期内 &
% #VALUE_ID: LEX-P0009-V0006
% #FIELD_VALUE: 第3项审核结果
% #HANDWRITTEN: ✓
\fieldvalue{\handwritten{\checkmark}} \\
\hline
4 &
企业信用及相关资信证明材料是否符合要求 &
% #VALUE_ID: LEX-P0009-V0007
% #FIELD_VALUE: 第4项审核结果
% #HANDWRITTEN: ✓
\fieldvalue{\handwritten{\checkmark}} \\
\hline
5 &
近三年是否存在重大违法、违规记录 &
% #VALUE_ID: LEX-P0009-V0008
% #FIELD_VALUE: 第5项审核结果
% #HANDWRITTEN: ✓
\fieldvalue{\handwritten{\checkmark}} \\
\hline
6 &
是否具有履行合同所必需的设备和专业技术能力 &
% #VALUE_ID: LEX-P0009-V0009
% #FIELD_VALUE: 第6项审核结果
% #HANDWRITTEN: ✓
\fieldvalue{\handwritten{\checkmark}} \\
\hline
7 &
是否具有良好的商业信誉和健全的财务会计制度 &
% #VALUE_ID: LEX-P0009-V0010
% #FIELD_VALUE: 第7项审核结果
% #HANDWRITTEN: ✓
\fieldvalue{\handwritten{\checkmark}} \\
\hline
8 &
是否依法缴纳税收和社会保障资金 &
% #VALUE_ID: LEX-P0009-V0011
% #FIELD_VALUE: 第8项审核结果
% #HANDWRITTEN: ✓
\fieldvalue{\handwritten{\checkmark}} \\
\hline
9 &
法定代表人或授权代表相关资料是否齐全 &
% #VALUE_ID: LEX-P0009-V0012
% #FIELD_VALUE: 第9项审核结果
% #HANDWRITTEN: ✓
\fieldvalue{\handwritten{\checkmark}} \\
\hline
10 &
供应商响应文件是否按要求编制并加盖公章 &
% #VALUE_ID: LEX-P0009-V0013
% #FIELD_VALUE: 第10项审核结果
% #HANDWRITTEN: ✓
\fieldvalue{\handwritten{\checkmark}} \\
\hline
11 &
供应商是否满足采购文件规定的其他资格条件 &
% #VALUE_ID: LEX-P0009-V0014
% #FIELD_VALUE: 第11项审核结果
% #HANDWRITTEN: ✓
\fieldvalue{\handwritten{\checkmark}} \\
\hline
12 &
其他审查事项 &
% #VALUE_ID: LEX-P0009-V0015
% #FIELD_VALUE: 第12项审核结果
% #HANDWRITTEN: ✓
\fieldvalue{\handwritten{\checkmark}} \\
\hline
\end{tabularx}

% TODO: 本页表格中的部分中文检查项目及蓝色手写内容因扫描方向和分辨率较低，需结合原始 PDF 复核。
% LEXOID_PAGE_COMPLETED: 9/78

\newpage

\section*{检验记录表}

% TODO: The page is rotated and several Chinese characters are low-resolution; transcription below preserves the visible table structure and legible text.

\begin{tabularx}{\textwidth}{|p{0.10\textwidth}|X|p{0.18\textwidth}|}
\hline
\textbf{项目} & \textbf{检验内容} & \textbf{结果} \\
\hline
产品名称及规格型号 &
% #VALUE_ID: LEX-P0010-V0001
% #FIELD_VALUE: 产品名称及规格型号
% #TODO #HANDWRITTEN: 202?; handwritten blue entry is partially illegible
\fieldvalue{\handwritten{202?}} &
\\
\hline
(1) &
% #VALUE_ID: LEX-P0010-V0002
% #FIELD_VALUE: (1) 检验结果
% #HANDWRITTEN: ✓
\fieldvalue{\handwritten{✓}} &
是\quad 否 \\
\hline
(2) &
% #VALUE_ID: LEX-P0010-V0003
% #FIELD_VALUE: (2) 检验结果
% #HANDWRITTEN: ✓
\fieldvalue{\handwritten{✓}} &
是\quad 否 \\
\hline
(3) &
% #VALUE_ID: LEX-P0010-V0004
% #FIELD_VALUE: (3) 检验结果
% #HANDWRITTEN: ✓
\fieldvalue{\handwritten{✓}} &
是\quad 否 \\
\hline
(4) &
% #VALUE_ID: LEX-P0010-V0005
% #FIELD_VALUE: (4) 检验结果
% #HANDWRITTEN: ✓
\fieldvalue{\handwritten{✓}} &
是\quad 否 \\
\hline
(5) &
% #VALUE_ID: LEX-P0010-V0006
% #FIELD_VALUE: (5) 检验结果
% #HANDWRITTEN: ✓
\fieldvalue{\handwritten{✓}} &
是\quad 否 \\
\hline
(6) &
% #VALUE_ID: LEX-P0010-V0007
% #FIELD_VALUE: (6) 检验结果
% #HANDWRITTEN: ✓
\fieldvalue{\handwritten{✓}} &
是\quad 否 \\
\hline
(7) &
% #VALUE_ID: LEX-P0010-V0008
% #FIELD_VALUE: (7) 检验结果
% #HANDWRITTEN: ✓
\fieldvalue{\handwritten{✓}} &
是\quad 否 \\
\hline
(8) &
% #VALUE_ID: LEX-P0010-V0009
% #FIELD_VALUE: (8) 检验结果
% #HANDWRITTEN: ✓
\fieldvalue{\handwritten{✓}} &
是\quad 否 \\
\hline
(9) &
% #VALUE_ID: LEX-P0010-V0010
% #FIELD_VALUE: (9) 检验结果
% #HANDWRITTEN: ✓
\fieldvalue{\handwritten{✓}} &
是\quad 否 \\
\hline
(10) &
% #VALUE_ID: LEX-P0010-V0011
% #FIELD_VALUE: (10) 检验结果
% #HANDWRITTEN: ✓
\fieldvalue{\handwritten{✓}} &
是\quad 否 \\
\hline
(11) &
% #VALUE_ID: LEX-P0010-V0012
% #FIELD_VALUE: (11) 检验结果
% #HANDWRITTEN: ✓
\fieldvalue{\handwritten{✓}} &
是\quad 否 \\
\hline
(12) &
% #VALUE_ID: LEX-P0010-V0013
% #FIELD_VALUE: (12) 检验结果
% #HANDWRITTEN: ✓
\fieldvalue{\handwritten{✓}} &
是\quad 否 \\
\hline
(13) &
检验项目及判定依据：见表内相关条款。 &
\\
\hline
(14) &
检查、检验及判定结果等内容。 &
\\
\hline
(15) &
审核工作日期：%
% #VALUE_ID: LEX-P0010-V0014
% #FIELD_VALUE: 审核工作日期
% #HANDWRITTEN: 11-10-20
\fieldvalue{\handwritten{11-10-20}} &
\\
\hline
(16) &
检验项目及判定依据、检验结果等。 &
\\
\hline
(17) &
检验结论及有关说明。 &
\\
\hline
\end{tabularx}

\vspace{0.5em}

\noindent
备注：检验项目、检验方法及判定依据等内容见表内说明。

\noindent
% TODO: The lower form header contains additional printed identifiers and a partially illegible handwritten entry; the visible printed code appears to be ``200-408''.
% LEXOID_PAGE_COMPLETED: 10/78

\newpage

\begin{center}
\textbf{水泥混凝土拌合物试验记录}
\end{center}

\noindent
\begin{tabularx}{\textwidth}{|l|X|l|X|}
\hline
工程名称 &  & 记录编号 &  \\ \hline
试验日期 &
% #VALUE_ID: LEX-P0011-V0001
% #FIELD_VALUE: 试验日期
\fieldvalue{2022-10-01}
& 试验人员 &  \\ \hline
\end{tabularx}

\vspace{2mm}

\noindent
\begin{tabularx}{\textwidth}{|l|X|}
\hline
试验项目及依据标准 &
（1）水泥混凝土拌合物性能试验；\\
& （2）试验方法按相关标准执行。 \\ \hline
试验结果 &
% #VALUE_ID: LEX-P0011-V0002
% #FIELD_VALUE: 试验结果判定
% #HANDWRITTEN: √ 合格
\fieldvalue{\handwritten{√ 合格}}
\\ \hline
备注 &
% #VALUE_ID: LEX-P0011-V0003
% #FIELD_VALUE: 备注
% #HANDWRITTEN: √ 合格
\fieldvalue{\handwritten{√ 合格}}
\\ \hline
\end{tabularx}

\vspace{2mm}

\noindent
\small
\begin{tabularx}{\textwidth}{|X|c|}
\hline
\multicolumn{2}{|c|}{试验数据记录} \\ \hline
项目 & 数值 \\ \hline
混凝土坍落度（mm） &
% #VALUE_ID: LEX-P0011-V0004
% #FIELD_VALUE: 混凝土坍落度（mm）
\fieldvalue{60}
\\ \hline
混凝土扩展度（mm） &
% #VALUE_ID: LEX-P0011-V0005
% #FIELD_VALUE: 混凝土扩展度（mm）
\fieldvalue{500}
\\ \hline
粗集料最大粒径（mm） &
% #VALUE_ID: LEX-P0011-V0006
% #FIELD_VALUE: 粗集料最大粒径（mm）
\fieldvalue{800}
\\ \hline
水泥用量（kg/m$^3$） &
% #VALUE_ID: LEX-P0011-V0007
% #FIELD_VALUE: 水泥用量（kg/m$^3$）
\fieldvalue{240}
\\ \hline
砂用量（kg/m$^3$） &
% #VALUE_ID: LEX-P0011-V0008
% #FIELD_VALUE: 砂用量（kg/m$^3$）
\fieldvalue{1800}
\\ \hline
石用量（kg/m$^3$） &
% #VALUE_ID: LEX-P0011-V0009
% #FIELD_VALUE: 石用量（kg/m$^3$）
\fieldvalue{150}
\\ \hline
水用量（kg/m$^3$） &
% #VALUE_ID: LEX-P0011-V0010
% #FIELD_VALUE: 水用量（kg/m$^3$）
\fieldvalue{400}
\\ \hline
外加剂用量（\%） &
% #VALUE_ID: LEX-P0011-V0011
% #FIELD_VALUE: 外加剂用量（\%）
\fieldvalue{2}
\\ \hline
含气量（\%） &
% #VALUE_ID: LEX-P0011-V0012
% #FIELD_VALUE: 含气量（\%）
\fieldvalue{180}
\\ \hline
试验温度（$^\circ$C） &
% #VALUE_ID: LEX-P0011-V0013
% #FIELD_VALUE: 试验温度（$^\circ$C）
\fieldvalue{18}
\\ \hline
试件尺寸（mm） &
% #VALUE_ID: LEX-P0011-V0014
% #FIELD_VALUE: 试件尺寸（mm）
\fieldvalue{150}
\\ \hline
试件数量（个） &
% #VALUE_ID: LEX-P0011-V0015
% #FIELD_VALUE: 试件数量（个）
\fieldvalue{3}
\\ \hline
养护龄期（d） &
% #VALUE_ID: LEX-P0011-V0016
% #FIELD_VALUE: 养护龄期（d）
\fieldvalue{4}
\\ \hline
\end{tabularx}
\normalsize

% TODO: The lower portion of the form contains additional tabular entries and printed values; some labels and digits are illegible in the provided page image.
% LEXOID_PAGE_COMPLETED: 11/78

\newpage

% TODO: The source page is rotated and several printed characters are faint or illegible.
\section*{會議提案／追蹤事項表}

\begin{flushright}
第 12 頁
\end{flushright}

\noindent
\textbf{國立臺灣大學\quad 會議紀錄及追蹤事項}\\[2pt]
\textbf{（22）\quad [文件標題部分難辨]}
\hfill
\textit{日期：[部分難辨]}

\small
\renewcommand{\arraystretch}{1.35}
\begin{tabularx}{\textwidth}{|>{\RaggedRight\arraybackslash}p{0.16\textwidth}|X|}
\hline
\textbf{案號／序號} & \textbf{案由及辦理情形} \\
\hline
\multicolumn{2}{|l|}{\textbf{承辦單位：}[固定欄位內容部分難辨]\quad
\textbf{日期：}[固定欄位內容部分難辨]}\\
\hline

11 &
學門、學術研究及相關業務事項；辦理情形及說明如下。\\
\cline{1-2}
&
% #VALUE_ID: LEX-P0012-V0001
% #FIELD_VALUE: 第11案辦理情形欄
% TODO #HANDWRITTEN: 5025?; 藍筆字跡不清
\fieldvalue{\handwritten{5025?}}\\
\hline

10 &
有關學術研究及教學相關事項。\\
\cline{1-2}
&
% #VALUE_ID: LEX-P0012-V0002
% #FIELD_VALUE: 第10案辦理情形欄
% TODO #HANDWRITTEN: 2?; 藍筆字跡不清
\fieldvalue{\handwritten{2?}}\\
\hline

9 &
有關教師聘任、聘期及相關規定之事項。\\
\cline{1-2}
&
% #VALUE_ID: LEX-P0012-V0003
% #FIELD_VALUE: 第9案辦理情形欄
% TODO #HANDWRITTEN: 2?; 藍筆字跡不清
\fieldvalue{\handwritten{2?}}\\
\hline

8 &
有關學生事務及教學行政之事項。\\
\cline{1-2}
&
% #VALUE_ID: LEX-P0012-V0004
% #FIELD_VALUE: 第8案辦理情形欄
% TODO #HANDWRITTEN: 2?; 藍筆字跡不清
\fieldvalue{\handwritten{2?}}\\
\hline

7 &
有關研究計畫、研究成果及經費使用事項。\\
\cline{1-2}
&
% #VALUE_ID: LEX-P0012-V0005
% #FIELD_VALUE: 第7案辦理情形欄
% TODO #HANDWRITTEN: 1?; 藍筆字跡不清
\fieldvalue{\handwritten{1?}}\\
\hline

6 &
有關校務發展及行政業務之事項。\\
\cline{1-2}
&
% #VALUE_ID: LEX-P0012-V0006
% #FIELD_VALUE: 第6案辦理情形欄
% TODO #HANDWRITTEN: 2?; 藍筆字跡不清
\fieldvalue{\handwritten{2?}}\\
\hline

5 &
有關會議決議及後續辦理情形之事項。\\
\cline{1-2}
&
% #VALUE_ID: LEX-P0012-V0007
% #FIELD_VALUE: 第5案辦理情形欄
% TODO #HANDWRITTEN: 18; 藍筆數字可辨識但欄位內容不清
\fieldvalue{\handwritten{18}}\\
\hline

4 &
有關校內規章及業務執行之事項。\\
\cline{1-2}
&
% #VALUE_ID: LEX-P0012-V0008
% #FIELD_VALUE: 第4案辦理情形欄
% TODO #HANDWRITTEN: 1?; 藍筆字跡不清
\fieldvalue{\handwritten{1?}}\\
\hline

3 &
有關人事及組織調整之事項。\\
\cline{1-2}
&
% #VALUE_ID: LEX-P0012-V0009
% #FIELD_VALUE: 第3案辦理情形欄
% TODO #HANDWRITTEN: 2?; 藍筆字跡不清
\fieldvalue{\handwritten{2?}}\\
\hline

2 &
有關經費、採購及財產管理之事項。\\
\cline{1-2}
&
% #VALUE_ID: LEX-P0012-V0010
% #FIELD_VALUE: 第2案辦理情形欄
% TODO #HANDWRITTEN: 18; 藍筆數字可辨識但欄位內容不清
\fieldvalue{\handwritten{18}}\\
\hline

1 &
有關前次會議提案及追蹤事項。\\
\cline{1-2}
&
% #VALUE_ID: LEX-P0012-V0011
% #FIELD_VALUE: 第1案辦理情形欄
% TODO #HANDWRITTEN: 19; 藍筆數字可辨識但欄位內容不清
\fieldvalue{\handwritten{19}}\\
\hline
\end{tabularx}

\vspace{4pt}

\noindent
\textbf{備註：} [頁面下方固定說明文字及印章內容部分難辨]

% #VALUE_ID: LEX-P0012-V0012
% #FIELD_VALUE: 頁首日期／編號欄
% TODO #HANDWRITTEN: 15/09/??; 藍筆日期字跡不清
\fieldvalue{\handwritten{15/09/??}}

% TODO: This page appears to continue a larger table; only the rows visible on this page are represented.
% LEXOID_PAGE_COMPLETED: 12/78

\newpage

% TODO: The page is largely blank and appears to contain an upside-down scanned Chinese construction form; the following reproduces the visible form content in corrected reading orientation.

\begin{center}
\textbf{安徽省建设工程监理资料表（22）}\\
% TODO: The complete title is partially illegible in the scan.
\end{center}

\begin{flushright}
% TODO: A red official seal is visible near the upper-left area of the form.
\fbox{\scriptsize 红色印章}
\end{flushright}

\small
\begin{tabularx}{\textwidth}{|p{2.6cm}|X|p{2.6cm}|X|}
\hline
工程名称 &
% #VALUE_ID: LEX-P0013-V0001
% #FIELD_VALUE: 工程名称
% #TODO #HANDWRITTEN: 难以辨认; 蓝色手写内容模糊
\fieldvalue{\handwritten{难以辨认}} &
建设单位 &
% #VALUE_ID: LEX-P0013-V0002
% #FIELD_VALUE: 建设单位
% #TODO #HANDWRITTEN: 难以辨认; 蓝色手写内容模糊
\fieldvalue{\handwritten{难以辨认}} \\
\hline
监理单位 &
% #VALUE_ID: LEX-P0013-V0003
% #FIELD_VALUE: 监理单位
% #TODO #HANDWRITTEN: 难以辨认; 蓝色手写内容模糊
\fieldvalue{\handwritten{难以辨认}} &
项目负责人 &
% #VALUE_ID: LEX-P0013-V0004
% #FIELD_VALUE: 项目负责人
% #TODO #HANDWRITTEN: 难以辨认; 蓝色手写内容模糊
\fieldvalue{\handwritten{难以辨认}} \\
\hline
工程地点 &
% #VALUE_ID: LEX-P0013-V0005
% #FIELD_VALUE: 工程地点
% #TODO #HANDWRITTEN: 难以辨认; 蓝色手写内容模糊
\fieldvalue{\handwritten{难以辨认}} &
合同编号 &
% #VALUE_ID: LEX-P0013-V0006
% #FIELD_VALUE: 合同编号
% #TODO #HANDWRITTEN: 难以辨认; 蓝色手写内容模糊
\fieldvalue{\handwritten{难以辨认}} \\
\hline
\multicolumn{4}{|l|}{\textbf{监理工作内容及相关说明}} \\
\hline
\multicolumn{4}{|p{\textwidth}|}{
监理工作内容：工程质量、进度及安全生产等相关监理工作。
% TODO: The detailed printed text in this cell is too blurred to transcribe reliably.
} \\
\hline
\multicolumn{2}{|p{0.47\textwidth}|}{
监理单位意见：\newline\rule{0pt}{1.2cm}
} &
\multicolumn{2}{p{0.47\textwidth}|}{
建设单位意见：\newline\rule{0pt}{1.2cm}
} \\
\hline
\end{tabularx}

\normalsize

% TODO: Additional printed footer text and page numbering are visible on the scan but are too indistinct to transcribe confidently.
% LEXOID_PAGE_COMPLETED: 13/78

\newpage

% TODO: The source page is rotated 180 degrees; the following reconstruction reflects the readable orientation.
\section*{医薬品関連様式}

% #VALUE_ID: LEX-P0014-V0001
% #FIELD_VALUE: 申請・届出欄
% #TODO #HANDWRITTEN: 判読不能; 赤色手書き文字が不鮮明
\fieldvalue{\handwritten{判読不能}}

\begin{tabularx}{\textwidth}{|p{0.24\textwidth}|X|}
\hline
文書番号 & 
% #VALUE_ID: LEX-P0014-V0002
% #FIELD_VALUE: 文書番号
% #TODO #HANDWRITTEN: 200-408; 数字列の一部が不鮮明
\fieldvalue{\handwritten{200-408}} \\
\hline
\end{tabularx}

\vspace{1em}

\noindent
\textbf{製造販売承認申請書類}

\begin{tabularx}{\textwidth}{|p{0.30\textwidth}|X|}
\hline
項目 & 記載内容 \\
\hline
販売名 & \\
\hline
一般名 & \\
\hline
剤形・含量 & \\
\hline
製造販売業者 & \\
\hline
承認番号 & \\
\hline
承認年月日 &
% #VALUE_ID: LEX-P0014-V0003
% #FIELD_VALUE: 承認年月日
% #TODO #HANDWRITTEN: 2025.06.??; 日付末尾が不鮮明
\fieldvalue{\handwritten{2025.06.??}} \\
\hline
\end{tabularx}

\vspace{1em}

\noindent
\textbf{添付資料}

\begin{tabularx}{\textwidth}{|c|X|c|}
\hline
\multicolumn{3}{|c|}{資料の種類} \\
\cline{1-3}
番号 & 資料名 & 確認 \\
\hline
1 & 医薬品の概要 & 
% #VALUE_ID: LEX-P0014-V0004
% #FIELD_VALUE: 資料1の確認
% #HANDWRITTEN: ✓
\fieldvalue{\handwritten{\checkmark}} \\
\hline
2 & 起源または発見の経緯 & 
% #VALUE_ID: LEX-P0014-V0005
% #FIELD_VALUE: 資料2の確認
% #HANDWRITTEN: ✓
\fieldvalue{\handwritten{\checkmark}} \\
\hline
3 & 物理的化学的性質ならびに規格及び試験方法 & 
% #VALUE_ID: LEX-P0014-V0006
% #FIELD_VALUE: 資料3の確認
% #HANDWRITTEN: ✓
\fieldvalue{\handwritten{\checkmark}} \\
\hline
4 & 安定性 & 
% #VALUE_ID: LEX-P0014-V0007
% #FIELD_VALUE: 資料4の確認
% #HANDWRITTEN: ✓
\fieldvalue{\handwritten{\checkmark}} \\
\hline
5 & 効能又は効果 & 
% #VALUE_ID: LEX-P0014-V0008
% #FIELD_VALUE: 資料5の確認
% #HANDWRITTEN: ✓
\fieldvalue{\handwritten{\checkmark}} \\
\hline
6 & 用法及び用量 & 
% #VALUE_ID: LEX-P0014-V0009
% #FIELD_VALUE: 資料6の確認
% #HANDWRITTEN: ✓
\fieldvalue{\handwritten{\checkmark}} \\
\hline
\end{tabularx}

\vspace{1em}

\noindent
\textbf{備考}

\begin{tabularx}{\textwidth}{|p{0.30\textwidth}|X|}
\hline
薬価 & \\
\hline
薬価基準収載年月日 &
% #VALUE_ID: LEX-P0014-V0010
% #FIELD_VALUE: 薬価基準収載年月日
% #TODO #HANDWRITTEN: 25/??; 手書き日付が不鮮明
\fieldvalue{\handwritten{25/??}} \\
\hline
その他 &
% #VALUE_ID: LEX-P0014-V0011
% #FIELD_VALUE: その他
% #TODO #HANDWRITTEN: 判読不能; 青色手書き文字が不鮮明
\fieldvalue{\handwritten{判読不能}} \\
\hline
\end{tabularx}

\vspace{1em}

\noindent
\textbf{記載年月日：}
% #VALUE_ID: LEX-P0014-V0012
% #FIELD_VALUE: 記載年月日
% #TODO #HANDWRITTEN: 8/??; 手書きの日付が不鮮明
\fieldvalue{\handwritten{8/??}}

\noindent
\textbf{様式番号：} Z5-MAP-01-0-01

\noindent
\textbf{関連様式番号：} Z5-MAP-04-0-01

% TODO: The lower portion of the page contains additional printed form rows and continuation content that is partially illegible in the supplied scan.
% LEXOID_PAGE_COMPLETED: 14/78

\newpage

\section{工艺记录}

% #VALUE_ID: LEX-P0015-V0001
% #HANDWRITTEN: 1/8
\handwritten{1/8}

\begin{tabularx}{\textwidth}{@{}lX lX@{}}
工艺编号 &
% #VALUE_ID: LEX-P0015-V0002
% #FIELD_VALUE: 工艺编号
\fieldvalue{PP-Ham-002}
&
工艺名称 &
% #VALUE_ID: LEX-P0015-V0003
% #FIELD_VALUE: 工艺名称
\fieldvalue{\textit{[无法辨认]}} \\

生产批号 &
% #VALUE_ID: LEX-P0015-V0004
% #FIELD_VALUE: 生产批号
\fieldvalue{2025-06-18 19:29:45}
&
记录时间 &
% #VALUE_ID: LEX-P0015-V0005
% #FIELD_VALUE: 记录时间
\fieldvalue{\textit{[无法辨认]}} \\

工艺版本 &
% #VALUE_ID: LEX-P0015-V0006
% #FIELD_VALUE: 工艺版本
\fieldvalue{V1.18}
&
产品编码 &
% #VALUE_ID: LEX-P0015-V0007
% #FIELD_VALUE: 产品编码
\fieldvalue{20240336} \\

产品型号 &
% #VALUE_ID: LEX-P0015-V0008
% #FIELD_VALUE: 产品型号
\fieldvalue{G32-291-L-005}
&
生产日期 &
% #VALUE_ID: LEX-P0015-V0009
% #FIELD_VALUE: 生产日期
\fieldvalue{2026.04.10} \\

产品批号 &
% #VALUE_ID: LEX-P0015-V0010
% #FIELD_VALUE: 产品批号
\fieldvalue{20260410}
&
产品规格 &
% #VALUE_ID: LEX-P0015-V0011
% #FIELD_VALUE: 产品规格
\fieldvalue{\textit{[无法辨认]}} \\

设备编号 &
% #VALUE_ID: LEX-P0015-V0012
% #FIELD_VALUE: 设备编号
\fieldvalue{\textit{[无法辨认]}}
&
设备名称 &
% #VALUE_ID: LEX-P0015-V0013
% #FIELD_VALUE: 设备名称
\fieldvalue{\textit{[无法辨认]}} \\

批次号 &
% #VALUE_ID: LEX-P0015-V0014
% #FIELD_VALUE: 批次号
\fieldvalue{\textit{[无法辨认]}}
&
操作员 &
% #VALUE_ID: LEX-P0015-V0015
% #FIELD_VALUE: 操作员
\fieldvalue{\textit{[无法辨认]}}
\end{tabularx}

\vspace{0.5em}

% #TODO #HANDWRITTEN: [无法辨认]; 蓝色手写批注
% #VALUE_ID: LEX-P0015-V0016
% #HANDWRITTEN: [无法辨认]
\noindent
\textbf{手写记录：}\ \handwritten{[无法辨认]}

% #TODO #HANDWRITTEN: [无法辨认]; 蓝色手写日期或批次信息
% #VALUE_ID: LEX-P0015-V0017
% #HANDWRITTEN: [无法辨认]
\noindent
\textbf{手写日期/批次：}\ \handwritten{[无法辨认]}

\subsection{物料记录}

\begin{tabularx}{\textwidth}{@{}lX@{}}
物料一 &
% #VALUE_ID: LEX-P0015-V0018
% #FIELD_VALUE: 物料一用量
\fieldvalue{7426.7\,mL} \\

物料二 &
% #VALUE_ID: LEX-P0015-V0019
% #FIELD_VALUE: 物料二用量
\fieldvalue{1437.8\,mL} \\

物料三 &
% #VALUE_ID: LEX-P0015-V0020
% #FIELD_VALUE: 物料三用量
\fieldvalue{188.5\,mL} \\

物料四 &
% #VALUE_ID: LEX-P0015-V0021
% #FIELD_VALUE: 物料四用量
\fieldvalue{1128.5\,mL} \\

物料五 &
% #VALUE_ID: LEX-P0015-V0022
% #FIELD_VALUE: 物料五用量
\fieldvalue{512.9\,mL} \\

物料六 &
% #VALUE_ID: LEX-P0015-V0023
% #FIELD_VALUE: 物料六用量
\fieldvalue{(1346.3\,mL/216.6\,mL)}
\end{tabularx}

\subsection{工艺操作记录}

\small
\begin{tabularx}{\textwidth}{@{}lX@{}}
\textbf{操作时间} & \textbf{工艺操作描述} \\ \hline
% #VALUE_ID: LEX-P0015-V0024
% #FIELD_VALUE: 操作时间
\fieldvalue{16:15:00} &
% #VALUE_ID: LEX-P0015-V0025
% #FIELD_VALUE: 工艺操作描述
\fieldvalue{\textit{[无法辨认]}} \\

% #VALUE_ID: LEX-P0015-V0026
% #FIELD_VALUE: 操作时间
\fieldvalue{16:15:00} &
% #VALUE_ID: LEX-P0015-V0027
% #FIELD_VALUE: 工艺操作描述
\fieldvalue{\textit{[无法辨认]}} \\

% #VALUE_ID: LEX-P0015-V0028
% #FIELD_VALUE: 操作时间
\fieldvalue{16:15:00} &
% #VALUE_ID: LEX-P0015-V0029
% #FIELD_VALUE: 工艺操作描述
\fieldvalue{\textit{[无法辨认]}} \\

% #VALUE_ID: LEX-P0015-V0030
% #FIELD_VALUE: 操作时间
\fieldvalue{16:15:00} &
% #VALUE_ID: LEX-P0015-V0031
% #FIELD_VALUE: 工艺操作描述
\fieldvalue{\textit{[无法辨认]}} \\

% #VALUE_ID: LEX-P0015-V0032
% #FIELD_VALUE: 操作时间
\fieldvalue{16:25:24} &
% #VALUE_ID: LEX-P0015-V0033
% #FIELD_VALUE: 工艺操作描述
\fieldvalue{\textit{[无法辨认]}} \\

% #VALUE_ID: LEX-P0015-V0034
% #FIELD_VALUE: 操作时间
\fieldvalue{16:27:46} &
% #VALUE_ID: LEX-P0015-V0035
% #FIELD_VALUE: 工艺操作描述
\fieldvalue{\textit{[无法辨认]}} \\

% #VALUE_ID: LEX-P0015-V0036
% #FIELD_VALUE: 操作时间
\fieldvalue{16:29:50} &
% #VALUE_ID: LEX-P0015-V0037
% #FIELD_VALUE: 工艺操作描述
\fieldvalue{\textit{[无法辨认]}} \\

% #VALUE_ID: LEX-P0015-V0038
% #FIELD_VALUE: 操作时间
\fieldvalue{16:30:07} &
% #VALUE_ID: LEX-P0015-V0039
% #FIELD_VALUE: 工艺操作描述
\fieldvalue{\textit{[无法辨认]}} \\

% #VALUE_ID: LEX-P0015-V0040
% #FIELD_VALUE: 操作时间
\fieldvalue{16:30:26} &
% #VALUE_ID: LEX-P0015-V0041
% #FIELD_VALUE: 工艺操作描述
\fieldvalue{\textit{[无法辨认]}} \\

% #VALUE_ID: LEX-P0015-V0042
% #FIELD_VALUE: 操作时间
\fieldvalue{16:37:50} &
% #VALUE_ID: LEX-P0015-V0043
% #FIELD_VALUE: 工艺操作描述
\fieldvalue{\textit{[无法辨认]}} \\

% #VALUE_ID: LEX-P0015-V0044
% #FIELD_VALUE: 操作时间
\fieldvalue{16:37:55} &
% #VALUE_ID: LEX-P0015-V0045
% #FIELD_VALUE: 工艺操作描述
\fieldvalue{\textit{[无法辨认]}} \\

% #VALUE_ID: LEX-P0015-V0046
% #FIELD_VALUE: 操作时间
\fieldvalue{16:38:01} &
% #VALUE_ID: LEX-P0015-V0047
% #FIELD_VALUE: 工艺操作描述
\fieldvalue{\textit{[无法辨认]}} \\

% #VALUE_ID: LEX-P0015-V0048
% #FIELD_VALUE: 操作时间
\fieldvalue{16:38:09} &
% #VALUE_ID: LEX-P0015-V0049
% #FIELD_VALUE: 工艺操作描述
\fieldvalue{\textit{[无法辨认]}} \\

% #VALUE_ID: LEX-P0015-V0050
% #FIELD_VALUE: 操作时间
\fieldvalue{16:38:15} &
% #VALUE_ID: LEX-P0015-V0051
% #FIELD_VALUE: 工艺操作描述
\fieldvalue{\textit{[无法辨认]}} \\

% #VALUE_ID: LEX-P0015-V0052
% #FIELD_VALUE: 操作时间
\fieldvalue{16:38:51} &
% #VALUE_ID: LEX-P0015-V0053
% #FIELD_VALUE: 工艺操作描述
\fieldvalue{\textit{[无法辨认]}} \\

% #VALUE_ID: LEX-P0015-V0054
% #FIELD_VALUE: 操作时间
\fieldvalue{16:38:52} &
% #VALUE_ID: LEX-P0015-V0055
% #FIELD_VALUE: 工艺操作描述
\fieldvalue{\textit{[无法辨认]}} \\

% #VALUE_ID: LEX-P0015-V0056
% #FIELD_VALUE: 操作时间
\fieldvalue{16:38:55} &
% #VALUE_ID: LEX-P0015-V0057
% #FIELD_VALUE: 工艺操作描述
\fieldvalue{\textit{[无法辨认]}}
\end{tabularx}
\normalsize

% TODO: The operation-record table continues beyond the visible portion of this page.
% LEXOID_PAGE_COMPLETED: 15/78

\newpage

% #VALUE_ID: LEX-P0016-V0001
% #HANDWRITTEN: 2/18
\marginpar{\handwritten{2/18}}

\rotatebox{90}{%
\begin{minipage}{0.96\textheight}
\scriptsize
\begin{tabularx}{\textwidth}{@{}p{1.2cm}X@{\hspace{1.5em}}p{1.2cm}X@{}}
16:39:05 & 進出料泵轉速：5.1 mL/min，1號泵 &
16:56:05 & 進出料泵轉速：5.1 mL/min，1號泵 \\ 
16:39:05 & 進出料泵轉速：6.3 mL/min，2號泵 &
16:56:05 & 進出料泵轉速：6.3 mL/min，2號泵 \\ 
16:39:08 & 進出料泵開始，90 mL/min，1號泵 &
16:56:05 & 進出料泵停止，90 mL/min，1號泵 \\ 
16:39:08 & 進出料泵開始，90 mL/min，2號泵 &
17:04:16 & 進出料泵停止，90 mL/min，2號泵 \\ 
16:39:25 & 進出料泵開始，90 mL/min，1號泵 &
17:04:16 & 進出料泵開始，300 mL/min，1號泵 \\ 
16:39:32 & 進出料泵停止，90 mL/min，1號泵 &
17:05:47 & 進出料泵停止，450 mL/min，1號泵 \\ 
16:39:32 & 進出料泵開始，90 mL/min，2號泵 &
17:05:47 & 進出料泵開始，450 mL/min，1號泵 \\ 
16:39:54 & 進出料泵停止，90 mL/min，2號泵 &
17:06:02 & 進出料泵停止，70 mL/min，1號泵 \\ 
16:39:56 & 進出料泵開始，90 mL/min，1號泵 &
17:06:02 & 進出料泵開始，70 mL/min，2號泵 \\ 
16:39:57 & 進出料泵開始，90 mL/min，2號泵 &
17:06:23 & 進出料泵停止，70 mL/min，2號泵 \\ 
16:40:14 & 進出料泵停止，90 mL/min，1號泵 &
17:09:23 & 進出料泵開始，200 mL/min，1號泵 \\ 
16:40:30 & 進出料泵停止，90 mL/min，2號泵 &
17:09:23 & 進出料泵停止，200 mL/min，1號泵 \\ 
16:40:31 & 進出料泵開始，90 mL/min，1號泵 &
17:09:23 & 進出料泵開始，280 mL/min，2號泵 \\ 
16:40:51 & 進出料泵開始，90 mL/min，2號泵 &
17:14:18 & 進出料泵停止，90 mL/min，1號泵 \\ 
16:41:16 & 進出料泵停止，90 mL/min，1號泵 &
17:14:18 & 進出料泵開始，66.675 mL/min，1號泵 \\ 
16:41:16 & 進出料泵開始，100--200 mL/min &
17:14:18 & 進出料泵停止，66.675 mL/min，1號泵 \\ 
16:41:17 & 進出料泵停止，100--200 mL/min &
17:14:18 & 進出料泵開始，99.675 mL/min，2號泵 \\ 
16:41:22 & 進出料泵開始，150 mL/min，1號泵 &
17:14:18 & 進出料泵停止，99.675 mL/min，2號泵 \\ 
16:55:22 & 進出料泵停止，150 mL/min，1號泵 &
17:14:18 & 進出料泵開始，70 mL/min，2號泵 \\ 
16:55:48 & 進出料泵開始，150 mL/min，2號泵 &
17:14:18 & 進出料泵停止，70 mL/min，2號泵 \\ 
16:55:46 & 進出料泵停止，150 mL/min，2號泵 &
17:14:18 & 進出料泵開始，70 mL/min，1號泵 \\ 
16:55:46 & 進出料泵開始，150 mL/min，1號泵 &
17:14:18 & 進出料泵停止，70 mL/min，1號泵 \\ 
16:56:05 & 進出料泵停止，150 mL/min，1號泵 &
17:45:48 & 進出料泵開始，90 mL/min，1號泵 \\ 
16:56:05 & 進出料泵開始，150 mL/min，2號泵 &
17:45:48 & 進出料泵開始，90 mL/min，2號泵 \\ 
\end{tabularx}

% TODO: The source page contains a dense rotated equipment-event log; several
% Chinese event descriptions and flow-rate digits are difficult to resolve.
\end{minipage}%
}
% LEXOID_PAGE_COMPLETED: 16/78

\newpage

% TODO: This page is rotated in the source PDF. The visible content is a chronological operation log; the accompanying Chinese descriptions are very small and rotated, so they require verification against the original PDF.

% #VALUE_ID: LEX-P0017-V0001
% #HANDWRITTEN: 3/8
\handwritten{3/8}

\begin{tabularx}{\textwidth}{@{}lX@{}}
\textbf{时间} & \textbf{操作记录} \\ \hline
17:45:48 & \text{中文操作记录（原页为旋转文字，待核对）}\\
17:46:10 & \text{中文操作记录（原页为旋转文字，待核对）}\\
17:46:11 & \text{中文操作记录（原页为旋转文字，待核对）}\\
17:46:11 & \text{中文操作记录（原页为旋转文字，待核对）}\\
17:46:26 & \text{中文操作记录（原页为旋转文字，待核对）}\\
17:46:27 & \text{中文操作记录（原页为旋转文字，待核对）}\\
17:46:27 & \text{中文操作记录（原页为旋转文字，待核对）}\\
17:46:27 & \text{中文操作记录（原页为旋转文字，待核对）}\\
17:46:38 & \text{中文操作记录（原页为旋转文字，待核对）}\\
17:46:39 & \text{中文操作记录（原页为旋转文字，待核对）}\\
17:46:39 & \text{中文操作记录（原页为旋转文字，待核对）}\\
17:46:39 & \text{中文操作记录（原页为旋转文字，待核对）}\\
17:47:24 & \text{中文操作记录（原页为旋转文字，待核对）}\\
17:47:24 & \text{中文操作记录（原页为旋转文字，待核对）}\\
17:47:25 & \text{中文操作记录（原页为旋转文字，待核对）}\\
17:47:26 & \text{中文操作记录（原页为旋转文字，待核对）}\\
17:47:42 & \text{中文操作记录（原页为旋转文字，待核对）}\\
17:47:42 & \text{中文操作记录（原页为旋转文字，待核对）}\\
17:47:42 & \text{中文操作记录（原页为旋转文字，待核对）}\\
17:47:56 & \text{中文操作记录（原页为旋转文字，待核对）}\\
17:47:57 & \text{中文操作记录（原页为旋转文字，待核对）}\\
17:47:57 & \text{中文操作记录（原页为旋转文字，待核对）}\\
17:47:57 & \text{中文操作记录（原页为旋转文字，待核对）}\\
17:47:57 & \text{中文操作记录（原页为旋转文字，待核对）}\\
17:50:02 & \text{中文操作记录（原页为旋转文字，待核对）}\\
17:50:02 & \text{中文操作记录（原页为旋转文字，待核对）}\\
17:50:23 & \text{中文操作记录（原页为旋转文字，待核对）}\\
17:50:23 & \text{中文操作记录（原页为旋转文字，待核对）}\\
17:50:23 & \text{中文操作记录（原页为旋转文字，待核对）}\\
17:50:23 & \text{中文操作记录（原页为旋转文字，待核对）}\\
17:54:23 & \text{中文操作记录（原页为旋转文字，待核对）}\\
17:54:23 & \text{中文操作记录（原页为旋转文字，待核对）}\\
17:56:30 & \text{中文操作记录（原页为旋转文字，待核对）}\\
17:56:30 & \text{中文操作记录（原页为旋转文字，待核对）}\\
17:56:30 & \text{中文操作记录（原页为旋转文字，待核对）}\\
17:56:30 & \text{中文操作记录（原页为旋转文字，待核对）}\\
17:56:40 & \text{中文操作记录（原页为旋转文字，待核对）}\\
17:56:47 & \text{中文操作记录（原页为旋转文字，待核对）}\\
17:56:48 & \text{中文操作记录（原页为旋转文字，待核对）}\\
17:56:48 & \text{中文操作记录（原页为旋转文字，待核对）}\\
17:56:48 & \text{中文操作记录（原页为旋转文字，待核对）}\\
17:56:49 & \text{中文操作记录（原页为旋转文字，待核对）}\\
17:57:01 & \text{中文操作记录（原页为旋转文字，待核对）}\\
17:57:01 & \text{中文操作记录（原页为旋转文字，待核对）}\\
17:57:01 & \text{中文操作记录（原页为旋转文字，待核对）}\\
17:57:01 & \text{中文操作记录（原页为旋转文字，待核对）}\\
17:57:01 & \text{中文操作记录（原页为旋转文字，待核对）}\\
17:57:01 & \text{中文操作记录（原页为旋转文字，待核对）}\\
17:57:02 & \text{中文操作记录（原页为旋转文字，待核对）}\\
17:57:02 & \text{中文操作记录（原页为旋转文字，待核对）}\\
17:57:03 & \text{中文操作记录（原页为旋转文字，待核对）}\\
17:58:03 & \text{中文操作记录（原页为旋转文字，待核对）}\\
17:58:31 & \text{中文操作记录（原页为旋转文字，待核对）}\\
18:01:03 & \text{中文操作记录（原页为旋转文字，待核对）}\\
18:01:03 & \text{中文操作记录（原页为旋转文字，待核对）}\\
18:01:03 & \text{中文操作记录（原页为旋转文字，待核对）}\\
18:04:09 & \text{中文操作记录（原页为旋转文字，待核对）}\\
18:04:09 & \text{中文操作记录（原页为旋转文字，待核对）}\\
18:04:09 & \text{中文操作记录（原页为旋转文字，待核对）}\\
18:04:09 & \text{中文操作记录（原页为旋转文字，待核对）}\\
18:06:14 & \text{中文操作记录（原页为旋转文字，待核对）}\\
18:06:14 & \text{中文操作记录（原页为旋转文字，待核对）}\\
18:06:14 & \text{中文操作记录（原页为旋转文字，待核对）}\\
18:06:14 & \text{中文操作记录（原页为旋转文字，待核对）}\\
18:06:14 & \text{中文操作记录（原页为旋转文字，待核对）}\\
18:06:14 & \text{中文操作记录（原页为旋转文字，待核对）}\\
18:06:14 & \text{中文操作记录（原页为旋转文字，待核对）}\\
18:06:30 & \text{中文操作记录（原页为旋转文字，待核对）}\\
18:06:30 & \text{中文操作记录（原页为旋转文字，待核对）}\\
18:06:30 & \text{中文操作记录（原页为旋转文字，待核对）}\\
18:06:30 & \text{中文操作记录（原页为旋转文字，待核对）}\\
\end{tabularx}
% LEXOID_PAGE_COMPLETED: 17/78

\newpage

% TODO: This page is a low-resolution, sideways scan. The printed Chinese log text is only partially legible; the visible time entries and recognizable medication-related text are transcribed below.

\begin{center}
\rotatebox{90}{%
\begin{minipage}{0.92\textheight}
\small
\begin{tabularx}{\linewidth}{@{}lX@{}}
18:06:43 & 输液泵相关记录；静脉输液及滴速信息。\\
18:06:43 & 输液泵相关记录；静脉输液及滴速信息。\\
18:06:43 & 输液泵相关记录；静脉输液及滴速信息。\\
18:06:43 & 输液泵相关记录；静脉输液及滴速信息。\\
18:06:43 & 输液泵相关记录；静脉输液及滴速信息。\\
18:11:13 & 静脉输液记录。\\
18:11:13 & 静脉输液记录。\\
18:11:13 & 静脉输液记录。\\
18:13:54 & 静脉输液记录。\\
18:13:54 & 静脉输液记录。\\
18:15:57 & 输液记录。\\
18:15:57 & 输液记录。\\
18:16:57 & 输液记录。\\
18:16:57 & 输液记录。\\
18:16:12 & 输液记录。\\
18:16:12 & 输液记录。\\
18:16:25 & 输液记录。\\
18:16:25 & 输液记录。\\
18:17:56 & 输液记录。\\
18:17:56 & 输液记录。\\
18:20:56 & 输液记录。\\
18:26:34 & 输液记录。\\
18:27:35 & 输液记录。\\
\end{tabularx}

\vspace{1em}

\begin{tabularx}{\linewidth}{@{}lX@{}}
18:27:35 & 静脉输液；可见滴速、剂量及输液泵信息。\\
18:27:35 & 静脉输液；可见滴速、剂量及输液泵信息。\\
18:27:35 & 静脉输液；可见滴速、剂量及输液泵信息。\\
18:27:35 & 静脉输液；可见滴速、剂量及输液泵信息。\\
18:27:35 & 静脉输液；可见滴速、剂量及输液泵信息。\\
18:27:35 & 静脉输液；可见滴速、剂量及输液泵信息。\\
18:27:54 & 静脉输液；可见滴速、剂量及输液泵信息。\\
18:27:54 & 静脉输液；可见滴速、剂量及输液泵信息。\\
18:28:07 & 输液记录。\\
18:28:07 & 输液记录。\\
18:28:07 & 输液记录。\\
18:28:07 & 输液记录。\\
18:30:37 & 输液记录。\\
18:30:37 & 输液记录。\\
18:30:37 & 输液记录。\\
18:31:39 & 输液记录。\\
18:31:39 & 输液记录。\\
18:35:45 & 输液记录。\\
18:35:45 & 输液记录。\\
18:36:50 & 输液记录。\\
18:36:51 & 输液记录。\\
18:38:31 & 输液记录。\\
18:38:31 & 输液记录。\\
18:38:45 & 输液记录。\\
18:38:45 & 输液记录。\\
18:39:00 & 输液记录。\\
18:39:00 & 输液记录。\\
\end{tabularx}
\end{minipage}%
}
\end{center}

% #VALUE_ID: LEX-P0018-V0001
% #HANDWRITTEN: 4/8
\handwritten{4/8}
% LEXOID_PAGE_COMPLETED: 18/78

\newpage

% #VALUE_ID: LEX-P0019-V0001
% #TODO #HANDWRITTEN: illegible; small blue handwritten mark near the upper-right edge of the page
\handwritten{[illegible blue mark]}

\begin{center}
\scriptsize
\begin{tabularx}{\textwidth}{@{}>{\ttfamily}lX@{}}
18:39:00 & 流量平均值 500\,rpm；开始进料\\
18:41:00 & 加料完成；重量约 125\,g\\
18:44:00 & 加料完成；重量约 457\,g\\
18:45:00 & 加料完成；重量约 245\,g\\
18:45:00 & 加料完成；重量约 457\,g\\
18:48:26 & 进料平均流量 19.2\,mL；开始进料\\
18:48:27 & 加料完成；重量约 257\,g\\
18:48:27 & 加料完成；重量约 257\,g\\
18:48:32 & 加料完成；重量约 532\,g\\
18:51:12 & 加料完成；重量约 99.9\,mL\\
18:51:12 & 出料平均流量 70\,mL；第 2 罐\\
18:51:12 & 出料平均流量 70\,mL；第 2 罐\\
18:51:26 & 出料平均流量 39\,mL；第 2 罐\\
18:51:40 & 出料平均流量 39\,mL；第 2 罐\\
18:53:40 & 出料平均流量 39\,mL；第 2 罐\\
18:53:40 & 出料完成；重量约 99.9\,mL\\
18:57:40 & 加料完成；重量约 99.9\,mL\\
18:57:40 & 出料平均流量 70\,mL；第 2 罐\\
18:57:40 & 出料平均流量 70\,mL；第 2 罐\\
18:59:16 & 加料完成；重量约 19.8\,mL\\
19:01:17 & 加料完成；重量约 257\,g\\
19:01:17 & 出料平均流量 70\,mL；第 2 罐\\
19:01:22 & 出料平均流量 70\,mL；第 2 罐\\
19:01:36 & 加料完成；重量约 257\,g\\
19:01:37 & 加料完成；重量约 257\,g\\
19:01:49 & 加料完成；重量约 99.9\,mL\\[2pt]
19:02:00 & 出料平均流量 20\,mL；第 2 罐\\
19:02:00 & 出料平均流量 70\,mL；第 2 罐\\
19:02:00 & 加料完成；重量约 100\,mL\\
19:02:17 & 加料完成；重量约 257\,g\\
19:03:34 & 加料完成；重量约 457\,g\\
19:03:44 & 加料完成；重量约 457\,g\\
19:03:44 & 出料完成；重量约 99.9\,mL\\
19:04:00 & 出料平均流量 800\,rpm；第 2 罐\\
19:05:00 & 加料完成；重量约 257\,g\\
19:05:00 & 出料平均流量 70\,mL；第 2 罐\\
19:06:00 & 出料平均流量 70\,mL；第 2 罐\\
19:07:00 & 加料完成；重量约 99.9\,mL\\
19:08:00 & 出料平均流量 800\,rpm；第 2 罐\\
19:09:00 & 加料完成；重量约 257\,g\\
19:10:00 & 出料平均流量 800\,rpm；第 2 罐\\
19:11:00 & 加料完成；重量约 99.9\,mL\\
19:12:44 & 出料平均流量 70\,mL；第 2 罐\\
19:12:48 & 出料平均流量 70\,mL；第 2 罐\\
19:13:08 & 加料完成；重量约 99.9\,mL\\
19:15:08 & 加料完成；重量约 99.9\,mL\\
19:15:33 & 出料平均流量 800\,rpm；第 2 罐\\
19:15:39 & 加料完成；重量约 257\,g\\
19:16:47 & 出料平均流量 70\,mL；第 2 罐\\
19:17:44 & 加料完成；重量约 257\,g\\
19:18:47 & 出料平均流量 800\,rpm；第 2 罐\\
19:19:45 & 出料平均流量 70\,mL；第 2 罐\\
19:19:45 & 加料完成；重量约 257\,g\\
19:19:45 & 加料完成；重量约 99.9\,mL\\
19:21:44 & 出料平均流量 70\,mL；第 2 罐\\
19:21:44 & 加料完成；重量约 19.8\,mL\\
19:24:44 & 出料平均流量 70\,mL；第 2 罐\\
\end{tabularx}
\end{center}

% TODO: The source page contains a very small, vertically oriented process log. Several Chinese event descriptions and numerical readings are difficult to distinguish at the supplied resolution; the visible rows above preserve the apparent chronological structure only.
% LEXOID_PAGE_COMPLETED: 19/78

\newpage

\begin{center}
\textbf{蒸汽流量（t/h）}
\end{center}

\noindent
\begin{tabularx}{\textwidth}{@{}X X@{}}
\textbullet\ 蒸汽组1 &
\textbullet\ 蒸汽组2 \\
\textbullet\ 蒸汽组3 &
\textbullet\ 蒸汽组4
\end{tabularx}

\begin{center}
\begin{tabular}{c|c|c|c|c|c}
2040 & 1530 & 1020 & 510 & 0 & -510 \\
\hline
16{:}30/26 & & 18{:}00/26 & & 19{:}30/26 &
\end{tabular}

\vspace{0.5em}
% TODO: Reproduce the visible multi-series steam-flow chart; source filename is unavailable.
\fbox{\parbox[c][0.25\textheight][c]{0.95\textwidth}{\centering
\textit{蒸汽流量曲线图（图形待补）}}}
\end{center}

\begin{center}
\textbf{蒸汽压力（MPa）}
\end{center}

\noindent
\begin{tabularx}{\textwidth}{@{}X X@{}}
\textbullet\ 蒸汽组1 &
\textbullet\ 蒸汽组2 \\
\textbullet\ 蒸汽组3 &
\textbullet\ 蒸汽组4
\end{tabularx}

\begin{center}
\begin{tabular}{c|c|c|c|c|c}
480 & 360 & 240 & 120 & 0 & -120 \\
\hline
16{:}30/26 & & 18{:}00/26 & & 19{:}30/26 &
\end{tabular}

\vspace{0.5em}
% TODO: Reproduce the visible steam-pressure chart; source filename is unavailable.
\fbox{\parbox[c][0.25\textheight][c]{0.95\textwidth}{\centering
\textit{蒸汽压力曲线图（图形待补）}}}
\end{center}

\vfill

\noindent
% #VALUE_ID: LEX-P0020-V0001
% #HANDWRITTEN: 1/8
\handwritten{1/8}

\bigskip

\noindent
【工艺参数】

\medskip

\noindent
192/2145

\noindent
192-3-27

\medskip

\noindent
炉膛出口温度：
% #VALUE_ID: LEX-P0020-V0002
% #FIELD_VALUE: 炉膛出口温度
\fieldvalue{$375.7\,^{\circ}\mathrm{C}$}

\medskip

\noindent
签字：
% #VALUE_ID: LEX-P0020-V0003
% #FIELD_VALUE: 签字
% TODO #HANDWRITTEN: 难辨；蓝色手写签名无法可靠识别
\fieldvalue{\handwritten{难辨}}

\noindent
日期：
% #VALUE_ID: LEX-P0020-V0004
% #FIELD_VALUE: 日期
% #HANDWRITTEN: 8
\fieldvalue{\handwritten{8}}
% LEXOID_PAGE_COMPLETED: 20/78

\newpage

\begin{figure}[ht]
\centering
\fbox{\parbox{0.92\textwidth}{%
\centering
\textbf{流量 (mL/min)}\\[2pt]
\textit{処置1}（橙色）\qquad \textit{処置2}（緑色）\\[2pt]
縦軸目盛：0, 66, 132, 198, 264, 350\\
横軸目盛：16:30:26,\quad 18:00:26,\quad 19:30:26\\[4pt]
% TODO: Insert the upper plotted chart; no source filename is available.
}}
\caption{流量の推移}
\end{figure}

% #VALUE_ID: LEX-P0021-V0001
% #HANDWRITTEN: 716
\handwritten{716}

\begin{figure}[ht]
\centering
\fbox{\parbox{0.92\textwidth}{%
\centering
\textbf{重さ (g)}\\[2pt]
\textit{含水分}（橙色）\\[2pt]
縦軸目盛：0, 85, 170, 255, 340, 425\\
横軸目盛：16:30:26,\quad 18:00:26,\quad 19:30:26\\[4pt]
% TODO: Insert the lower plotted chart; no source filename is available.
}}
\caption{重さの推移}
\end{figure}
% LEXOID_PAGE_COMPLETED: 21/78

\newpage

% TODO: The source page contains two plotted graphs; no source filename is available for direct inclusion.
\begin{figure}[ht]
\centering
\fbox{\begin{minipage}{0.92\textwidth}
\centering
\textbf{Pressure plot}\\[2pt]
\small
Vertical axis: 压力 (kPa), with visible tick marks at
$-140$, $0$, $140$, $280$, $420$, and $560$.\\
Horizontal axis time marks: 16:30:26, 18:00:26, and 19:30:26.\\
The plot contains multiple colored traces and a legend with four entries
(label text partially unclear in the scan).
\end{minipage}}
\end{figure}

% TODO: The source page contains a plotted graph; no source filename is available for direct inclusion.
\begin{figure}[ht]
\centering
\fbox{\begin{minipage}{0.92\textwidth}
\centering
\textbf{Temperature plot}\\[2pt]
\small
Vertical axis: 温度 ($^\circ$C), with visible tick marks at
$-10$, $0$, $10$, $20$, $30$, and $40$.\\
Horizontal axis time marks: 16:30:26, 18:00:26, and 19:30:26.\\
The legend contains one orange trace labeled 温度.
\end{minipage}}
\end{figure}

% #VALUE_ID: LEX-P0022-V0001
% #TODO #HANDWRITTEN: 8/8; blue handwritten annotation at the upper right margin
\handwritten{8/8}

% #VALUE_ID: LEX-P0022-V0002
% #TODO #HANDWRITTEN: illegible Chinese text followed by 8/8; blue handwritten annotation at the right margin
\handwritten{[illegible Chinese text] 8/8}
% LEXOID_PAGE_COMPLETED: 22/78

\newpage

\section*{工艺记录}

% #VALUE_ID: LEX-P0023-V0001
% #FIELD_VALUE: 出炉日期
\fieldvalue{2025-06-18 19:28:00}

\begin{tabularx}{\textwidth}{@{}lX@{}}
工艺编号 & 
% #VALUE_ID: LEX-P0023-V0002
% #FIELD_VALUE: 工艺编号
\fieldvalue{Pro-202207300014} \\
工艺版本 & 
% #VALUE_ID: LEX-P0023-V0003
% #FIELD_VALUE: 工艺版本
\fieldvalue{V1.18} \\
产品编码 & 
% #VALUE_ID: LEX-P0023-V0004
% #FIELD_VALUE: 产品编码
\fieldvalue{202?} \\
产品名称 & 
% #VALUE_ID: LEX-P0023-V0005
% #FIELD_VALUE: 产品名称
\fieldvalue{年产5000吨聚醚} \\
原料批号 & 
% #VALUE_ID: LEX-P0023-V0006
% #FIELD_VALUE: 原料批号
\fieldvalue{2025-06-18 19:18:49} \\
\end{tabularx}

\vspace{1em}

\begin{tabularx}{\textwidth}{@{}l l X@{}}
\textbf{工艺步骤} & \textbf{时间} & \textbf{工艺参数} \\
\hline
投料 & 
% #VALUE_ID: LEX-P0023-V0007
% #FIELD_VALUE: 投料时间
\fieldvalue{16:14:59} &
% #VALUE_ID: LEX-P0023-V0008
% #FIELD_VALUE: 投料量
\fieldvalue{690.0\,mL} \\

投料 & 
% #VALUE_ID: LEX-P0023-V0009
% #FIELD_VALUE: 投料时间
\fieldvalue{16:15:05} &
% #VALUE_ID: LEX-P0023-V0010
% #FIELD_VALUE: 投料量
\fieldvalue{143.6\,mL} \\

投料 & 
% #VALUE_ID: LEX-P0023-V0011
% #FIELD_VALUE: 投料时间
\fieldvalue{16:15:52} &
% #VALUE_ID: LEX-P0023-V0012
% #FIELD_VALUE: 投料量
\fieldvalue{183.6\,mL} \\

投料 & 
% #VALUE_ID: LEX-P0023-V0013
% #FIELD_VALUE: 投料时间
\fieldvalue{16:20:18} &
% #VALUE_ID: LEX-P0023-V0014
% #FIELD_VALUE: 投料量
\fieldvalue{133.0\,mL} \\

投料 & 
% #VALUE_ID: LEX-P0023-V0015
% #FIELD_VALUE: 投料时间
\fieldvalue{16:26:05} &
% #VALUE_ID: LEX-P0023-V0016
% #FIELD_VALUE: 投料量
\fieldvalue{507.0\,mL} \\

投料 & 
% #VALUE_ID: LEX-P0023-V0017
% #FIELD_VALUE: 投料时间
\fieldvalue{16:24:53} &
% #VALUE_ID: LEX-P0023-V0018
% #FIELD_VALUE: 投料量
\fieldvalue{10344.4\,mL（2163\,mL）} \\

搅拌 & 
% #VALUE_ID: LEX-P0023-V0019
% #FIELD_VALUE: 搅拌时间
\fieldvalue{16:26:38} &
搅拌开始 \\

搅拌 & 
% #VALUE_ID: LEX-P0023-V0020
% #FIELD_VALUE: 搅拌时间
\fieldvalue{16:26:43} &
搅拌继续 \\

升温 & 
% #VALUE_ID: LEX-P0023-V0021
% #FIELD_VALUE: 升温时间
\fieldvalue{16:26:52} &
升温 \\

升温 & 
% #VALUE_ID: LEX-P0023-V0022
% #FIELD_VALUE: 升温时间
\fieldvalue{16:27:33} &
升温 \\

加热 & 
% #VALUE_ID: LEX-P0023-V0023
% #FIELD_VALUE: 加热时间
\fieldvalue{16:37:10} &
% #VALUE_ID: LEX-P0024-V0024
% #FIELD_VALUE: 温度
\fieldvalue{$90\,^{\circ}\mathrm{C}$} \\

加热 & 
% #VALUE_ID: LEX-P0023-V0025
% #FIELD_VALUE: 加热时间
\fieldvalue{16:37:12} &
% #VALUE_ID: LEX-P0023-V0026
% #FIELD_VALUE: 温度
\fieldvalue{$90\,^{\circ}\mathrm{C}$} \\

加热 & 
% #VALUE_ID: LEX-P0023-V0027
% #FIELD_VALUE: 加热时间
\fieldvalue{16:37:16} &
% #VALUE_ID: LEX-P0023-V0028
% #FIELD_VALUE: 温度
\fieldvalue{$90\,^{\circ}\mathrm{C}$} \\

加热 & 
% #VALUE_ID: LEX-P0023-V0029
% #FIELD_VALUE: 加热时间
\fieldvalue{16:37:21} &
% #VALUE_ID: LEX-P0023-V0030
% #FIELD_VALUE: 温度
\fieldvalue{$90\,^{\circ}\mathrm{C}$} \\

加热 & 
% #VALUE_ID: LEX-P0023-V0031
% #FIELD_VALUE: 加热时间
\fieldvalue{16:37:29} &
% #VALUE_ID: LEX-P0023-V0032
% #FIELD_VALUE: 温度
\fieldvalue{$90\,^{\circ}\mathrm{C}$} \\

加热 & 
% #VALUE_ID: LEX-P0023-V0033
% #FIELD_VALUE: 加热时间
\fieldvalue{16:37:32} &
% #VALUE_ID: LEX-P0023-V0034
% #FIELD_VALUE: 温度
\fieldvalue{$90\,^{\circ}\mathrm{C}$} \\

加热 & 
% #VALUE_ID: LEX-P0023-V0035
% #FIELD_VALUE: 加热时间
\fieldvalue{16:37:42} &
% #VALUE_ID: LEX-P0023-V0036
% #FIELD_VALUE: 温度
\fieldvalue{$90\,^{\circ}\mathrm{C}$} \\

加热 & 
% #VALUE_ID: LEX-P0023-V0037
% #FIELD_VALUE: 加热时间
\fieldvalue{16:37:51} &
% #VALUE_ID: LEX-P0023-V0038
% #FIELD_VALUE: 温度
\fieldvalue{$90\,^{\circ}\mathrm{C}$} \\

加热 & 
% #VALUE_ID: LEX-P0023-V0039
% #FIELD_VALUE: 加热时间
\fieldvalue{16:37:53} &
% #VALUE_ID: LEX-P0023-V0040
% #FIELD_VALUE: 温度
\fieldvalue{$90\,^{\circ}\mathrm{C}$} \\

加热 & 
% #VALUE_ID: LEX-P0023-V0041
% #FIELD_VALUE: 加热时间
\fieldvalue{16:37:54} &
% #VALUE_ID: LEX-P0023-V0042
% #FIELD_VALUE: 温度
\fieldvalue{$90\,^{\circ}\mathrm{C}$} \\
\end{tabularx}

\vspace{1em}

% #VALUE_ID: LEX-P0023-V0043
% #HANDWRITTEN: 11/06?
\handwritten{11/06?}

% #VALUE_ID: LEX-P0023-V0044
% #HANDWRITTEN: 2025 09?16
\handwritten{2025 09?16}

\noindent
\textbf{备注：} 自检通过

% TODO: Several small Chinese process labels and parameter descriptions are partially illegible in the source image; verify against the PDF at higher resolution.
% LEXOID_PAGE_COMPLETED: 23/78

\newpage

\rotatebox{-90}{%
\begin{minipage}{0.98\textheight}
\scriptsize
\begin{tabularx}{\textwidth}{@{}p{1.6cm}X@{}}
16:38:13 & 进出液体积：50.0mL，1号泵\\
16:38:13 & 进出液体积：600.0mL，2号泵\\
16:38:16 & 进出液体积：50.0mL，1号泵\\
16:38:33 & 进出液体积：50.0mL，1号泵\\
16:38:40 & 进出液体积：50.0mL，2号泵\\
16:38:40 & 进出液体积：50.0mL，2号泵\\
16:38:52 & 进出液体积：50.0mL，1号泵\\
16:38:56 & 进出液体积：50.0mL，1号泵\\
16:39:58 & 进出液体积：300mL/min，1号泵\\
16:40:14 & 进出液体积：450.0mL，1号泵\\
16:40:31 & 进出液体积：90mL/min，1号泵\\
16:40:31 & 进出液体积：16.0mL，1号泵\\
16:40:31 & 进出液体积：90mL/min，1号泵\\
16:40:31 & 进出液体积：16.0mL，1号泵\\
16:40:45 & 进出液体积：16.0mL，1号泵\\
16:41:16 & 进出液体积：16.0mL，1号泵\\
16:41:16 & 进出液体积：90mL/min，1号泵\\
16:41:21 & 进出液体积：16.0mL，1号泵\\
16:41:21 & 进出液体积：16.0mL，1号泵\\
16:55:21 & 进出液体积：16.0mL，1号泵\\
16:55:41 & 进出液体积：16.0mL，1号泵\\
16:55:41 & 进出液体积：16.0mL，1号泵\\
16:55:41 & 进出液体积：16.0mL，1号泵\\
16:55:41 & 进出液体积：16.0mL，1号泵\\
16:55:56 & 进出液体积：16.0mL，1号泵\\
16:55:56 & 进出液体积：16.0mL，1号泵\\[2pt]
16:55:25 & 进出液体积：50.0mL，1号泵\\
16:55:25 & 进出液体积：600.0mL，2号泵\\
16:55:30 & 进出液体积：50.0mL，1号泵\\
16:55:30 & 进出液体积：50.0mL，1号泵\\
16:55:35 & 进出液体积：50.0mL，2号泵\\
16:55:35 & 进出液体积：50.0mL，2号泵\\
16:55:40 & 进出液体积：50.0mL，1号泵\\
16:55:40 & 进出液体积：50.0mL，1号泵\\
16:55:55 & 进出液体积：300mL/min，1号泵\\
16:56:00 & 进出液体积：450.0mL，1号泵\\
16:56:05 & 进出液体积：90mL/min，1号泵\\
16:56:05 & 进出液体积：16.0mL，1号泵\\
16:56:10 & 进出液体积：90mL/min，1号泵\\
16:56:10 & 进出液体积：16.0mL，1号泵\\
17:44:06 & 进出液体积：16.0mL，1号泵\\
17:44:07 & 进出液体积：16.0mL，1号泵\\
\end{tabularx}

% TODO: The page contains a rotated Chinese process log. Several characters in
% the event descriptions are too small to distinguish reliably from the scan.
\end{minipage}%
}
% LEXOID_PAGE_COMPLETED: 24/78

\newpage

% TODO: The source page is rotated 90 degrees and contains a very small, densely printed event log. The time entries are transcribed below; the adjacent Chinese event descriptions are not reliably legible at the supplied resolution.

\begin{center}
\scriptsize
\begin{tabularx}{\textwidth}{@{}lX@{\qquad}lX@{}}
\textbf{时间} & \textbf{记录内容} & \textbf{时间} & \textbf{记录内容} \\
\hline
17:44:07 & \textit{（扫描文字过小，无法可靠辨识）} &
17:49:42 & \textit{（扫描文字过小，无法可靠辨识）} \\[1pt]
17:44:22 & \textit{（扫描文字过小，无法可靠辨识）} &
17:50:43 & \textit{（扫描文字过小，无法可靠辨识）} \\[1pt]
17:44:22 & \textit{（扫描文字过小，无法可靠辨识）} &
17:53:45 & \textit{（扫描文字过小，无法可靠辨识）} \\[1pt]
17:44:22 & \textit{（扫描文字过小，无法可靠辨识）} &
17:53:45 & \textit{（扫描文字过小，无法可靠辨识）} \\[1pt]
17:44:22 & \textit{（扫描文字过小，无法可靠辨识）} &
17:53:58 & \textit{（扫描文字过小，无法可靠辨识）} \\[1pt]
17:45:42 & \textit{（扫描文字过小，无法可靠辨识）} &
17:53:58 & \textit{（扫描文字过小，无法可靠辨识）} \\[1pt]
17:45:42 & \textit{（扫描文字过小，无法可靠辨识）} &
17:53:58 & \textit{（扫描文字过小，无法可靠辨识）} \\[1pt]
17:45:47 & \textit{（扫描文字过小，无法可靠辨识）} &
17:53:58 & \textit{（扫描文字过小，无法可靠辨识）} \\[1pt]
17:45:47 & \textit{（扫描文字过小，无法可靠辨识）} &
17:53:58 & \textit{（扫描文字过小，无法可靠辨识）} \\[1pt]
17:46:02 & \textit{（扫描文字过小，无法可靠辨识）} &
17:54:13 & \textit{（扫描文字过小，无法可靠辨识）} \\[1pt]
17:46:02 & \textit{（扫描文字过小，无法可靠辨识）} &
17:54:13 & \textit{（扫描文字过小，无法可靠辨识）} \\[1pt]
17:46:02 & \textit{（扫描文字过小，无法可靠辨识）} &
17:54:13 & \textit{（扫描文字过小，无法可靠辨识）} \\[1pt]
17:46:02 & \textit{（扫描文字过小，无法可靠辨识）} &
17:54:13 & \textit{（扫描文字过小，无法可靠辨识）} \\[1pt]
17:46:25 & \textit{（扫描文字过小，无法可靠辨识）} &
17:56:13 & \textit{（扫描文字过小，无法可靠辨识）} \\[1pt]
17:46:25 & \textit{（扫描文字过小，无法可靠辨识）} &
17:56:13 & \textit{（扫描文字过小，无法可靠辨识）} \\[1pt]
17:47:13 & \textit{（扫描文字过小，无法可靠辨识）} &
17:56:13 & \textit{（扫描文字过小，无法可靠辨识）} \\[1pt]
17:47:13 & \textit{（扫描文字过小，无法可靠辨识）} &
17:56:13 & \textit{（扫描文字过小，无法可靠辨识）} \\[1pt]
17:47:13 & \textit{（扫描文字过小，无法可靠辨识）} &
17:56:28 & \textit{（扫描文字过小，无法可靠辨识）} \\[1pt]
17:47:13 & \textit{（扫描文字过小，无法可靠辨识）} &
17:56:28 & \textit{（扫描文字过小，无法可靠辨识）} \\[1pt]
17:47:27 & \textit{（扫描文字过小，无法可靠辨识）} &
17:57:40 & \textit{（扫描文字过小，无法可靠辨识）} \\[1pt]
17:47:27 & \textit{（扫描文字过小，无法可靠辨识）} &
17:57:40 & \textit{（扫描文字过小，无法可靠辨识）} \\[1pt]
17:47:27 & \textit{（扫描文字过小，无法可靠辨识）} &
17:57:55 & \textit{（扫描文字过小，无法可靠辨识）} \\[1pt]
17:47:42 & \textit{（扫描文字过小，无法可靠辨识）} &
17:57:55 & \textit{（扫描文字过小，无法可靠辨识）} \\[1pt]
17:47:42 & \textit{（扫描文字过小，无法可靠辨识）} &
17:57:55 & \textit{（扫描文字过小，无法可靠辨识）} \\[1pt]
17:49:22 & \textit{（扫描文字过小，无法可靠辨识）} &
17:57:55 & \textit{（扫描文字过小，无法可靠辨识）} \\[1pt]
17:49:22 & \textit{（扫描文字过小，无法可靠辨识）} &
17:58:10 & \textit{（扫描文字过小，无法可靠辨识）} \\[1pt]
 & &
17:58:10 & \textit{（扫描文字过小，无法可靠辨识）} \\[1pt]
 & &
17:58:10 & \textit{（扫描文字过小，无法可靠辨识）} \\[1pt]
 & &
17:58:25 & \textit{（扫描文字过小，无法可靠辨识）} \\[1pt]
 & &
17:58:25 & \textit{（扫描文字过小，无法可靠辨识）} \\[1pt]
 & &
18:00:05 & \textit{（扫描文字过小，无法可靠辨识）} \\[1pt]
 & &
18:00:05 & \textit{（扫描文字过小，无法可靠辨识）} \\[1pt]
 & &
18:03:25 & \textit{（扫描文字过小，无法可靠辨识）} \\[1pt]
 & &
18:03:25 & \textit{（扫描文字过小，无法可靠辨识）} \\[1pt]
 & &
18:03:25 & \textit{（扫描文字过小，无法可靠辨识）} \\[1pt]
 & &
18:03:25 & \textit{（扫描文字过小，无法可靠辨识）} \\[1pt]
 & &
18:03:25 & \textit{（扫描文字过小，无法可靠辨识）} \\[1pt]
 & &
18:05:21 & \textit{（扫描文字过小，无法可靠辨识）} \\[1pt]
 & &
18:05:21 & \textit{（扫描文字过小，无法可靠辨识）} \\[1pt]
 & &
18:05:21 & \textit{（扫描文字过小，无法可靠辨识）} \\[1pt]
 & &
18:05:21 & \textit{（扫描文字过小，无法可靠辨识）} \\
\end{tabularx}
\end{center}
% LEXOID_PAGE_COMPLETED: 25/78

\newpage

\begin{flushright}
% #VALUE_ID: LEX-P0026-V0001
% #HANDWRITTEN: 4/8
\handwritten{4/8}
\end{flushright}

\noindent
\begin{tabularx}{\textwidth}{@{}lX@{\hspace{1.5em}}lX@{}}
\scriptsize
18:05:21 & 给药中；静脉泵：速度约 2.5\,mL/min & 18:16:42 & 给药中；静脉泵：速度约 2.5\,mL/min \\
18:05:21 & 输液；累计量约 0.7\,mL & 18:17:08 & 输液；累计量约 0.7\,mL \\
18:05:21 & 输液；累计量约 3.25\,mL & 18:17:42 & 输液；累计量约 3.25\,mL \\
18:05:36 & 注射器剩余量约 150\,mL；流速约 2\,mL/min & 18:19:24 & 注射器剩余量约 150\,mL；流速约 2\,mL/min \\
18:05:36 & 给药完成；静脉泵停止 & 18:20:08 & 给药完成；静脉泵停止 \\
18:05:36 & 设备状态正常 & 18:22:03 & 设备状态正常 \\
18:05:49 & 输液；累计量约 2.5\,mL & 18:22:08 & 输液；累计量约 2.5\,mL \\
18:05:49 & 输液；累计量约 3.25\,mL & 18:22:42 & 输液；累计量约 3.25\,mL \\
18:05:49 & 给药中；静脉泵：速度约 2.5\,mL/min & 18:23:45 & 给药中；静脉泵：速度约 2.5\,mL/min \\
18:07:19 & 给药中；静脉泵：速度约 2.5\,mL/min & 18:23:45 & 输液；累计量约 3.25\,mL \\
18:07:19 & 注射器剩余量约 150\,mL；流速约 2\,mL/min & 18:24:13 & 注射器剩余量约 150\,mL；流速约 2\,mL/min \\
18:07:19 & 给药完成；静脉泵停止 & 18:24:35 & 给药完成；静脉泵停止 \\
18:12:48 & 给药中；静脉泵：速度约 2.5\,mL/min & 18:24:35 & 给药中；静脉泵：速度约 2.5\,mL/min \\
18:12:49 & 输液；累计量约 2.5\,mL & 18:25:05 & 输液；累计量约 2.5\,mL \\
18:12:49 & 输液；累计量约 3.25\,mL & 18:25:09 & 输液；累计量约 3.25\,mL \\
18:12:49 & 注射器剩余量约 150\,mL；流速约 2\,mL/min & 18:26:05 & 注射器剩余量约 150\,mL；流速约 2\,mL/min \\
18:12:49 & 给药完成；静脉泵停止 & 18:29:05 & 给药完成；静脉泵停止 \\
18:14:19 & 给药中；静脉泵：速度约 2.5\,mL/min & 18:31:37 & 给药中；静脉泵：速度约 2.5\,mL/min \\
18:14:43 & 输液；累计量约 2.5\,mL & 18:33:37 & 输液；累计量约 2.5\,mL \\
18:14:43 & 输液；累计量约 3.25\,mL & 18:33:37 & 输液；累计量约 3.25\,mL \\
18:14:43 & 注射器剩余量约 150\,mL；流速约 2\,mL/min & 18:35:42 & 注射器剩余量约 150\,mL；流速约 2\,mL/min \\
18:14:57 & 给药完成；静脉泵停止 & 18:35:42 & 给药完成；静脉泵停止 \\
\end{tabularx}

% TODO: The source page is rotated and the log text is very small; several Chinese event descriptions may require verification against a higher-resolution scan.
% LEXOID_PAGE_COMPLETED: 26/78

\newpage

% #VALUE_ID: LEX-P0027-V0001
% #HANDWRITTEN: 4
\handwritten{4}

\small
\begin{tabularx}{\textwidth}{@{}p{0.18\textwidth}X@{}}
18:36:19 & 抽出液体，累计量约 410.2\,mL。\\
18:36:19 & 继续抽出液体，约 410.2\,mL。\\
18:36:33 & 抽出液体，累计量约 418.0\,mL。\\
18:36:33 & 继续抽出液体，约 418.0\,mL。\\
18:36:46 & 抽出液体，累计量约 459.5\,mL。\\
18:36:46 & 继续抽出液体，约 459.5\,mL。\\
18:36:46 & 抽出液体，累计量约 459.5\,mL。\\
18:42:46 & 抽出液体，累计量约 567.5\,mL。\\
18:42:46 & 继续抽出液体，约 567.5\,mL。\\
18:42:57 & 抽出液体，累计量约 475.0\,mL。\\
18:45:58 & 抽出液体，累计量约 475.0\,mL。\\
18:45:58 & 抽出液体，累计量约 475.0\,mL。\\
18:46:03 & 抽出液体，累计量约 475.0\,mL。\\
18:46:03 & 抽出液体，累计量约 475.0\,mL。\\
18:48:40 & 抽出液体，累计量约 401.5\,mL。\\
18:48:40 & 抽出液体，累计量约 401.5\,mL。\\
18:48:53 & 抽出液体，累计量约 700.2\,mL。\\
18:48:53 & 继续抽出液体，约 700.2\,mL。\\
18:49:05 & 抽出液体，累计量约 97.0\,mL。\\
18:49:07 & 抽出液体，累计量约 97.0\,mL。\\
18:49:07 & 继续抽出液体，约 97.0\,mL。\\
18:51:07 & 抽出液体，累计量约 150\,mL。\\
18:51:07 & 继续抽出液体，约 150\,mL。\\
18:55:07 & 抽出液体，累计量约 150\,mL。\\
18:55:07 & 继续抽出液体，约 150\,mL。\\
18:58:20 & 抽出液体，累计量约 198.0\,mL。\\
18:58:20 & 继续抽出液体，约 198.0\,mL。\\
\end{tabularx}

\vspace{1em}

\begin{tabularx}{\textwidth}{@{}p{0.18\textwidth}X@{}}
18:58:20 & 抽出液体，累计量约 150\,mL。\\
18:58:22 & 继续抽出液体，约 150\,mL。\\
18:58:24 & 抽出液体，累计量约 50\,mL。\\
18:58:44 & 继续抽出液体，约 50\,mL。\\
18:58:46 & 抽出液体，累计量约 90\,mL。\\
18:58:52 & 继续抽出液体，约 90\,mL。\\
18:59:04 & 抽出液体，累计量约 90\,mL。\\
19:00:03 & 抽出液体，累计量约 150\,mL。\\
19:00:08 & 继续抽出液体，约 150\,mL。\\
19:00:38 & 抽出液体，累计量约 300\,mL。\\
19:00:38 & 继续抽出液体，约 300\,mL。\\
19:06:38 & 抽出液体，累计量约 300\,mL。\\
19:06:48 & 继续抽出液体，约 300\,mL。\\
19:07:48 & 抽出液体，累计量约 300\,mL。\\
19:08:51 & 继续抽出液体，约 300\,mL。\\
19:09:52 & 抽出液体，累计量约 300\,mL。\\
19:09:52 & 继续抽出液体，约 300\,mL。\\
19:10:01 & 抽出液体，累计量约 150\,mL。\\
19:10:01 & 继续抽出液体，约 150\,mL。\\
19:12:50 & 抽出液体，累计量约 544.4\,mL。\\
19:12:50 & 继续抽出液体，约 544.4\,mL。\\
19:12:50 & 抽出液体，累计量约 544.4\,mL。\\
19:13:58 & 抽出液体，累计量约 150\,mL。\\
19:13:58 & 继续抽出液体，约 150\,mL。\\
19:15:58 & 抽出液体，累计量约 198.0\,mL。\\
19:15:58 & 继续抽出液体，约 198.0\,mL。\\
\end{tabularx}

% TODO: The source page is rotated and the small Chinese log descriptions are
% partially illegible; verify the transcribed entries against a higher-resolution scan.
% LEXOID_PAGE_COMPLETED: 27/78

\newpage

\section*{测试记录}

\begin{tabularx}{\textwidth}{@{}lX@{}}
时间 & 记录值 \\
\hline
% #VALUE_ID: LEX-P0028-V0001
% #FIELD_VALUE: 时间
\fieldvalue{19:15:58} & \textit{预设平均转速：800 rpm} \\

% #VALUE_ID: LEX-P0028-V0002
% #FIELD_VALUE: 时间
\fieldvalue{19:16:08} & \textit{温度记录：0.0 s} \\

% #VALUE_ID: LEX-P0028-V0003
% #FIELD_VALUE: 时间
\fieldvalue{19:17:04} & \textit{燃烧温度：1800 mm} \\

% #VALUE_ID: LEX-P0028-V0004
% #FIELD_VALUE: 时间
\fieldvalue{19:17:48} & \textit{温度记录：1600 mm} \\

% #VALUE_ID: LEX-P0028-V0005
% #FIELD_VALUE: 时间
\fieldvalue{19:17:48} & \textit{温度记录：165.1 mm} \\

% #VALUE_ID: LEX-P0028-V0006
% #FIELD_VALUE: 时间
\fieldvalue{19:18:08} & \textit{温度记录：165.1 mm} \\

% #VALUE_ID: LEX-P0028-V0007
% #FIELD_VALUE: 时间
\fieldvalue{19:18:49} & \textit{温度记录：4.1°C} \\

% #VALUE_ID: LEX-P0028-V0008
% #FIELD_VALUE: 时间
\fieldvalue{19:18:49} & \\

% #VALUE_ID: LEX-P0028-V0009
% #FIELD_VALUE: 时间
\fieldvalue{19:20:50} & \\
\end{tabularx}

\vspace{1em}

\noindent
签名：
% #VALUE_ID: LEX-P0028-V0010
% #FIELD_VALUE: 签名
% #TODO #HANDWRITTEN: 疑似“孙”；字迹不清
\fieldvalue{\handwritten{疑似“孙”}}

\medskip

\noindent
日期：
% #VALUE_ID: LEX-P0028-V0011
% #FIELD_VALUE: 日期
% #TODO #HANDWRITTEN: 6月？日；末位数字不清
\fieldvalue{\handwritten{6月？日}}

\vspace{1em}

\begin{figure}[ht]
\centering
% TODO: 当前页面包含上下排列的两幅温度曲线图；未提供可用图像文件名。
\fbox{\parbox{0.9\textwidth}{\centering
上图：测温点1--4的温度曲线\\[0.5em]
下图：测温点5--8的温度曲线
}}
\caption{温度记录曲线}
\end{figure}

% TODO: 页面中的曲线图及部分中文记录字段因扫描方向和分辨率限制，需依据原始 PDF 图像进一步核对。
% LEXOID_PAGE_COMPLETED: 28/78

\newpage

% #VALUE_ID: LEX-P0029-V0001
% #HANDWRITTEN: 18
\handwritten{18}

% TODO: Reproduce the two plotted charts visible on this page; no source image filename is available.
\begin{figure}[htbp]
    \centering
    \fbox{\parbox{0.92\textwidth}{\centering
        \textbf{Upper chart}\\[2pt]
        Chart with vertical axis labeled ``流量 (mL/min)'',\\
        tick marks from 0 to 330, and time labels approximately
        16:27:30, 17:57:30, and 19:27:30.\\
        Two plotted series are shown in orange and green.
    }}
\end{figure}

\begin{figure}[htbp]
    \centering
    \fbox{\parbox{0.92\textwidth}{\centering
        \textbf{Lower chart}\\[2pt]
        Chart with a vertical axis labeled ``電力 (g)'',\\
        tick marks from 0 to 435, and time labels approximately
        16:27:30, 17:57:30, and 19:27:30.\\
        An orange plotted series is shown.
    }}
\end{figure}
% LEXOID_PAGE_COMPLETED: 29/78

\newpage

% TODO: The page contains two rotated graph panels. No source image filename is available for direct inclusion.
\begin{figure}[p]
\centering
\fbox{\parbox[c][0.38\textheight][c]{0.92\textwidth}{\centering
\textbf{压力 (kPa)}\\[2pt]
Graph with four plotted series: \(\#1\), \(\#2\), \(\#3\), and \(\#4\).\\
Visible time axis: 16:27:30, 17:57:30, 19:27:30.\\
Visible vertical scale: \(-80, 0, 80, 160, 240, 320\).\\[4pt]
\textit{TODO: Replace this placeholder with the upper pressure graph.}
}}
\end{figure}

\begin{figure}[p]
\centering
\fbox{\parbox[c][0.38\textheight][c]{0.92\textwidth}{\centering
\textbf{温度 (°C)}\\[2pt]
Single plotted temperature curve.\\
Visible time axis: 16:27:30, 17:57:30, 19:27:30.\\
Visible vertical scale: \(-10, 0, 10, 20, 30, 40\).\\[4pt]
\textit{TODO: Replace this placeholder with the lower temperature graph.}
}}
\end{figure}

\begin{flushright}
% #VALUE_ID: LEX-P0030-V0001
% #HANDWRITTEN: 8/8
\handwritten{8/8}

% #VALUE_ID: LEX-P0030-V0002
% #TODO #HANDWRITTEN: 3/10?; final character is unclear
\handwritten{3/10?}

% #VALUE_ID: LEX-P0030-V0003
% #TODO #HANDWRITTEN: 2023.2.1?; handwritten date is partially illegible
\handwritten{2023.2.1?}
\end{flushright}
% LEXOID_PAGE_COMPLETED: 30/78

\newpage

\section*{职业健康检查表（22）}

\noindent
\textbf{职业病危害因素接触史及检查结果记录}

\begin{tabularx}{\textwidth}{|p{3.2cm}|X|}
\hline
检查日期 &
% #VALUE_ID: LEX-P0031-V0001
% #FIELD_VALUE: 检查日期
% #HANDWRITTEN: [日期，辨认不清]
\fieldvalue{\handwritten{\text{[日期，辨认不清]}}}
\quad
% #VALUE_ID: LEX-P0031-V0002
% #FIELD_VALUE: 检查日期附加记录
% #HANDWRITTEN: [中文手写姓名或备注，辨认不清]
\fieldvalue{\handwritten{\text{[中文手写姓名或备注，辨认不清]}}}
\\
\hline
采样日期 &
% #VALUE_ID: LEX-P0031-V0003
% #FIELD_VALUE: 采样日期
% #HANDWRITTEN: [日期，辨认不清]
\fieldvalue{\handwritten{\text{[日期，辨认不清]}}}
\\
\hline
报告日期 &
% #VALUE_ID: LEX-P0031-V0004
% #FIELD_VALUE: 报告日期
% #HANDWRITTEN: [日期，辨认不清]
\fieldvalue{\handwritten{\text{[日期，辨认不清]}}}
\\
\hline
职业病危害因素接触史 &
\\[-2pt]
\hline
\multicolumn{2}{|p{\dimexpr\textwidth-2\tabcolsep-2\arrayrulewidth\relax}|}{
\textbf{既往职业史及职业病危害因素接触情况：}

% #VALUE_ID: LEX-P0031-V0005
% #FIELD_VALUE: 职业史记录（一）
% #HANDWRITTEN: [手写记录，辨认不清]
\fieldvalue{\handwritten{\text{[手写记录，辨认不清]}}}

\par\smallskip
% #VALUE_ID: LEX-P0031-V0006
% #FIELD_VALUE: 职业史记录（二）
% #HANDWRITTEN: [手写记录，辨认不清]
\fieldvalue{\handwritten{\text{[手写记录，辨认不清]}}}

\par\smallskip
% #VALUE_ID: LEX-P0031-V0007
% #FIELD_VALUE: 职业史记录（三）
% #HANDWRITTEN: [手写记录，辨认不清]
\fieldvalue{\handwritten{\text{[手写记录，辨认不清]}}}
}
\\
\hline
检查结论 &
% #VALUE_ID: LEX-P0031-V0008
% #FIELD_VALUE: 检查结论
% #HANDWRITTEN: [手写结论，辨认不清]
\fieldvalue{\handwritten{\text{[手写结论，辨认不清]}}}
\\
\hline
\end{tabularx}

\medskip

\noindent
复核意见：\hrulefill

\medskip

\noindent
医师签名：\hrulefill

% #VALUE_ID: LEX-P0031-V0009
% #FIELD_VALUE: 医师签名或日期
% #HANDWRITTEN: [手写签名及日期，辨认不清]
\fieldvalue{\handwritten{\text{[手写签名及日期，辨认不清]}}}

\medskip

\noindent
备注：\hrulefill

% TODO：本页扫描图像倒置且分辨率较低，部分表头、日期及手写内容无法可靠辨认；表格仅保留本页可见区域。
% LEXOID_PAGE_COMPLETED: 31/78

\newpage

\begin{center}
\textbf{水利水电工程}

\medskip

\textbf{单位工程施工质量评定表（22）}
\end{center}

\bigskip

\noindent
\begin{tabularx}{\textwidth}{@{}lX@{}}
工程名称： & \hrulefill \\[1.2em]
单位工程名称： & \hrulefill \\[1.2em]
施工单位： &
% #VALUE_ID: LEX-P0032-V0001
% #FIELD_VALUE: 施工单位
% #HANDWRITTEN: [数字串，无法完全辨认]
\fieldvalue{\handwritten{[数字串，无法完全辨认]}} \\[1.2em]
编号： &
% #VALUE_ID: LEX-P0032-V0002
% #FIELD_VALUE: 编号
% #HANDWRITTEN: 200--480
\fieldvalue{\handwritten{200--480}} \\
\end{tabularx}

\bigskip

\noindent
评定日期：
% #VALUE_ID: LEX-P0032-V0003
% #FIELD_VALUE: 评定日期
% TODO #HANDWRITTEN: 202?年?月?日；蓝色手写日期辨认不清
\fieldvalue{\handwritten{202?年?月?日}}

\vfill

% TODO：页面下方可见印章、单位标识及其他固定页眉信息，文字过小且倒置，无法可靠辨认。
% LEXOID_PAGE_COMPLETED: 32/78

\newpage

% TODO: The source page is rotated and low-resolution; several labels and handwritten entries are only partially legible.

\begin{center}
\textbf{实验报告（22）}
\end{center}

\begin{tabularx}{\textwidth}{|l|X|l|X|l|X|}
\hline
实验名称 & \text{（表格内容部分不可辨识）} &
姓名 & \text{（未辨识）} &
学号 & 
% #VALUE_ID: LEX-P0033-V0001
% #FIELD_VALUE: 学号
\fieldvalue{200-048}
\\ \hline
\end{tabularx}

\medskip

\begin{tabularx}{\textwidth}{|c|X|c|}
\hline
序号 & 检查项目 & 检查结果 \\ \hline
1 & \text{实验前是否按要求预习及准备} &
% #VALUE_ID: LEX-P0033-V0002
% #FIELD_VALUE: 检查结果（第1项）
% #HANDWRITTEN: 否（勾选）
\fieldvalue{\handwritten{否（勾选）}}
\\ \hline
2 & \text{实验过程中是否遵守实验室规章制度} &
% #VALUE_ID: LEX-P0033-V0003
% #FIELD_VALUE: 检查结果（第2项）
% #HANDWRITTEN: 否（勾选）
\fieldvalue{\handwritten{否（勾选）}}
\\ \hline
3 & \text{是否按要求完成实验操作} &
% #VALUE_ID: LEX-P0033-V0004
% #FIELD_VALUE: 检查结果（第3项）
% #HANDWRITTEN: 否（勾选）
\fieldvalue{\handwritten{否（勾选）}}
\\ \hline
4 & \text{实验数据记录是否完整、准确} &
% #VALUE_ID: LEX-P0033-V0005
% #FIELD_VALUE: 检查结果（第4项）
% #HANDWRITTEN: 否（勾选）
\fieldvalue{\handwritten{否（勾选）}}
\\ \hline
5 & \text{实验报告是否按时完成} &
% #VALUE_ID: LEX-P0033-V0006
% #FIELD_VALUE: 检查结果（第5项）
% #HANDWRITTEN: 否（勾选）
\fieldvalue{\handwritten{否（勾选）}}
\\ \hline
6 & \text{实验结束后是否整理实验场地} &
% #VALUE_ID: LEX-P0033-V0007
% #FIELD_VALUE: 检查结果（第6项）
% #HANDWRITTEN: 否（勾选）
\fieldvalue{\handwritten{否（勾选）}}
\\ \hline
\end{tabularx}

\medskip

\begin{tabularx}{\textwidth}{|c|X|X|c|}
\hline
\multicolumn{4}{|c|}{实验过程检查记录} \\ \cline{1-4}
序号 & 检查内容 & 记录 & 结果 \\ \hline
1 & \text{实验前准备情况} &
% #VALUE_ID: LEX-P0033-V0008
% #FIELD_VALUE: 记录
% #TODO #HANDWRITTEN: 190-522（疑似编号或日期，字迹不清）
\fieldvalue{\handwritten{190-522}} &
% #VALUE_ID: LEX-P0033-V0009
% #FIELD_VALUE: 结果
% #HANDWRITTEN: 否（勾选）
\fieldvalue{\handwritten{否（勾选）}}
\\ \hline
2 & \text{实验操作及安全情况} &
% #VALUE_ID: LEX-P0033-V0010
% #FIELD_VALUE: 记录
% #TODO #HANDWRITTEN: 190-?（疑似编号，字迹不清）
\fieldvalue{\handwritten{190-?}} &
% #VALUE_ID: LEX-P0033-V0011
% #FIELD_VALUE: 结果
% #HANDWRITTEN: 否（勾选）
\fieldvalue{\handwritten{否（勾选）}}
\\ \hline
\end{tabularx}

\medskip

\begin{tabularx}{\textwidth}{|c|X|X|X|}
\hline
\multicolumn{4}{|c|}{实验数据及结果} \\ \cline{1-4}
项目 & 实验数据 & 计算结果 & 备注 \\ \hline
1 &
% #VALUE_ID: LEX-P0033-V0012
% #FIELD_VALUE: 实验数据
% #TODO #HANDWRITTEN: 200120（疑似数值，字迹不清）
\fieldvalue{\handwritten{200120}} &
% #VALUE_ID: LEX-P0033-V0013
% #FIELD_VALUE: 计算结果
% #TODO #HANDWRITTEN: 52?（字迹不清）
\fieldvalue{\handwritten{52?}} &
\text{（空白）}
\\ \hline
2 &
% #VALUE_ID: LEX-P0033-V0014
% #FIELD_VALUE: 实验数据
% #TODO #HANDWRITTEN: 190-?（字迹不清）
\fieldvalue{\handwritten{190-?}} &
% #VALUE_ID: LEX-P0033-V0015
% #FIELD_VALUE: 计算结果
% #TODO #HANDWRITTEN: 52?（字迹不清）
\fieldvalue{\handwritten{52?}} &
\text{（空白）}
\\ \hline
\end{tabularx}

% TODO: The bottom portion of the form continues beyond the visible page area.
% LEXOID_PAGE_COMPLETED: 33/78

\newpage

\begin{center}
\textbf{附件：设备检查记录表}
\end{center}

\noindent
审核日期：
% #VALUE_ID: LEX-P0034-V0001
% #FIELD_VALUE: 审核日期
% #HANDWRITTEN: 2025.06.19
\fieldvalue{\handwritten{2025.06.19}}
\hfill
检查日期：
% #VALUE_ID: LEX-P0034-V0002
% #FIELD_VALUE: 检查日期
% #HANDWRITTEN: 2025.06.19
\fieldvalue{\handwritten{2025.06.19}}

\vspace{0.5em}

\begin{tabularx}{\textwidth}{|c|X|X|}
\hline
序号 & 检查项目 & 检查结果及备注 \\
\hline
3 & 
% #TODO: 检查项目文字部分因图像方向及分辨率不清。
设备及其附件应齐全，外观应无明显损伤。 &
% #VALUE_ID: LEX-P0034-V0003
% #FIELD_VALUE: 第3项检查结果
% #HANDWRITTEN: 有
\fieldvalue{\handwritten{有}} \\
\hline
4 &
设备铭牌、型号及主要参数应与相关资料一致。 &
% #VALUE_ID: LEX-P0034-V0004
% #FIELD_VALUE: 第4项检查结果
% #HANDWRITTEN: 是
\fieldvalue{\handwritten{是}} \\
\hline
5 &
设备的外观、标识及接线应符合要求。 &
% #VALUE_ID: LEX-P0034-V0005
% #FIELD_VALUE: 第5项检查结果
% #HANDWRITTEN: 2025.06.19
\fieldvalue{\handwritten{2025.06.19}} \\
\hline
6 &
检查资料及随设备提供的文件是否齐全。 &
% #VALUE_ID: LEX-P0034-V0006
% #FIELD_VALUE: 第6项检查结果
% #HANDWRITTEN: 有
\fieldvalue{\handwritten{有}} \\
\hline
7 &
设备运行、操作及维护要求应明确。 &
% #VALUE_ID: LEX-P0034-V0007
% #FIELD_VALUE: 第7项检查结果
% #HANDWRITTEN: 2025.06.19
\fieldvalue{\handwritten{2025.06.19}} \\
\hline
8 &
% #TODO: 表内第8项文字无法从当前页面清晰辨识。
检查项目及验收要求。 &
% #VALUE_ID: LEX-P0034-V0008
% #FIELD_VALUE: 第8项检查结果
% #HANDWRITTEN: 合格
\fieldvalue{\handwritten{合格}} \\
\hline
9 &
% #TODO: 表内第9项文字无法从当前页面清晰辨识。
检查项目及验收要求。 &
% #VALUE_ID: LEX-P0034-V0009
% #FIELD_VALUE: 第9项检查结果
% #HANDWRITTEN: 合格
\fieldvalue{\handwritten{合格}} \\
\hline
10 &
% #TODO: 表内第10项文字无法从当前页面清晰辨识。
检查项目及验收要求。 &
% #VALUE_ID: LEX-P0034-V0010
% #FIELD_VALUE: 第10项检查结果
% #HANDWRITTEN: 合格
\fieldvalue{\handwritten{合格}} \\
\hline
11 &
% #TODO: 表内第11项文字无法从当前页面清晰辨识。
检查项目及验收要求。 &
% #VALUE_ID: LEX-P0034-V0011
% #FIELD_VALUE: 第11项检查结果
% #HANDWRITTEN: 合格
\fieldvalue{\handwritten{合格}} \\
\hline
12 &
% #TODO: 表内第12项文字无法从当前页面清晰辨识。
检查项目及验收要求。 &
% #VALUE_ID: LEX-P0034-V0012
% #FIELD_VALUE: 第12项检查结果
% #HANDWRITTEN: 合格
\fieldvalue{\handwritten{合格}} \\
\hline
\end{tabularx}

\vspace{1em}

\begin{tabularx}{\textwidth}{|c|X|X|X|}
\hline
序号 & 设备名称 & 规格型号 & 备注 \\
\hline
3 &
% #VALUE_ID: LEX-P0034-V0013
% #FIELD_VALUE: 设备名称
% #TODO #HANDWRITTEN: 不清楚；蓝色手写内容分辨率不足
\fieldvalue{\handwritten{不清楚}} &
% #VALUE_ID: LEX-P0034-V0014
% #FIELD_VALUE: 规格型号
% #TODO #HANDWRITTEN: 不清楚；蓝色手写内容分辨率不足
\fieldvalue{\handwritten{不清楚}} &
% #VALUE_ID: LEX-P0034-V0015
% #FIELD_VALUE: 备注
% #TODO #HANDWRITTEN: 不清楚；蓝色手写内容分辨率不足
\fieldvalue{\handwritten{不清楚}} \\
\hline
\end{tabularx}

% TODO: 本页表格内容疑似承接前页或继续至后页；仅保留当前页面可见部分。
% LEXOID_PAGE_COMPLETED: 34/78

\newpage

\section*{（表单标题，旋转/字迹较难辨认）}
% TODO：本页扫描图像为倒置页面，部分标题及表格文字无法可靠辨识。

\noindent
\textbf{受理编号：}%
% #VALUE_ID: LEX-P0035-V0001
% #FIELD_VALUE: 受理编号
\fieldvalue{XZ-2024-01-02-04}
\hfill
\textbf{日期：}%
% #VALUE_ID: LEX-P0035-V0002
% #FIELD_VALUE: 日期
% #HANDWRITTEN: 2024.11.15
\fieldvalue{\handwritten{2024.11.15}}

\begin{tabularx}{\textwidth}{|p{2.3cm}|X|p{2.3cm}|X|}
\hline
项目 &
填写内容 &
项目 &
填写内容 \\
\hline
姓名 &
% #VALUE_ID: LEX-P0035-V0003
% #FIELD_VALUE: 姓名
% #HANDWRITTEN: 难辨
\fieldvalue{\handwritten{难辨}} &
证件号码 &
% #VALUE_ID: LEX-P0035-V0004
% #FIELD_VALUE: 证件号码
% #HANDWRITTEN: 难辨
\fieldvalue{\handwritten{难辨}} \\
\hline
申请日期 &
% #VALUE_ID: LEX-P0035-V0005
% #FIELD_VALUE: 申请日期
% #HANDWRITTEN: 2024.10.15
\fieldvalue{\handwritten{2024.10.15}} &
联系电话 &
% #VALUE_ID: LEX-P0035-V0006
% #FIELD_VALUE: 联系电话
% #HANDWRITTEN: 难辨
\fieldvalue{\handwritten{难辨}} \\
\hline
票据/凭证号码 &
% #VALUE_ID: LEX-P0035-V0007
% #FIELD_VALUE: 票据/凭证号码
% #HANDWRITTEN: 难辨
\fieldvalue{\handwritten{难辨}} &
金额 &
% #VALUE_ID: LEX-P0035-V0008
% #FIELD_VALUE: 金额
% #HANDWRITTEN: 12
\fieldvalue{\handwritten{12}} \\
\hline
备注 &
% #VALUE_ID: LEX-P0035-V0009
% #FIELD_VALUE: 备注
% #HANDWRITTEN: 难辨
\fieldvalue{\handwritten{难辨}} &
审核意见 &
% #VALUE_ID: LEX-P0035-V0010
% #FIELD_VALUE: 审核意见
% #HANDWRITTEN: 难辨
\fieldvalue{\handwritten{难辨}} \\
\hline
\end{tabularx}

\vspace{1em}

\noindent
\textbf{票据/证明材料：}

\begin{center}
% TODO：此处为粘贴或扫描的票据/凭证，原图中文字及条形码可见，但无法在当前分辨率下完整辨识。
\fbox{\parbox[c][7.5cm][c]{0.86\textwidth}{\centering
票据/凭证扫描件\\[0.5em]
（含条形码及印章）}}
\end{center}

\vspace{0.5em}

\noindent
附注：页面中可见蓝色手写批注及数字，部分内容因原始页面倒置且分辨率不足而无法准确转录。

% #VALUE_ID: LEX-P0035-V0011
% #HANDWRITTEN: 难辨
\handwritten{难辨}

% TODO：页面底部/背面区域存在多处蓝色手写文字、日期及数字；由于其与倒置的表格和票据重叠，无法可靠区分为独立字段。
% LEXOID_PAGE_COMPLETED: 35/78

\newpage

\section*{职业技能鉴定记录}

\begin{tabularx}{\textwidth}{|l|X|}
\hline
职业资格证书编号 &
% #VALUE_ID: LEX-P0036-V0001
% #FIELD_VALUE: 证书编号
% #TODO #HANDWRITTEN: 不可辨认；蓝色手写编号
\fieldvalue{\handwritten{[不可辨认]}} \\
\hline
姓名 &
% #VALUE_ID: LEX-P0036-V0002
% #FIELD_VALUE: 姓名
% #TODO #HANDWRITTEN: 不可辨认；个人信息栏中的手写姓名
\fieldvalue{\handwritten{[不可辨认]}} \\
\hline
性别 &
% #VALUE_ID: LEX-P0036-V0003
% #FIELD_VALUE: 性别
\fieldvalue{男} \\
\hline
民族 &
% #VALUE_ID: LEX-P0036-V0004
% #FIELD_VALUE: 民族
\fieldvalue{汉族} \\
\hline
出生日期 &
% #VALUE_ID: LEX-P0036-V0005
% #FIELD_VALUE: 出生日期
% #TODO #HANDWRITTEN: 不可辨认；日期栏中的手写内容
\fieldvalue{\handwritten{[不可辨认]}} \\
\hline
\end{tabularx}

\vspace{0.5em}

\begin{tabularx}{\textwidth}{|X|c|}
\hline
鉴定项目 & 鉴定结果 \\
\hline
% #VALUE_ID: LEX-P0036-V0006
% #FIELD_VALUE: 鉴定项目一结果
% #HANDWRITTEN: ✓
职业道德 & \fieldvalue{\handwritten{\checkmark}} 合格 \\
\hline
% #VALUE_ID: LEX-P0036-V0007
% #FIELD_VALUE: 鉴定项目二结果
% #HANDWRITTEN: ✓
基础知识 & \fieldvalue{\handwritten{\checkmark}} 合格 \\
\hline
% #VALUE_ID: LEX-P0036-V0008
% #FIELD_VALUE: 鉴定项目三结果
% #HANDWRITTEN: ✓
专业知识 & \fieldvalue{\handwritten{\checkmark}} 合格 \\
\hline
% #VALUE_ID: LEX-P0036-V0009
% #FIELD_VALUE: 鉴定项目四结果
% #HANDWRITTEN: ✓
操作技能 & \fieldvalue{\handwritten{\checkmark}} 合格 \\
\hline
% #VALUE_ID: LEX-P0036-V0010
% #FIELD_VALUE: 鉴定项目五结果
% #HANDWRITTEN: ✓
综合评审 & \fieldvalue{\handwritten{\checkmark}} 合格 \\
\hline
% #VALUE_ID: LEX-P0036-V0011
% #FIELD_VALUE: 鉴定项目六结果
% #HANDWRITTEN: ✓
相关技能 & \fieldvalue{\handwritten{\checkmark}} 合格 \\
\hline
% #VALUE_ID: LEX-P0036-V0012
% #FIELD_VALUE: 鉴定项目七结果
% #HANDWRITTEN: ✓
其他项目 & \fieldvalue{\handwritten{\checkmark}} 合格 \\
\hline
\end{tabularx}

\vspace{0.5em}

\begin{tabularx}{\textwidth}{|l|X|}
\hline
鉴定成绩 &
% #VALUE_ID: LEX-P0036-V0013
% #FIELD_VALUE: 鉴定成绩
% #TODO #HANDWRITTEN: 不可辨认；成绩栏中的蓝色手写数字
\fieldvalue{\handwritten{[不可辨认]}} \\
\hline
鉴定日期 &
% #VALUE_ID: LEX-P0036-V0014
% #FIELD_VALUE: 鉴定日期
% #HANDWRITTEN: 2015/09/19
\fieldvalue{\handwritten{2015/09/19}} \\
\hline
鉴定机构意见 &
% #VALUE_ID: LEX-P0036-V0015
% #FIELD_VALUE: 鉴定机构意见
% #TODO #HANDWRITTEN: 不可辨认；意见栏中的蓝色手写内容
\fieldvalue{\handwritten{[不可辨认]}} \\
\hline
\end{tabularx}

\vspace{0.5em}

\noindent
鉴定机构（盖章）：\hfill 日期：\underline{\hspace{3cm}}

\vspace{1em}

\noindent
备注：本页内容为鉴定记录及结果登记栏。% TODO: 页面底部印章及部分个人信息因图像倒置且分辨率不足，无法完整辨识。
% LEXOID_PAGE_COMPLETED: 36/78

\newpage

\begin{center}
\textbf{医疗机构相关信息登记表}
\end{center}

\begin{tabularx}{\textwidth}{|l|X|l|X|}
\hline
医疗机构名称 & 
% #VALUE_ID: LEX-P0037-V0001
% #FIELD_VALUE: 医疗机构名称
% #HANDWRITTEN: （字迹难以辨认）
\fieldvalue{\handwritten{（字迹难以辨认）}} &
填表日期 &
% #VALUE_ID: LEX-P0037-V0002
% #FIELD_VALUE: 填表日期
% #HANDWRITTEN: 2025.06.09（疑似）
\fieldvalue{\handwritten{2025.06.09（疑似）}} \\
\hline
\multicolumn{4}{|l|}{
% #VALUE_ID: LEX-P0037-V0003
% #FIELD_VALUE: 编号
\fieldvalue{1}
} \\
\hline
\end{tabularx}

\vspace{0.5em}

\begin{tabularx}{\textwidth}{|p{0.28\textwidth}|X|}
\hline
项目 & 检查内容及要求 \\
\hline
一、基本情况 &
\textit{（表内文字部分因扫描方向及分辨率较低，部分内容难以辨认。）} \\
\hline
检查项目一 &
是否符合要求：\quad 是 $\square$\quad 否 $\square$ \\
\hline
检查项目二 &
是否符合要求：\quad 是 $\square$\quad 否 $\square$ \\
\hline
检查项目三 &
是否符合要求：\quad 是 $\square$\quad 否 $\square$ \\
\hline
检查项目四 &
是否符合要求：\quad 是 $\square$\quad 否 $\square$ \\
\hline
检查项目五 &
是否符合要求：\quad 是 $\square$\quad 否 $\square$ \\
\hline
检查项目六 &
是否符合要求：\quad 是 $\square$\quad 否 $\square$ \\
\hline
备注 &
% #VALUE_ID: LEX-P0037-V0004
% #FIELD_VALUE: 备注
% #HANDWRITTEN: （字迹难以辨认）
\fieldvalue{\handwritten{（字迹难以辨认）}} \\
\hline
\end{tabularx}

\vspace{0.5em}

\noindent
\textbf{2、检查结果及整改情况}

\begin{tabularx}{\textwidth}{|p{0.22\textwidth}|X|}
\hline
项目 & 结果及说明 \\
\hline
检查结果 &
% #VALUE_ID: LEX-P0037-V0005
% #FIELD_VALUE: 检查结果
% #HANDWRITTEN: （蓝色手写内容难以辨认）
\fieldvalue{\handwritten{（蓝色手写内容难以辨认）}} \\
\hline
整改情况 &
% #VALUE_ID: LEX-P0037-V0006
% #FIELD_VALUE: 整改情况
% #HANDWRITTEN: （蓝色手写内容难以辨认）
\fieldvalue{\handwritten{（蓝色手写内容难以辨认）}} \\
\hline
\end{tabularx}

\vspace{0.75em}

\begin{tabularx}{\textwidth}{|c|X|c|c|c|}
\hline
序号 & 药品名称 & 规格 & 单位 & 数量 \\
\hline
1 & \textit{（文字难以辨认）} & \textit{（难以辨认）} &
\textit{（难以辨认）} & \textit{（难以辨认）} \\
\hline
2 & 
% #VALUE_ID: LEX-P0037-V0007
% #FIELD_VALUE: 药品名称
% #HANDWRITTEN: （蓝色手写药品名称难以辨认）
\fieldvalue{\handwritten{（蓝色手写药品名称难以辨认）}} &
\textit{（难以辨认）} & \textit{（难以辨认）} &
% #VALUE_ID: LEX-P0037-V0008
% #FIELD_VALUE: 数量
% #HANDWRITTEN: （蓝色手写数字难以辨认）
\fieldvalue{\handwritten{（蓝色手写数字难以辨认）}} \\
\hline
\end{tabularx}

% TODO: 本页扫描图像为倒置页面，且正文及手写内容分辨率不足；需依据更高清晰度原图校正标题、表格字段及各项手写值。
% LEXOID_PAGE_COMPLETED: 37/78

\newpage

\section*{受试者相关记录}

\begin{tabularx}{\textwidth}{|l|X|l|X|}
\hline
姓名 & 
% #VALUE_ID: LEX-P0038-V0001
% #FIELD_VALUE: 姓名
% #TODO #HANDWRITTEN: 2025...; handwriting is partially illegible
\fieldvalue{\handwritten{2025...}} &
性别 & 
% #VALUE_ID: LEX-P0038-V0002
% #FIELD_VALUE: 性别
% #TODO #HANDWRITTEN: 不详; handwritten entry is unclear
\fieldvalue{\handwritten{不详}} \\
\hline
年龄 & 
% #VALUE_ID: LEX-P0038-V0003
% #FIELD_VALUE: 年龄
% #HANDWRITTEN: 49
\fieldvalue{\handwritten{49}} &
受试者编号 & 
% #VALUE_ID: LEX-P0038-V0004
% #FIELD_VALUE: 受试者编号
% #TODO #HANDWRITTEN: 不详; red handwritten/stamped value is illegible
\fieldvalue{\handwritten{不详}} \\
\hline
\end{tabularx}

\vspace{0.5em}

\section*{2.1.1 实验室检查/检验}

\begin{tabularx}{\textwidth}{|c|X|X|X|c|}
\hline
序号 & 检查名称 & 观察结果 & 结果 & 是否异常 \\
\hline
1 & 心电图检查 & 
% #VALUE_ID: LEX-P0038-V0005
% #FIELD_VALUE: 观察结果，第1项
% #TODO #HANDWRITTEN: 10/11; handwriting is partially illegible
\fieldvalue{\handwritten{10/11}} &
% #VALUE_ID: LEX-P0038-V0006
% #FIELD_VALUE: 是否异常，第1项
% #HANDWRITTEN: \checkmark
\fieldvalue{\handwritten{\checkmark}} & 否 \\
\hline
2 & 实验室检查 & 
% #VALUE_ID: LEX-P0038-V0007
% #FIELD_VALUE: 观察结果，第2项
% #TODO #HANDWRITTEN: 10/11; handwriting is partially illegible
\fieldvalue{\handwritten{10/11}} &
% #VALUE_ID: LEX-P0038-V0008
% #FIELD_VALUE: 是否异常，第2项
% #HANDWRITTEN: \checkmark
\fieldvalue{\handwritten{\checkmark}} & 否 \\
\hline
3 & 实验室检查 & 
% #VALUE_ID: LEX-P0038-V0009
% #FIELD_VALUE: 观察结果，第3项
% #TODO #HANDWRITTEN: 490/11; handwriting is partially illegible
\fieldvalue{\handwritten{490/11}} &
% #VALUE_ID: LEX-P0038-V0010
% #FIELD_VALUE: 是否异常，第3项
% #HANDWRITTEN: \checkmark
\fieldvalue{\handwritten{\checkmark}} & 否 \\
\hline
4 & 实验室检查 & 
% #VALUE_ID: LEX-P0038-V0011
% #FIELD_VALUE: 观察结果，第4项
% #HANDWRITTEN: /
\fieldvalue{\handwritten{/}} &
% #VALUE_ID: LEX-P0038-V0012
% #FIELD_VALUE: 是否异常，第4项
% #HANDWRITTEN: \checkmark
\fieldvalue{\handwritten{\checkmark}} & 否 \\
\hline
5 & 实验室检查 & 
% #VALUE_ID: LEX-P0038-V0013
% #FIELD_VALUE: 观察结果，第5项
% #TODO #HANDWRITTEN: 11/??; handwriting is unclear
\fieldvalue{\handwritten{11/??}} &
% #VALUE_ID: LEX-P0038-V0014
% #FIELD_VALUE: 是否异常，第5项
% #HANDWRITTEN: \checkmark
\fieldvalue{\handwritten{\checkmark}} & 否 \\
\hline
6 & 实验室检查 & 
% #VALUE_ID: LEX-P0038-V0015
% #FIELD_VALUE: 观察结果，第6项
% #TODO #HANDWRITTEN: 111/??; handwriting is unclear
\fieldvalue{\handwritten{111/??}} &
% #VALUE_ID: LEX-P0038-V0016
% #FIELD_VALUE: 是否异常，第6项
% #HANDWRITTEN: \checkmark
\fieldvalue{\handwritten{\checkmark}} & 否 \\
\hline
7 & 实验室检查 & 
% #VALUE_ID: LEX-P0038-V0017
% #FIELD_VALUE: 观察结果，第7项
% #TODO #HANDWRITTEN: 不详; handwriting is illegible
\fieldvalue{\handwritten{不详}} &
% #VALUE_ID: LEX-P0038-V0018
% #FIELD_VALUE: 是否异常，第7项
% #HANDWRITTEN: \checkmark
\fieldvalue{\handwritten{\checkmark}} & 否 \\
\hline
8 & 实验室检查 & 
% #VALUE_ID: LEX-P0038-V0019
% #FIELD_VALUE: 观察结果，第8项
% #TODO #HANDWRITTEN: 不详; handwriting is illegible
\fieldvalue{\handwritten{不详}} &
% #VALUE_ID: LEX-P0038-V0020
% #FIELD_VALUE: 是否异常，第8项
% #HANDWRITTEN: \checkmark
\fieldvalue{\handwritten{\checkmark}} & 否 \\
\hline
9 & 实验室检查 & 
% #VALUE_ID: LEX-P0038-V0021
% #FIELD_VALUE: 观察结果，第9项
% #TODO #HANDWRITTEN: 不详; handwriting is illegible
\fieldvalue{\handwritten{不详}} &
% #VALUE_ID: LEX-P0038-V0022
% #FIELD_VALUE: 是否异常，第9项
% #HANDWRITTEN: \checkmark
\fieldvalue{\handwritten{\checkmark}} & 否 \\
\hline
10 & 实验室检查 & 
% #VALUE_ID: LEX-P0038-V0023
% #FIELD_VALUE: 观察结果，第10项
% #TODO #HANDWRITTEN: 不详; handwriting is illegible
\fieldvalue{\handwritten{不详}} &
% #VALUE_ID: LEX-P0038-V0024
% #FIELD_VALUE: 是否异常，第10项
% #HANDWRITTEN: \checkmark
\fieldvalue{\handwritten{\checkmark}} & 否 \\
\hline
\end{tabularx}

% TODO: The page is a rotated scan and several table labels and handwritten entries are not fully legible.

\section*{2.1.2 实验室检查结果}

\begin{tabularx}{\textwidth}{|c|X|X|X|c|}
\hline
序号 & 检查项目 & 检查方法 & 检查结果 & 是否异常 \\
\hline
1 & 实验室检查 & CRP、血常规等 & 
% #VALUE_ID: LEX-P0038-V0025
% #FIELD_VALUE: 检查结果，第1项
% #TODO #HANDWRITTEN: 不详; blue handwriting is illegible
\fieldvalue{\handwritten{不详}} &
% #VALUE_ID: LEX-P0038-V0026
% #FIELD_VALUE: 是否异常，第1项
% #HANDWRITTEN: \checkmark
\fieldvalue{\handwritten{\checkmark}} \\
\hline
2 & 实验室检查 & CRP、血常规等 & 
% #VALUE_ID: LEX-P0038-V0027
% #FIELD_VALUE: 检查结果，第2项
% #TODO #HANDWRITTEN: 不详; blue handwriting is illegible
\fieldvalue{\handwritten{不详}} &
% #VALUE_ID: LEX-P0038-V0028
% #FIELD_VALUE: 是否异常，第2项
% #HANDWRITTEN: \checkmark
\fieldvalue{\handwritten{\checkmark}} \\
\hline
3 & 实验室检查 & CRP、血常规等 & 
% #VALUE_ID: LEX-P0038-V0029
% #FIELD_VALUE: 检查结果，第3项
% #TODO #HANDWRITTEN: 不详; blue handwriting is illegible
\fieldvalue{\handwritten{不详}} &
% #VALUE_ID: LEX-P0038-V0030
% #FIELD_VALUE: 是否异常，第3项
% #HANDWRITTEN: \checkmark
\fieldvalue{\handwritten{\checkmark}} \\
\hline
4 & 实验室检查 & CRP、血常规等 & 
% #VALUE_ID: LEX-P0038-V0031
% #FIELD_VALUE: 检查结果，第4项
% #TODO #HANDWRITTEN: 不详; blue handwriting is illegible
\fieldvalue{\handwritten{不详}} &
% #VALUE_ID: LEX-P0038-V0032
% #FIELD_VALUE: 是否异常，第4项
% #HANDWRITTEN: \checkmark
\fieldvalue{\handwritten{\checkmark}} \\
\hline
5 & 实验室检查 & CRP、血常规等 & 
% #VALUE_ID: LEX-P0038-V0033
% #FIELD_VALUE: 检查结果，第5项
% #TODO #HANDWRITTEN: 不详; blue handwriting is illegible
\fieldvalue{\handwritten{不详}} &
% #VALUE_ID: LEX-P0038-V0034
% #FIELD_VALUE: 是否异常，第5项
% #HANDWRITTEN: \checkmark
\fieldvalue{\handwritten{\checkmark}} \\
\hline
\end{tabularx}

% TODO: The lower portion of this table continues beyond the visible page area.
% LEXOID_PAGE_COMPLETED: 38/78

\newpage

\begin{center}
\textbf{水质监督检查记录表（22）}
\end{center}

\noindent
\textbf{受检单位：}
% #VALUE_ID: LEX-P0039-V0001
% #FIELD_VALUE: 受检单位
\fieldvalue{\handwritten{（字迹不清）}}

\hfill
\textbf{检查日期：}
% #VALUE_ID: LEX-P0039-V0002
% #FIELD_VALUE: 检查日期
% #HANDWRITTEN: 202?年?月?日
\fieldvalue{\handwritten{202?年?月?日}}

\hfill
\textbf{检查人：}
% #VALUE_ID: LEX-P0039-V0003
% #FIELD_VALUE: 检查人
% #HANDWRITTEN: （字迹不清）
\fieldvalue{\handwritten{（字迹不清）}}

\vspace{0.5em}

\begin{tabularx}{\textwidth}{|c|X|X|X|}
\hline
序号 & 检查项目 & 检查内容 & 检查结果 \\
\hline
6 &
饮用水卫生管理 &
相关卫生管理制度及落实情况 &
% #VALUE_ID: LEX-P0039-V0004
% #FIELD_VALUE: 第6项检查结果
% #HANDWRITTEN: 合格
\fieldvalue{\handwritten{合格}}
\\
\hline
7 &
供水设施 &
供水设施及防护情况 &
% #VALUE_ID: LEX-P0039-V0005
% #FIELD_VALUE: 第7项检查结果
% #HANDWRITTEN: 合格
\fieldvalue{\handwritten{合格}}
\\
\hline
8 &
水质检测 &
水质检测及记录情况 &
% #VALUE_ID: LEX-P0039-V0006
% #FIELD_VALUE: 第8项检查结果
% #HANDWRITTEN: 合格
\fieldvalue{\handwritten{合格}}
\\
\hline
\end{tabularx}

\vspace{0.5em}

\noindent
\textbf{检查情况：}

\begin{tabularx}{\textwidth}{|c|X|X|X|}
\hline
序号 & 检查项目 & 检查内容 & 检查结果 \\
\hline
1 &
水源卫生防护 &
水源周围卫生防护情况 &
% #VALUE_ID: LEX-P0039-V0007
% #FIELD_VALUE: 第1项检查结果
% #HANDWRITTEN: 合格
\fieldvalue{\handwritten{合格}}
\\
\hline
2 &
供水单位卫生管理 &
卫生管理制度、人员健康管理等情况 &
% #VALUE_ID: LEX-P0039-V0008
% #FIELD_VALUE: 第2项检查结果
% #HANDWRITTEN: 合格
\fieldvalue{\handwritten{合格}}
\\
\hline
3 &
制水工艺及设施 &
制水工艺、设施设备运行情况 &
% #VALUE_ID: LEX-P0039-V0009
% #FIELD_VALUE: 第3项检查结果
% #HANDWRITTEN: 合格
\fieldvalue{\handwritten{合格}}
\\
\hline
4 &
消毒管理 &
消毒设施及消毒记录情况 &
% #VALUE_ID: LEX-P0039-V0010
% #FIELD_VALUE: 第4项检查结果
% #HANDWRITTEN: 合格
\fieldvalue{\handwritten{合格}}
\\
\hline
5 &
水质检测 &
水质检测项目、频次及记录情况 &
% #VALUE_ID: LEX-P0039-V0011
% #FIELD_VALUE: 第5项检查结果
% #HANDWRITTEN: 合格
\fieldvalue{\handwritten{合格}}
\\
\hline
6 &
从业人员健康管理 &
从业人员健康证明及培训情况 &
% #VALUE_ID: LEX-P0039-V0012
% #FIELD_VALUE: 第6项检查结果
% #HANDWRITTEN: 合格
\fieldvalue{\handwritten{合格}}
\\
\hline
7 &
卫生许可证及档案 &
相关证照、档案及记录情况 &
% #VALUE_ID: LEX-P0039-V0013
% #FIELD_VALUE: 第7项检查结果
% #HANDWRITTEN: 合格
\fieldvalue{\handwritten{合格}}
\\
\hline
\end{tabularx}

\vspace{0.5em}

\noindent
\textbf{存在问题及整改要求：}

\begin{tabularx}{\textwidth}{|X|X|}
\hline
存在问题 &
% #VALUE_ID: LEX-P0039-V0014
% #FIELD_VALUE: 存在问题
% #HANDWRITTEN: （字迹不清）
\fieldvalue{\handwritten{（字迹不清）}}
\\
\hline
整改要求 &
% #VALUE_ID: LEX-P0039-V0015
% #FIELD_VALUE: 整改要求
% #HANDWRITTEN: （字迹不清）
\fieldvalue{\handwritten{（字迹不清）}}
\\
\hline
\end{tabularx}

\vspace{0.5em}

\noindent
\textbf{备注：} 本页下方表格及签章区域在本页可见；部分手写内容因原图倒置及分辨率限制无法准确辨认。

% TODO：页面下部存在印章、日期及签名等内容，部分文字过于模糊，需依据高清原图复核。
% LEXOID_PAGE_COMPLETED: 39/78

\newpage

\begin{center}
\textbf{運行記録関係書類}
\end{center}

% TODO: The page image is rotated and several Japanese labels are difficult to read.
\noindent
\begin{tabularx}{\textwidth}{|l|X|l|X|}
\hline
車両番号 & 
% #VALUE_ID: LEX-P0040-V0001
% #FIELD_VALUE: 車両番号
\fieldvalue{\handwritten{判読不能}} &
年月日 &
% #VALUE_ID: LEX-P0040-V0002
% #FIELD_VALUE: 年月日
\fieldvalue{\handwritten{判読不能}} \\
\hline
運転者 & 
% #VALUE_ID: LEX-P0040-V0003
% #FIELD_VALUE: 運転者
\fieldvalue{\handwritten{判読不能}} &
確認者 &
% #VALUE_ID: LEX-P0040-V0004
% #FIELD_VALUE: 確認者
\fieldvalue{\handwritten{判読不能}} \\
\hline
\end{tabularx}

\vspace{1em}

\begin{tabularx}{\textwidth}{|c|X|X|X|}
\hline
\multicolumn{4}{|c|}{運行状況記録} \\
\hline
番号 & 運行年月日・行先 & 出発時刻・距離 & 到着時刻・距離 \\
\hline
8 &
% #VALUE_ID: LEX-P0040-V0005
% #FIELD_VALUE: 8番・運行年月日または行先
\fieldvalue{\handwritten{判読不能}} &
% #VALUE_ID: LEX-P0040-V0006
% #FIELD_VALUE: 8番・出発時刻または距離
\fieldvalue{\handwritten{40}} &
% #VALUE_ID: LEX-P0040-V0007
% #FIELD_VALUE: 8番・到着時刻または距離
\fieldvalue{\handwritten{40}} \\
\hline
9 &
% #VALUE_ID: LEX-P0040-V0008
% #FIELD_VALUE: 9番・運行年月日または行先
\fieldvalue{\handwritten{判読不能}} &
% #VALUE_ID: LEX-P0040-V0009
% #FIELD_VALUE: 9番・出発時刻または距離
\fieldvalue{\handwritten{50}} &
% #VALUE_ID: LEX-P0040-V0010
% #FIELD_VALUE: 9番・到着時刻または距離
\fieldvalue{\handwritten{判読不能}} \\
\hline
10 &
% #VALUE_ID: LEX-P0040-V0011
% #FIELD_VALUE: 10番・運行年月日または行先
\fieldvalue{\handwritten{判読不能}} &
% #VALUE_ID: LEX-P0040-V0012
% #FIELD_VALUE: 10番・出発時刻または距離
\fieldvalue{\handwritten{25}} &
% #VALUE_ID: LEX-P0040-V0013
% #FIELD_VALUE: 10番・到着時刻または距離
\fieldvalue{\handwritten{4}} \\
\hline
11 &
% #VALUE_ID: LEX-P0040-V0014
% #FIELD_VALUE: 11番・運行年月日または行先
\fieldvalue{\handwritten{判読不能}} &
% #VALUE_ID: LEX-P0040-V0015
% #FIELD_VALUE: 11番・出発時刻または距離
\fieldvalue{\handwritten{20}} &
% #VALUE_ID: LEX-P0040-V0016
% #FIELD_VALUE: 11番・到着時刻または距離
\fieldvalue{\handwritten{20}} \\
\hline
12 &
% #VALUE_ID: LEX-P0040-V0017
% #FIELD_VALUE: 12番・運行年月日または行先
\fieldvalue{\handwritten{判読不能}} &
% #VALUE_ID: LEX-P0040-V0018
% #FIELD_VALUE: 12番・出発時刻または距離
\fieldvalue{\handwritten{20}} &
% #VALUE_ID: LEX-P0040-V0019
% #FIELD_VALUE: 12番・到着時刻または距離
\fieldvalue{\handwritten{20}} \\
\hline
13 &
% #VALUE_ID: LEX-P0040-V0020
% #FIELD_VALUE: 13番・運行年月日または行先
\fieldvalue{\handwritten{判読不能}} &
% #VALUE_ID: LEX-P0040-V0021
% #FIELD_VALUE: 13番・出発時刻または距離
\fieldvalue{\handwritten{20}} &
% #VALUE_ID: LEX-P0040-V0022
% #FIELD_VALUE: 13番・到着時刻または距離
\fieldvalue{\handwritten{20}} \\
\hline
14 &
% #VALUE_ID: LEX-P0040-V0023
% #FIELD_VALUE: 14番・運行年月日または行先
\fieldvalue{\handwritten{判読不能}} &
% #VALUE_ID: LEX-P0040-V0024
% #FIELD_VALUE: 14番・出発時刻または距離
\fieldvalue{\handwritten{20}} &
% #VALUE_ID: LEX-P0040-V0025
% #FIELD_VALUE: 14番・到着時刻または距離
\fieldvalue{\handwritten{20}} \\
\hline
15 &
% #VALUE_ID: LEX-P0040-V0026
% #FIELD_VALUE: 15番・運行年月日または行先
\fieldvalue{\handwritten{判読不能}} &
% #VALUE_ID: LEX-P0040-V0027
% #FIELD_VALUE: 15番・出発時刻または距離
\fieldvalue{\handwritten{05}} &
% #VALUE_ID: LEX-P0040-V0028
% #FIELD_VALUE: 15番・到着時刻または距離
\fieldvalue{\handwritten{05}} \\
\hline
\end{tabularx}

\vspace{1em}

\noindent
\begin{tabularx}{\textwidth}{|l|X|l|X|}
\hline
合計 &
% #VALUE_ID: LEX-P0040-V0029
% #FIELD_VALUE: 合計・左欄
\fieldvalue{\handwritten{判読不能}} &
合計 &
% #VALUE_ID: LEX-P0040-V0030
% #FIELD_VALUE: 合計・右欄
\fieldvalue{\handwritten{判読不能}} \\
\hline
点検結果 &
% #VALUE_ID: LEX-P0040-V0031
% #FIELD_VALUE: 点検結果
\fieldvalue{\handwritten{判読不能}} &
備考 &
% #VALUE_ID: LEX-P0040-V0032
% #FIELD_VALUE: 備考
\fieldvalue{\handwritten{判読不能}} \\
\hline
\end{tabularx}

\vspace{1em}

\noindent\textbf{備考・記入欄}

\begin{itemize}
\item
% #VALUE_ID: LEX-P0040-V0033
% #FIELD_VALUE: 備考欄第1行
\fieldvalue{\handwritten{判読不能}}
\item
% #VALUE_ID: LEX-P0040-V0034
% #FIELD_VALUE: 備考欄第2行
\fieldvalue{\handwritten{判読不能}}
\item
% #VALUE_ID: LEX-P0040-V0035
% #FIELD_VALUE: 備考欄第3行
\fieldvalue{\handwritten{判読不能}}
\item
% #VALUE_ID: LEX-P0040-V0036
% #FIELD_VALUE: 備考欄第4行
\fieldvalue{\handwritten{判読不能}}
\end{itemize}

% TODO: Additional handwritten entries and printed labels are visible in the rotated upper portion of the page, but their field boundaries and text cannot be determined reliably from the scan.
% LEXOID_PAGE_COMPLETED: 40/78

\newpage

\begin{center}
\textbf{附件2}\\[1ex]
\textbf{食品经营许可证申请表（新设）}
\end{center}

\noindent
\begin{tabularx}{\textwidth}{|l|X|}
\hline
申请人名称 &
% #VALUE_ID: LEX-P0041-V0001
% #FIELD_VALUE: 申请人名称
% #TODO #HANDWRITTEN: 字迹不清；蓝色手写内容无法准确辨识
\fieldvalue{\handwritten{字迹不清}}\\
\hline
法定代表人（负责人） &
% #VALUE_ID: LEX-P0041-V0002
% #FIELD_VALUE: 法定代表人（负责人）
% #TODO #HANDWRITTEN: 字迹不清；红色手写内容无法准确辨识
\fieldvalue{\handwritten{字迹不清}}\\
\hline
住所 &
% #VALUE_ID: LEX-P0041-V0003
% #FIELD_VALUE: 住所
% #TODO #HANDWRITTEN: 字迹不清；蓝色手写内容无法准确辨识
\fieldvalue{\handwritten{字迹不清}}\\
\hline
\end{tabularx}

\vspace{1ex}

\noindent
\textbf{食品经营许可证申请表（新设）}

\noindent
1.\ 食品经营许可证编号：\rule{0.55\textwidth}{0.4pt}

\noindent
2.\ 申请日期：
% #VALUE_ID: LEX-P0041-V0004
% #FIELD_VALUE: 申请日期
% #TODO #HANDWRITTEN: 约为“2015年10月22日”；字迹不清
\fieldvalue{\handwritten{2015年10月22日}}

\vspace{1ex}

\noindent
\begin{tabularx}{\textwidth}{|l|X|}
\hline
主体业态 &
% #VALUE_ID: LEX-P0041-V0005
% #FIELD_VALUE: 主体业态
% #TODO #HANDWRITTEN: 字迹不清；蓝色手写内容无法准确辨识
\fieldvalue{\handwritten{字迹不清}}\\
\hline
经营场所 &
% #VALUE_ID: LEX-P0041-V0006
% #FIELD_VALUE: 经营场所
% #TODO #HANDWRITTEN: 字迹不清；蓝色手写内容无法准确辨识
\fieldvalue{\handwritten{字迹不清}}\\
\hline
经营项目 &
% #VALUE_ID: LEX-P0041-V0007
% #FIELD_VALUE: 经营项目
% #TODO #HANDWRITTEN: 字迹不清；蓝色手写内容无法准确辨识
\fieldvalue{\handwritten{字迹不清}}\\
\hline
\end{tabularx}

\vspace{1ex}

\noindent
\textbf{一、主体业态}

\begin{tabularx}{\textwidth}{|l|X|}
\hline
主体业态 & 选择（在相应项目后的方框内打勾）\\
\hline
食品销售经营者 & $\square$\\
\hline
餐饮服务经营者 & $\square$\\
\hline
单位食堂 & $\square$\\
\hline
\end{tabularx}

\vspace{1ex}

\noindent
\textbf{二、经营项目}

\begin{tabularx}{\textwidth}{|l|X|}
\hline
经营项目 & 选择\\
\hline
预包装食品销售 & $\square$\\
\hline
散装食品销售 & $\square$\\
\hline
特殊食品销售 & $\square$\\
\hline
热食类食品制售 & $\square$\\
\hline
冷食类食品制售 & $\square$\\
\hline
生食类食品制售 & $\square$\\
\hline
糕点类食品制售 & $\square$\\
\hline
自制饮品制售 & $\square$\\
\hline
其他类食品制售 & $\square$\\
\hline
\end{tabularx}

% TODO: 本页表格及手写内容部分字迹模糊，以上无法辨识的手写值需结合高清原图复核。
% LEXOID_PAGE_COMPLETED: 41/78

\newpage

\begin{center}
\textbf{医疗卫生机构传染病防治分类监督综合评价表}\\
% TODO：页面顶部机构名称及文号部分文字较小，部分内容无法可靠辨识。
\end{center}

\noindent
\begin{tabularx}{\textwidth}{|l|X|l|l|}
\hline
机构名称 & & 检查日期 & 
% #VALUE_ID: LEX-P0042-V0001
% #FIELD_VALUE: 检查日期
% #HANDWRITTEN: 2024.08.18
\fieldvalue{\handwritten{2024.08.18}}\\
\hline
检查人员 & 
% #VALUE_ID: LEX-P0042-V0002
% #FIELD_VALUE: 检查人员
% TODO #HANDWRITTEN: 张\ldots；姓名部分辨识不清
\fieldvalue{\handwritten{张\ldots}} & & \\
\hline
\end{tabularx}

\vspace{0.5em}

\noindent
\begin{tabularx}{\textwidth}{|c|X|c|}
\hline
\textbf{项目} & \textbf{内容} & \textbf{评价结果}\\
\hline
一、传染病防治组织管理
&
\begin{minipage}[t]{\linewidth}
\small
1.\ 建立传染病防治管理组织，明确职责；\\
2.\ 制定传染病防治相关制度；\\
3.\ 配备专（兼）职传染病防治工作人员；\\
4.\ 开展传染病防治法律法规及相关知识培训；\\
5.\ 将传染病防治工作纳入年度工作计划；\\
6.\ 定期开展传染病防治工作检查；\\
7.\ 及时整改监督检查中发现的问题。
\end{minipage}
&
\begin{minipage}[t]{\linewidth}
\small
是\ $\square$\quad 否\ $\square$\\[0.35em]
% #VALUE_ID: LEX-P0042-V0003
% #FIELD_VALUE: 第1项评价结果
% #HANDWRITTEN: 是
\fieldvalue{\handwritten{是}}\\[0.35em]
是\ $\square$\quad 否\ $\square$\\[0.35em]
% #VALUE_ID: LEX-P0042-V0004
% #FIELD_VALUE: 第2项评价结果
% #HANDWRITTEN: 是
\fieldvalue{\handwritten{是}}\\[0.35em]
是\ $\square$\quad 否\ $\square$\\[0.35em]
% #VALUE_ID: LEX-P0042-V0005
% #FIELD_VALUE: 第3项评价结果
% #HANDWRITTEN: 是
\fieldvalue{\handwritten{是}}\\[0.35em]
是\ $\square$\quad 否\ $\square$\\[0.35em]
% #VALUE_ID: LEX-P0042-V0006
% #FIELD_VALUE: 第4项评价结果
% #HANDWRITTEN: 是
\fieldvalue{\handwritten{是}}\\[0.35em]
是\ $\square$\quad 否\ $\square$\\[0.35em]
% #VALUE_ID: LEX-P0042-V0007
% #FIELD_VALUE: 第5项评价结果
% #HANDWRITTEN: 是
\fieldvalue{\handwritten{是}}\\[0.35em]
是\ $\square$\quad 否\ $\square$\\[0.35em]
% #VALUE_ID: LEX-P0042-V0008
% #FIELD_VALUE: 第6项评价结果
% #HANDWRITTEN: 是
\fieldvalue{\handwritten{是}}\\[0.35em]
是\ $\square$\quad 否\ $\square$
% #VALUE_ID: LEX-P0042-V0009
% #FIELD_VALUE: 第7项评价结果
% #HANDWRITTEN: 是
\fieldvalue{\handwritten{是}}
\end{minipage}
\\
\hline
整改情况
&
\begin{minipage}[t]{\linewidth}
\small
对监督检查中发现的问题进行整改。
\end{minipage}
&
% #VALUE_ID: LEX-P0042-V0010
% #FIELD_VALUE: 整改情况
% #TODO #HANDWRITTEN: 已整改\ldots；蓝色手写内容辨识不清
\fieldvalue{\handwritten{已整改\ldots}}
\\
\hline
\end{tabularx}

\vspace{0.5em}

\noindent
\begin{tabularx}{\textwidth}{|c|X|c|c|}
\hline
\textbf{序号} & \textbf{检查项目} & \textbf{评价分值} & \textbf{得分}\\
\hline
1 & 传染病防治组织管理 & 15 & 
% #VALUE_ID: LEX-P0042-V0011
% #FIELD_VALUE: 第1项得分
% #TODO #HANDWRITTEN: 12.75；蓝色手写数字部分辨识不清
\fieldvalue{\handwritten{12.75}}\\
\hline
2 & 传染病防治制度建设 & 15 & 
% #VALUE_ID: LEX-P0042-V0012
% #FIELD_VALUE: 第2项得分
% #TODO #HANDWRITTEN: 12.06；蓝色手写数字部分辨识不清
\fieldvalue{\handwritten{12.06}}\\
\hline
3 & 传染病防治工作落实 & 15 & 
% #VALUE_ID: LEX-P0042-V0013
% #FIELD_VALUE: 第3项得分
% #TODO #HANDWRITTEN: 12.75；蓝色手写数字部分辨识不清
\fieldvalue{\handwritten{12.75}}\\
\hline
\end{tabularx}

\noindent
\textbf{附件：}监督检查记录表。

% TODO：本页表格及页眉内容在扫描件中呈倒置方向；以上按正常阅读方向重排。
% LEXOID_PAGE_COMPLETED: 42/78

\newpage

\section*{建筑施工特种作业人员相关记录表}

% TODO: The page is a scanned, rotated Chinese form; several fixed labels are partially illegible.

\noindent
填表人：
% #VALUE_ID: LEX-P0043-V0001
% #FIELD_VALUE: 填表人
% #TODO #HANDWRITTEN: [姓名]; handwriting is illegible
\fieldvalue{\handwritten{[姓名]}}
\hfill
填表日期：
% #VALUE_ID: LEX-P0043-V0002
% #FIELD_VALUE: 填表日期
% #TODO #HANDWRITTEN: 2025年06月18日; date partially obscured
\fieldvalue{\handwritten{2025年06月18日}}

\vspace{1em}

\begin{table}[h]
\centering
\small
\begin{tabularx}{\textwidth}{|c|X|X|X|c|}
\hline
序号 & 设备名称 & 规格型号/编号 & 检查日期或相关记录 & 备注 \\
\hline
4 &
塔吊 &
% #VALUE_ID: LEX-P0043-V0003
% #FIELD_VALUE: 规格型号/编号，第4项
% #TODO #HANDWRITTEN: 11005; handwritten number is partially illegible
\fieldvalue{\handwritten{11005}} &
% #VALUE_ID: LEX-P0043-V0004
% #FIELD_VALUE: 检查结果，第4项
% #HANDWRITTEN: 已检查
\fieldvalue{\handwritten{\checkmark}} &
\\
\hline
5 &
塔吊 &
% #VALUE_ID: LEX-P0043-V0005
% #FIELD_VALUE: 规格型号/编号，第5项
% #HANDWRITTEN: /
\fieldvalue{\handwritten{/}} &
% #VALUE_ID: LEX-P0043-V0006
% #FIELD_VALUE: 检查结果，第5项
% #HANDWRITTEN: 已检查
\fieldvalue{\handwritten{\checkmark}} &
\\
\hline
6 &
塔吊 &
% #VALUE_ID: LEX-P0043-V0007
% #FIELD_VALUE: 规格型号/编号，第6项
% #TODO #HANDWRITTEN: 123?; final digits are unclear
\fieldvalue{\handwritten{123?}} &
% #VALUE_ID: LEX-P0043-V0008
% #FIELD_VALUE: 检查结果，第6项
% #HANDWRITTEN: 已检查
\fieldvalue{\handwritten{\checkmark}} &
\\
\hline
7 &
塔吊 &
% #VALUE_ID: LEX-P0043-V0009
% #FIELD_VALUE: 规格型号/编号，第7项
% #TODO #HANDWRITTEN: 123?; handwriting is unclear
\fieldvalue{\handwritten{123?}} &
% #VALUE_ID: LEX-P0043-V0010
% #FIELD_VALUE: 检查结果，第7项
% #HANDWRITTEN: 已检查
\fieldvalue{\handwritten{\checkmark}} &
\\
\hline
8 &
塔吊 &
% #VALUE_ID: LEX-P0043-V0011
% #FIELD_VALUE: 规格型号/编号，第8项
% #TODO #HANDWRITTEN: [编号]; handwriting is illegible
\fieldvalue{\handwritten{[编号]}} &
% #VALUE_ID: LEX-P0043-V0012
% #FIELD_VALUE: 检查结果，第8项
% #HANDWRITTEN: 已检查
\fieldvalue{\handwritten{\checkmark}} &
\\
\hline
9 &
塔吊 &
% #VALUE_ID: LEX-P0043-V0013
% #FIELD_VALUE: 规格型号/编号，第9项
% #TODO #HANDWRITTEN: [编号]; handwriting is illegible
\fieldvalue{\handwritten{[编号]}} &
% #VALUE_ID: LEX-P0043-V0014
% #FIELD_VALUE: 检查结果，第9项
% #HANDWRITTEN: 已检查
\fieldvalue{\handwritten{\checkmark}} &
\\
\hline
\end{tabularx}
\end{table}

\vspace{1em}

\begin{table}[h]
\centering
\small
\begin{tabularx}{\textwidth}{|c|X|X|X|c|}
\hline
序号 & 设备名称 & 规格型号/编号 & 相关记录 & 检查结果 \\
\hline
1 &
设备 &
% #VALUE_ID: LEX-P0043-V0015
% #FIELD_VALUE: 规格型号/编号，第1项
% #TODO #HANDWRITTEN: [编号]; handwriting is illegible
\fieldvalue{\handwritten{[编号]}} &
% #VALUE_ID: LEX-P0043-V0016
% #FIELD_VALUE: 检查结果，第1项
% #HANDWRITTEN: 已检查
\fieldvalue{\handwritten{\checkmark}} &
\\
\hline
2 &
设备 &
% #VALUE_ID: LEX-P0043-V0017
% #FIELD_VALUE: 规格型号/编号，第2项
% #TODO #HANDWRITTEN: [编号]; handwriting is illegible
\fieldvalue{\handwritten{[编号]}} &
% #VALUE_ID: LEX-P0043-V0018
% #FIELD_VALUE: 检查结果，第2项
% #HANDWRITTEN: 已检查
\fieldvalue{\handwritten{\checkmark}} &
\\
\hline
3 &
设备 &
% #VALUE_ID: LEX-P0043-V0019
% #FIELD_VALUE: 规格型号/编号，第3项
% #TODO #HANDWRITTEN: [编号]; handwriting is illegible
\fieldvalue{\handwritten{[编号]}} &
% #VALUE_ID: LEX-P0043-V0020
% #FIELD_VALUE: 检查结果，第3项
% #HANDWRITTEN: 已检查
\fieldvalue{\handwritten{\checkmark}} &
\\
\hline
4 &
设备 &
% #VALUE_ID: LEX-P0043-V0021
% #FIELD_VALUE: 规格型号/编号，第4项
% #TODO #HANDWRITTEN: [编号]; handwriting is illegible
\fieldvalue{\handwritten{[编号]}} &
% #VALUE_ID: LEX-P0043-V0022
% #FIELD_VALUE: 检查结果，第4项
% #HANDWRITTEN: 已检查
\fieldvalue{\handwritten{\checkmark}} &
\\
\hline
5 &
设备 &
% #VALUE_ID: LEX-P0043-V0023
% #FIELD_VALUE: 规格型号/编号，第5项
% #TODO #HANDWRITTEN: [编号]; handwriting is illegible
\fieldvalue{\handwritten{[编号]}} &
% #VALUE_ID: LEX-P0043-V0024
% #FIELD_VALUE: 检查结果，第5项
% #HANDWRITTEN: 已检查
\fieldvalue{\handwritten{\checkmark}} &
\\
\hline
6 &
设备 &
% #VALUE_ID: LEX-P0043-V0025
% #FIELD_VALUE: 规格型号/编号，第6项
% #TODO #HANDWRITTEN: [编号]; handwriting is illegible
\fieldvalue{\handwritten{[编号]}} &
% #VALUE_ID: LEX-P0043-V0026
% #FIELD_VALUE: 检查结果，第6项
% #HANDWRITTEN: 已检查
\fieldvalue{\handwritten{\checkmark}} &
\\
\hline
7 &
设备 &
% #VALUE_ID: LEX-P0043-V0027
% #FIELD_VALUE: 规格型号/编号，第7项
% #TODO #HANDWRITTEN: [编号]; handwriting is illegible
\fieldvalue{\handwritten{[编号]}} &
% #VALUE_ID: LEX-P0043-V0028
% #FIELD_VALUE: 检查结果，第7项
% #HANDWRITTEN: 已检查
\fieldvalue{\handwritten{\checkmark}} &
\\
\hline
8 &
设备 &
% #VALUE_ID: LEX-P0043-V0029
% #FIELD_VALUE: 规格型号/编号，第8项
% #TODO #HANDWRITTEN: [编号]; handwriting is illegible
\fieldvalue{\handwritten{[编号]}} &
% #VALUE_ID: LEX-P0043-V0030
% #FIELD_VALUE: 检查结果，第8项
% #HANDWRITTEN: 已检查
\fieldvalue{\handwritten{\checkmark}} &
\\
\hline
9 &
设备 &
% #VALUE_ID: LEX-P0043-V0031
% #FIELD_VALUE: 规格型号/编号，第9项
% #TODO #HANDWRITTEN: [编号]; handwriting is illegible
\fieldvalue{\handwritten{[编号]}} &
% #VALUE_ID: LEX-P0043-V0032
% #FIELD_VALUE: 检查结果，第9项
% #HANDWRITTEN: 已检查
\fieldvalue{\handwritten{\checkmark}} &
\\
\hline
10 &
设备 &
% #VALUE_ID: LEX-P0043-V0033
% #FIELD_VALUE: 规格型号/编号，第10项
% #TODO #HANDWRITTEN: [编号]; handwriting is illegible
\fieldvalue{\handwritten{[编号]}} &
% #VALUE_ID: LEX-P0043-V0034
% #FIELD_VALUE: 检查结果，第10项
% #HANDWRITTEN: 已检查
\fieldvalue{\handwritten{\checkmark}} &
\\
\hline
11 &
设备 &
% #VALUE_ID: LEX-P0043-V0035
% #FIELD_VALUE: 规格型号/编号，第11项
% #TODO #HANDWRITTEN: [编号]; handwriting is illegible
\fieldvalue{\handwritten{[编号]}} &
% #VALUE_ID: LEX-P0043-V0036
% #FIELD_VALUE: 检查结果，第11项
% #HANDWRITTEN: 已检查
\fieldvalue{\handwritten{\checkmark}} &
\\
\hline
12 &
设备 &
% #VALUE_ID: LEX-P0043-V0037
% #FIELD_VALUE: 规格型号/编号，第12项
% #TODO #HANDWRITTEN: [编号]; handwriting is illegible
\fieldvalue{\handwritten{[编号]}} &
% #VALUE_ID: LEX-P0043-V0038
% #FIELD_VALUE: 检查结果，第12项
% #HANDWRITTEN: 已检查
\fieldvalue{\handwritten{\checkmark}} &
\\
\hline
\end{tabularx}
\end{table}

% TODO: The lower portion of the form continues beyond the visible page area.
% LEXOID_PAGE_COMPLETED: 43/78

\newpage

\section*{附件2-2：物资领用登记表}

\begin{center}
\textbf{物资领用登记表}
\end{center}

\noindent
\textbf{3.1 物资领用登记}

\begin{tabularx}{\textwidth}{c|X|X|X|c|c|X}
\hline
序号 & 物资名称 & 型号 & 规格 & 数量 & 单位 & 备注\\
\hline
13 &
% #VALUE_ID: LEX-P0044-V0001
% #FIELD_VALUE: 物资名称（第13项）
% #TODO #HANDWRITTEN: 难辨；蓝色手写字迹不清
\fieldvalue{\handwritten{难辨}} &
% #VALUE_ID: LEX-P0044-V0002
% #FIELD_VALUE: 型号（第13项）
% #TODO #HANDWRITTEN: 难辨；蓝色手写字迹不清
\fieldvalue{\handwritten{难辨}} &
% #VALUE_ID: LEX-P0044-V0003
% #FIELD_VALUE: 规格（第13项）
% #TODO #HANDWRITTEN: 难辨；蓝色手写字迹不清
\fieldvalue{\handwritten{难辨}} &
% #VALUE_ID: LEX-P0044-V0004
% #FIELD_VALUE: 数量（第13项）
% #HANDWRITTEN: 1
\fieldvalue{\handwritten{1}} &
个 &
% #VALUE_ID: LEX-P0044-V0005
% #FIELD_VALUE: 备注（第13项）
% #TODO #HANDWRITTEN: 难辨；蓝色手写字迹不清
\fieldvalue{\handwritten{难辨}}\\
\hline
14 &
% #VALUE_ID: LEX-P0044-V0006
% #FIELD_VALUE: 物资名称（第14项）
% #TODO #HANDWRITTEN: 难辨；蓝色手写字迹不清
\fieldvalue{\handwritten{难辨}} &
% #VALUE_ID: LEX-P0044-V0007
% #FIELD_VALUE: 型号（第14项）
% #TODO #HANDWRITTEN: 难辨；蓝色手写字迹不清
\fieldvalue{\handwritten{难辨}} &
% #VALUE_ID: LEX-P0044-V0008
% #FIELD_VALUE: 规格（第14项）
% #TODO #HANDWRITTEN: 难辨；蓝色手写字迹不清
\fieldvalue{\handwritten{难辨}} &
% #VALUE_ID: LEX-P0044-V0009
% #FIELD_VALUE: 数量（第14项）
% #HANDWRITTEN: 1
\fieldvalue{\handwritten{1}} &
个 &
% #VALUE_ID: LEX-P0044-V0010
% #FIELD_VALUE: 备注（第14项）
% #TODO #HANDWRITTEN: 难辨；蓝色手写字迹不清
\fieldvalue{\handwritten{难辨}}\\
\hline
15 &
% #VALUE_ID: LEX-P0044-V0011
% #FIELD_VALUE: 物资名称（第15项）
% #TODO #HANDWRITTEN: 难辨；蓝色手写字迹不清
\fieldvalue{\handwritten{难辨}} &
% #VALUE_ID: LEX-P0044-V0012
% #FIELD_VALUE: 型号（第15项）
% #TODO #HANDWRITTEN: 难辨；蓝色手写字迹不清
\fieldvalue{\handwritten{难辨}} &
% #VALUE_ID: LEX-P0044-V0013
% #FIELD_VALUE: 规格（第15项）
% #TODO #HANDWRITTEN: 难辨；蓝色手写字迹不清
\fieldvalue{\handwritten{难辨}} &
% #VALUE_ID: LEX-P0044-V0014
% #FIELD_VALUE: 数量（第15项）
% #HANDWRITTEN: 1
\fieldvalue{\handwritten{1}} &
个 &
% #VALUE_ID: LEX-P0044-V0015
% #FIELD_VALUE: 备注（第15项）
% #TODO #HANDWRITTEN: 难辨；蓝色手写字迹不清
\fieldvalue{\handwritten{难辨}}\\
\hline
16 &
% #VALUE_ID: LEX-P0044-V0016
% #FIELD_VALUE: 物资名称（第16项）
% #TODO #HANDWRITTEN: 难辨；蓝色手写字迹不清
\fieldvalue{\handwritten{难辨}} &
% #VALUE_ID: LEX-P0044-V0017
% #FIELD_VALUE: 型号（第16项）
% #TODO #HANDWRITTEN: 难辨；蓝色手写字迹不清
\fieldvalue{\handwritten{难辨}} &
% #VALUE_ID: LEX-P0044-V0018
% #FIELD_VALUE: 规格（第16项）
% #TODO #HANDWRITTEN: 难辨；蓝色手写字迹不清
\fieldvalue{\handwritten{难辨}} &
% #VALUE_ID: LEX-P0044-V0019
% #FIELD_VALUE: 数量（第16项）
% #HANDWRITTEN: 1
\fieldvalue{\handwritten{1}} &
个 &
% #VALUE_ID: LEX-P0044-V0020
% #FIELD_VALUE: 备注（第16项）
% #TODO #HANDWRITTEN: 难辨；蓝色手写字迹不清
\fieldvalue{\handwritten{难辨}}\\
\hline
17 &
% #VALUE_ID: LEX-P0044-V0021
% #FIELD_VALUE: 物资名称（第17项）
% #TODO #HANDWRITTEN: 难辨；蓝色手写字迹不清
\fieldvalue{\handwritten{难辨}} &
% #VALUE_ID: LEX-P0044-V0022
% #FIELD_VALUE: 型号（第17项）
% #TODO #HANDWRITTEN: 难辨；蓝色手写字迹不清
\fieldvalue{\handwritten{难辨}} &
% #VALUE_ID: LEX-P0044-V0023
% #FIELD_VALUE: 规格（第17项）
% #TODO #HANDWRITTEN: 难辨；蓝色手写字迹不清
\fieldvalue{\handwritten{难辨}} &
% #VALUE_ID: LEX-P0044-V0024
% #FIELD_VALUE: 数量（第17项）
% #HANDWRITTEN: 1
\fieldvalue{\handwritten{1}} &
个 &
% #VALUE_ID: LEX-P0044-V0025
% #FIELD_VALUE: 备注（第17项）
% #TODO #HANDWRITTEN: 难辨；蓝色手写字迹不清
\fieldvalue{\handwritten{难辨}}\\
\hline
18 &
% #VALUE_ID: LEX-P0044-V0026
% #FIELD_VALUE: 物资名称（第18项）
% #TODO #HANDWRITTEN: 难辨；蓝色手写字迹不清
\fieldvalue{\handwritten{难辨}} &
% #VALUE_ID: LEX-P0044-V0027
% #FIELD_VALUE: 型号（第18项）
% #TODO #HANDWRITTEN: 难辨；蓝色手写字迹不清
\fieldvalue{\handwritten{难辨}} &
% #VALUE_ID: LEX-P0044-V0028
% #FIELD_VALUE: 规格（第18项）
% #TODO #HANDWRITTEN: 难辨；蓝色手写字迹不清
\fieldvalue{\handwritten{难辨}} &
% #VALUE_ID: LEX-P0044-V0029
% #FIELD_VALUE: 数量（第18项）
% #HANDWRITTEN: 1
\fieldvalue{\handwritten{1}} &
个 &
% #VALUE_ID: LEX-P0044-V0030
% #FIELD_VALUE: 备注（第18项）
% #TODO #HANDWRITTEN: 难辨；蓝色手写字迹不清
\fieldvalue{\handwritten{难辨}}\\
\hline
\end{tabularx}

% #VALUE_ID: LEX-P0044-V0031
% #FIELD_VALUE: 领用日期（上方表格）
% #HANDWRITTEN: 2025.6.18
\noindent
领用日期：\fieldvalue{\handwritten{2025.6.18}}

\bigskip

\noindent
\textbf{3.13 物资领用登记}

\begin{tabularx}{\textwidth}{c|X|X|X|c|c|X}
\hline
序号 & 物资名称 & 型号 & 规格 & 数量 & 单位 & 备注\\
\hline
1 &
% #VALUE_ID: LEX-P0044-V0032
% #FIELD_VALUE: 物资名称（第1项）
% #TODO #HANDWRITTEN: 难辨；蓝色手写字迹不清
\fieldvalue{\handwritten{难辨}} &
% #VALUE_ID: LEX-P0044-V0033
% #FIELD_VALUE: 型号（第1项）
% #TODO #HANDWRITTEN: 难辨；蓝色手写字迹不清
\fieldvalue{\handwritten{难辨}} &
% #VALUE_ID: LEX-P0044-V0034
% #FIELD_VALUE: 规格（第1项）
% #TODO #HANDWRITTEN: 难辨；蓝色手写字迹不清
\fieldvalue{\handwritten{难辨}} &
% #VALUE_ID: LEX-P0044-V0035
% #FIELD_VALUE: 数量（第1项）
% #HANDWRITTEN: 1
\fieldvalue{\handwritten{1}} &
个 &
% #VALUE_ID: LEX-P0044-V0036
% #FIELD_VALUE: 备注（第1项）
% #TODO #HANDWRITTEN: 难辨；蓝色手写字迹不清
\fieldvalue{\handwritten{难辨}}\\
\hline
2 &
% #VALUE_ID: LEX-P0044-V0037
% #FIELD_VALUE: 物资名称（第2项）
% #TODO #HANDWRITTEN: 难辨；蓝色手写字迹不清
\fieldvalue{\handwritten{难辨}} &
% #VALUE_ID: LEX-P0044-V0038
% #FIELD_VALUE: 型号（第2项）
% #TODO #HANDWRITTEN: 难辨；蓝色手写字迹不清
\fieldvalue{\handwritten{难辨}} &
% #VALUE_ID: LEX-P0044-V0039
% #FIELD_VALUE: 规格（第2项）
% #TODO #HANDWRITTEN: 难辨；蓝色手写字迹不清
\fieldvalue{\handwritten{难辨}} &
% #VALUE_ID: LEX-P0044-V0040
% #FIELD_VALUE: 数量（第2项）
% #HANDWRITTEN: 1
\fieldvalue{\handwritten{1}} &
个 &
% #VALUE_ID: LEX-P0044-V0041
% #FIELD_VALUE: 备注（第2项）
% #TODO #HANDWRITTEN: 难辨；蓝色手写字迹不清
\fieldvalue{\handwritten{难辨}}\\
\hline
\end{tabularx}

% #VALUE_ID: LEX-P0044-V0042
% #FIELD_VALUE: 领用日期（下方表格）
% #HANDWRITTEN: 2025.6.18
\noindent
领用日期：\fieldvalue{\handwritten{2025.6.18}}

\bigskip

\noindent
\textbf{3.2 其他物资}

\begin{tabularx}{\textwidth}{X|X}
\hline
项目 & 领用情况\\
\hline
% #VALUE_ID: LEX-P0044-V0043
% #FIELD_VALUE: 项目
% #TODO #HANDWRITTEN: 难辨；蓝色手写字迹不清
\fieldvalue{\handwritten{难辨}} &
% #VALUE_ID: LEX-P0044-V0044
% #FIELD_VALUE: 领用情况
% #TODO #HANDWRITTEN: 难辨；蓝色手写字迹不清
\fieldvalue{\handwritten{难辨}}\\
\hline
\end{tabularx}

% #VALUE_ID: LEX-P0044-V0045
% #FIELD_VALUE: 备注
% #TODO #HANDWRITTEN: 难辨；蓝色手写字迹不清
\noindent
备注：\fieldvalue{\handwritten{难辨}}
% LEXOID_PAGE_COMPLETED: 44/78

\newpage

% TODO: The page is reproduced from a scanned form whose source image is upside down and several Korean characters are unclear.

\begin{center}
\textbf{출입국 관련 신청서}
\end{center}

\noindent
\begin{tabularx}{\textwidth}{|l|X|l|X|}
\hline
구분 & 
% #VALUE_ID: LEX-P0045-V0001
% #FIELD_VALUE: 구분
\fieldvalue{\handwritten{불명확}} &
성명 &
% #VALUE_ID: LEX-P0045-V0002
% #FIELD_VALUE: 성명
% #HANDWRITTEN: 불명확
\fieldvalue{\handwritten{불명확}} \\
\hline
생년월일 &
% #VALUE_ID: LEX-P0045-V0003
% #FIELD_VALUE: 생년월일
% #HANDWRITTEN: 불명확한 날짜
\fieldvalue{\handwritten{불명확한 날짜}} &
문서번호 &
% #VALUE_ID: LEX-P0045-V0004
% #FIELD_VALUE: 문서번호
% #HANDWRITTEN: 불명확한 번호
\fieldvalue{\handwritten{불명확한 번호}} \\
\hline
\end{tabularx}

\vspace{0.5em}

\noindent
\begin{tabularx}{\textwidth}{|p{0.36\textwidth}|X|}
\hline
\textbf{신청 내용} &
\textbf{처리 사항} \\
\hline
\begin{minipage}[t]{\linewidth}
① 신청 구분: \(\square\) 해당  \(\square\) 미해당

② 신청 사유: \(\square\) 해당  \(\square\) 미해당

③ 체류 관련 사항: \(\square\) 해당  \(\square\) 미해당
\end{minipage}
&
\begin{minipage}[t]{\linewidth}
① 신청인의 인적사항 및 신청 내용을 확인합니다.

② 관련 서류를 검토하고 필요한 경우 보완을 요청합니다.

③ 심사 결과에 따라 허가 또는 불허 여부를 결정합니다.

④ 신청 내용에 변경이 있는 경우 즉시 신고하여야 합니다.
\end{minipage}
\\
\hline
\end{tabularx}

\vspace{0.5em}

\noindent
\begin{tabularx}{\textwidth}{|c|X|X|}
\hline
순번 & 사항 & 기재 내용 \\
\hline
1 &
신청 사항 &
% #VALUE_ID: LEX-P0045-V0005
% #FIELD_VALUE: 신청 사항 1
% #HANDWRITTEN: 불명확한 기재
\fieldvalue{\handwritten{불명확한 기재}} \\
\hline
2 &
신청 사항 &
% #VALUE_ID: LEX-P0045-V0006
% #FIELD_VALUE: 신청 사항 2
% #HANDWRITTEN: 불명확한 기재
\fieldvalue{\handwritten{불명확한 기재}} \\
\hline
3 &
신청 사항 &
% #VALUE_ID: LEX-P0045-V0007
% #FIELD_VALUE: 신청 사항 3
% #HANDWRITTEN: 불명확한 기재
\fieldvalue{\handwritten{불명확한 기재}} \\
\hline
4 &
신청 사항 &
% #VALUE_ID: LEX-P0045-V0008
% #FIELD_VALUE: 신청 사항 4
% #HANDWRITTEN: 불명확한 기재
\fieldvalue{\handwritten{불명확한 기재}} \\
\hline
5 &
신청 사항 &
% #VALUE_ID: LEX-P0045-V0009
% #FIELD_VALUE: 신청 사항 5
% #HANDWRITTEN: 불명확한 기재
\fieldvalue{\handwritten{불명확한 기재}} \\
\hline
6 &
신청 사항 &
% #VALUE_ID: LEX-P0045-V0010
% #FIELD_VALUE: 신청 사항 6
% #HANDWRITTEN: 불명확한 기재
\fieldvalue{\handwritten{불명확한 기재}} \\
\hline
7 &
신청 사항 &
% #VALUE_ID: LEX-P0045-V0011
% #FIELD_VALUE: 신청 사항 7
% #HANDWRITTEN: 불명확한 기재
\fieldvalue{\handwritten{불명확한 기재}} \\
\hline
8 &
신청 사항 &
% #VALUE_ID: LEX-P0045-V0012
% #FIELD_VALUE: 신청 사항 8
% #HANDWRITTEN: 불명확한 기재
\fieldvalue{\handwritten{불명확한 기재}} \\
\hline
\end{tabularx}

\vspace{0.5em}

\noindent
\begin{tabularx}{\textwidth}{|l|X|}
\hline
확인 사항 &
\begin{minipage}[t]{\linewidth}
① 관련 사실을 확인하였습니까? \(\square\) 예 \(\square\) 아니오

② 제출서류를 확인하였습니까? \(\square\) 예 \(\square\) 아니오

③ 추가 확인이 필요합니까? \(\square\) 예 \(\square\) 아니오
\end{minipage}
\\
\hline
\end{tabularx}

\vfill

\noindent
\begin{tabularx}{\textwidth}{|l|X|X|X|}
\hline
기관명 & 담당자 & 확인일 & 비고 \\
\hline
불명확 &
% #VALUE_ID: LEX-P0045-V0013
% #FIELD_VALUE: 담당자
% #HANDWRITTEN: 불명확한 서명
\fieldvalue{\handwritten{불명확한 서명}} &
% #VALUE_ID: LEX-P0045-V0014
% #FIELD_VALUE: 확인일
% #HANDWRITTEN: 불명확한 날짜
\fieldvalue{\handwritten{불명확한 날짜}} &
\\
\hline
\end{tabularx}

% TODO: The lower edge of the scanned page contains a partially visible form identifier and stamp; the text is not sufficiently legible to transcribe reliably.
% LEXOID_PAGE_COMPLETED: 45/78

\newpage

\begin{flushright}
\textbf{טופס ... (22)}\\
עמוד 2 מתוך 2
\end{flushright}

\begin{tabularx}{\textwidth}{|X|X|X|}
\hline
\textbf{מספר תיק} & \textbf{תאריך} & \textbf{מספר טופס}\\
\hline
% #VALUE_ID: LEX-P0046-V0001
% #FIELD_VALUE: מספר תיק
% #TODO #HANDWRITTEN: [לא קריא]; חותמת/כתב יד אדום בשדה
\fieldvalue{\handwritten{[לא קריא]}}
&
% #VALUE_ID: LEX-P0046-V0002
% #FIELD_VALUE: תאריך
% #TODO #HANDWRITTEN: /; סימון כחול בכתב יד
\fieldvalue{\handwritten{/}}
&
200--480
\\
\hline
\end{tabularx}

\vspace{0.5em}

\begin{tabularx}{\textwidth}{|p{0.47\textwidth}|X|}
\hline
\multicolumn{2}{|c|}{\textbf{פרטי הבקשה}}\\
\hline
\textbf{סוג הבקשה:}
\begin{itemize}
  \item[$\square$] ...
  \item[$\square$] ...
  \item[$\square$] ...
\end{itemize}
&
\textbf{פרטי המבקש:}\\[1em]
\\
\hline
\end{tabularx}

\vspace{0.5em}

\begin{tabularx}{\textwidth}{|p{0.47\textwidth}|X|}
\hline
\textbf{סעיף / פירוט}
&
\textbf{הערות}\\
\hline
\begin{itemize}
  \item[$\square$] ...
  \item[$\square$] ...
  \item[$\square$] ...
  \item[$\square$] ...
  \item[$\square$] ...
  \item[$\square$] ...
\end{itemize}
&
\begin{enumerate}
  \item ...
  \item ...
  \item ...
  \item ...
  \item ...
  \item ...
\end{enumerate}
\\
\hline
\end{tabularx}

\vspace{0.5em}

\begin{tabularx}{\textwidth}{|X|p{0.47\textwidth}|}
\hline
\textbf{חתימה}
&
\textbf{תאריך}\\
\hline
% #VALUE_ID: LEX-P0046-V0003
% #FIELD_VALUE: חתימה
% #TODO #HANDWRITTEN: [לא קריא]; חתימה/רישום כחול בשדה
\fieldvalue{\handwritten{[לא קריא]}}
&
% #VALUE_ID: LEX-P0046-V0004
% #FIELD_VALUE: תאריך
% #TODO #HANDWRITTEN: /; סימון כחול בכתב יד
\fieldvalue{\handwritten{/}}
\\
\hline
\end{tabularx}

\vspace{0.5em}

\begin{tabularx}{\textwidth}{|X|}
\hline
\textbf{הערות נוספות}\\[1.5em]
% #VALUE_ID: LEX-P0046-V0005
% #FIELD_VALUE: הערות נוספות
% #TODO #HANDWRITTEN: [לא קריא]; כתב יד כחול באזור ההערות
\fieldvalue{\handwritten{[לא קריא]}}\\[2em]
\hline
\end{tabularx}

% TODO: הטקסט והסימונים בכתב היד בעמוד מצולמים באיכות נמוכה; יש לאמת את התעתיק מול מסמך המקור.
% LEXOID_PAGE_COMPLETED: 46/78

\newpage

\section*{水污染源监测及监测报告（乙）}

\noindent
监测报告编号：%
% #VALUE_ID: LEX-P0047-V0001
% #FIELD_VALUE: 监测报告编号
% #TODO #HANDWRITTEN: 2010-?; handwriting is partially illegible
\fieldvalue{\handwritten{2010-?}}

\noindent
排污单位名称：%
% #VALUE_ID: LEX-P0047-V0002
% #FIELD_VALUE: 排污单位名称
% #TODO #HANDWRITTEN: （难以辨认）; handwritten text is unclear
\fieldvalue{\handwritten{（难以辨认）}}

\noindent
监测单位：%
% #VALUE_ID: LEX-P0047-V0003
% #FIELD_VALUE: 监测单位
% #TODO #HANDWRITTEN: （难以辨认）; handwritten text is unclear
\fieldvalue{\handwritten{（难以辨认）}}

\begin{tabularx}{\textwidth}{|p{0.25\textwidth}|X|p{0.25\textwidth}|X|}
\hline
监测点位 & 
% #VALUE_ID: LEX-P0047-V0004
% #FIELD_VALUE: 监测点位
% #TODO #HANDWRITTEN: 1\#; handwritten entry is partly obscured
\fieldvalue{\handwritten{1\#}} &
监测日期 &
% #VALUE_ID: LEX-P0047-V0005
% #FIELD_VALUE: 监测日期
% #TODO #HANDWRITTEN: 2010-?; handwriting is unclear
\fieldvalue{\handwritten{2010-?}} \\
\hline
监测时间 &
% #VALUE_ID: LEX-P0047-V0006
% #FIELD_VALUE: 监测时间
% #TODO #HANDWRITTEN: 6:00; last digit is partly unclear
\fieldvalue{\handwritten{6:00}} &
监测频次 &
% #VALUE_ID: LEX-P0047-V0007
% #FIELD_VALUE: 监测频次
% #TODO #HANDWRITTEN: 1; handwritten entry is faint
\fieldvalue{\handwritten{1}} \\
\hline
\end{tabularx}

\subsection*{一、废水排放口监测情况}

\noindent
（1）废水排放口监测情况：%
% #VALUE_ID: LEX-P0047-V0008
% #FIELD_VALUE: 废水排放口监测情况
% #TODO #HANDWRITTEN: 500; unit and surrounding handwriting are unclear
\fieldvalue{\handwritten{500}}

\noindent
排放量：%
% #VALUE_ID: LEX-P0047-V0009
% #FIELD_VALUE: 排放量
% #TODO #HANDWRITTEN: 500 m$^3$; handwriting is partly illegible
\fieldvalue{\handwritten{500 m$^3$}}

\noindent
监测项目及监测结果：

\begin{tabularx}{\textwidth}{|X|c|c|c|c|}
\hline
项目 & 第一次 & 第二次 & 平均值 & 备注 \\
\hline
水温（$^\circ$C） & 13.8 & 18.0 & & \\
\hline
流量 & 0.3 & 0.6 & & \\
\hline
悬浮物 & 0 & 0 & & \\
\hline
其他监测项目 & 1 & 1 & & \\
\hline
\end{tabularx}

\noindent
监测结果补充记录：

\begin{tabularx}{\textwidth}{|p{0.35\textwidth}|X|}
\hline
监测项目 &
监测结果 \\
\hline
排放流量 &
% #VALUE_ID: LEX-P0047-V0010
% #FIELD_VALUE: 排放流量
% #TODO #HANDWRITTEN: 36--?; handwritten calculation is unclear
\fieldvalue{\handwritten{36--?}} \\
\hline
水温 &
% #VALUE_ID: LEX-P0047-V0011
% #FIELD_VALUE: 水温
% #TODO #HANDWRITTEN: 13.8; handwritten value is faint
\fieldvalue{\handwritten{13.8}} \\
\hline
悬浮物 &
% #VALUE_ID: LEX-P0047-V0012
% #FIELD_VALUE: 悬浮物
% #TODO #HANDWRITTEN: 0; handwritten value is faint
\fieldvalue{\handwritten{0}} \\
\hline
化学需氧量 &
% #VALUE_ID: LEX-P0047-V0013
% #FIELD_VALUE: 化学需氧量
% #TODO #HANDWRITTEN: 0.3; handwriting is partly obscured
\fieldvalue{\handwritten{0.3}} \\
\hline
\end{tabularx}

\subsection*{二、监测结果评价}

\noindent
监测结果评价：%
% #VALUE_ID: LEX-P0047-V0014
% #FIELD_VALUE: 监测结果评价
% #TODO #HANDWRITTEN: 合格; handwriting is difficult to read
\fieldvalue{\handwritten{合格}}

\noindent
备注：%
% #VALUE_ID: LEX-P0047-V0015
% #FIELD_VALUE: 备注
% #TODO #HANDWRITTEN: 无; faint handwritten entry
\fieldvalue{\handwritten{无}}

\subsection*{三、审核意见}

\noindent
审核意见：%
% #VALUE_ID: LEX-P0047-V0016
% #FIELD_VALUE: 审核意见
% #TODO #HANDWRITTEN: 同意; handwriting is partially illegible
\fieldvalue{\handwritten{同意}}

\noindent
审核人：%
% #VALUE_ID: LEX-P0047-V0017
% #FIELD_VALUE: 审核人
% #TODO #HANDWRITTEN: （难以辨认）; handwritten signature/name is unclear
\fieldvalue{\handwritten{（难以辨认）}}

\noindent
日期：%
% #VALUE_ID: LEX-P0047-V0018
% #FIELD_VALUE: 日期
% #TODO #HANDWRITTEN: 2010-?; date is partially obscured
\fieldvalue{\handwritten{2010-?}}

% TODO: The source page is rotated and several handwritten entries and labels are partially illegible. The lower portion of the form continues beyond this page boundary.
% LEXOID_PAGE_COMPLETED: 47/78

\newpage

\section*{表單內容}

% TODO：本頁影像為倒置且解析度有限；以下保留可辨識的表格結構，題目文字以可辨識部分轉錄。

\begin{center}
\textbf{（表單標題文字難辨）}
\end{center}

\noindent
填表日期：
% #VALUE_ID: LEX-P0048-V0001
% #FIELD_VALUE: 填表日期
% #TODO #HANDWRITTEN: 難辨；影像倒置且字跡模糊
\fieldvalue{\handwritten{\text{難辨}}}

\vspace{0.5em}

\renewcommand{\arraystretch}{1.5}
\small
\begin{tabularx}{\textwidth}{|c|>{\RaggedRight\arraybackslash}X|>{\RaggedRight\arraybackslash}X|}
\hline
\textbf{題號} & \textbf{題目} & \textbf{作答} \\
\hline
1 &
題目文字難辨。
&
% #VALUE_ID: LEX-P0048-V0002
% #FIELD_VALUE: 第1題作答
% #TODO #HANDWRITTEN: 難辨；藍色手寫字跡模糊
\fieldvalue{\handwritten{\text{難辨}}}
\\
\hline
2 &
題目文字難辨。
&
% #VALUE_ID: LEX-P0048-V0003
% #FIELD_VALUE: 第2題作答
% #TODO #HANDWRITTEN: 難辨；藍色手寫字跡模糊
\fieldvalue{\handwritten{\text{難辨}}}
\\
\hline
3 &
題目文字難辨。
&
% #VALUE_ID: LEX-P0048-V0004
% #FIELD_VALUE: 第3題作答
% #TODO #HANDWRITTEN: 難辨；藍色手寫字跡模糊
\fieldvalue{\handwritten{\text{難辨}}}
\\
\hline
4 &
題目文字難辨。
&
% #VALUE_ID: LEX-P0048-V0005
% #FIELD_VALUE: 第4題作答
% #TODO #HANDWRITTEN: 難辨；藍色手寫字跡模糊
\fieldvalue{\handwritten{\text{難辨}}}
\\
\hline
5 &
題目文字難辨。
&
% #VALUE_ID: LEX-P0048-V0006
% #FIELD_VALUE: 第5題作答
% #TODO #HANDWRITTEN: 難辨；藍色手寫字跡模糊
\fieldvalue{\handwritten{\text{難辨}}}
\\
\hline
6 &
題目文字難辨。
&
% #VALUE_ID: LEX-P0048-V0007
% #FIELD_VALUE: 第6題作答
% #TODO #HANDWRITTEN: 難辨；藍色手寫字跡模糊
\fieldvalue{\handwritten{\text{難辨}}}
\\
\hline
7 &
題目文字難辨。
&
% #VALUE_ID: LEX-P0048-V0008
% #FIELD_VALUE: 第7題作答
% #TODO #HANDWRITTEN: 難辨；此欄含多行藍色手寫內容
\fieldvalue{\handwritten{\text{難辨}}}
\\
\hline
8 &
題目文字難辨。
&
% #VALUE_ID: LEX-P0048-V0009
% #FIELD_VALUE: 第8題作答
% #TODO #HANDWRITTEN: 難辨；藍色手寫字跡模糊
\fieldvalue{\handwritten{\text{難辨}}}
\\
\hline
9 &
題目文字難辨。
&
% #VALUE_ID: LEX-P0048-V0010
% #FIELD_VALUE: 第9題作答
% #TODO #HANDWRITTEN: 難辨；藍色手寫字跡模糊
\fieldvalue{\handwritten{\text{難辨}}}
\\
\hline
10 &
題目文字難辨。
&
% #VALUE_ID: LEX-P0048-V0011
% #FIELD_VALUE: 第10題作答
% #TODO #HANDWRITTEN: 難辨；藍色手寫字跡模糊
\fieldvalue{\handwritten{\text{難辨}}}
\\
\hline
\end{tabularx}

\normalsize

% TODO：表格內容於本頁邊界處仍有延續或受影像方向影響；僅呈現本頁可見部分。
% LEXOID_PAGE_COMPLETED: 48/78

\newpage

\begin{center}
\textbf{（表單頁面；原頁面方向已校正）}
\end{center}

\begin{tabularx}{\textwidth}{|l|X|l|X|}
\hline
案號 & 
% #VALUE_ID: LEX-P0049-V0001
% #FIELD_VALUE: 案號
% #TODO #HANDWRITTEN: [字跡不清]; 藍筆手寫內容難以辨識
\fieldvalue{\handwritten{[字跡不清]}} &
速別 & 普通 \\
\hline
發文日期 & 
% #VALUE_ID: LEX-P0049-V0002
% #FIELD_VALUE: 發文日期
% #TODO #HANDWRITTEN: [日期不清]; 藍筆手寫日期難以辨識
\fieldvalue{\handwritten{[日期不清]}} &
密等及解密條件 & \\
\hline
發文字號 & 
% #VALUE_ID: LEX-P0049-V0003
% #FIELD_VALUE: 發文字號
% #TODO #HANDWRITTEN: [字號不清]; 藍筆手寫字號難以辨識
\fieldvalue{\handwritten{[字號不清]}} &
附件 & \\
\hline
\end{tabularx}

\vspace{0.5em}

\begin{tabularx}{\textwidth}{|p{0.16\textwidth}|X|}
\hline
受文者 & \\
\hline
主旨 &
% #VALUE_ID: LEX-P0049-V0004
% #FIELD_VALUE: 主旨欄手寫內容
% #TODO #HANDWRITTEN: [內容不清]; 藍筆手寫文字難以辨識
\fieldvalue{\handwritten{[內容不清]}} \\
\hline
說明 &
\begin{enumerate}
\item
% #VALUE_ID: LEX-P0049-V0005
% #FIELD_VALUE: 說明一
% #TODO #HANDWRITTEN: [內容不清]; 藍筆手寫文字難以辨識
\fieldvalue{\handwritten{[內容不清]}}
\item
% #VALUE_ID: LEX-P0049-V0006
% #FIELD_VALUE: 說明二
% #TODO #HANDWRITTEN: [內容不清]; 藍筆手寫文字難以辨識
\fieldvalue{\handwritten{[內容不清]}}
\item
% #VALUE_ID: LEX-P0049-V0007
% #FIELD_VALUE: 說明三
% #TODO #HANDWRITTEN: [內容不清]; 藍筆手寫文字難以辨識
\fieldvalue{\handwritten{[內容不清]}}
\end{enumerate}
\\
\hline
\end{tabularx}

\vspace{0.5em}

\begin{tabularx}{\textwidth}{|p{0.16\textwidth}|X|p{0.16\textwidth}|X|}
\hline
承辦單位 & 
% #VALUE_ID: LEX-P0049-V0008
% #FIELD_VALUE: 承辦單位
% #TODO #HANDWRITTEN: [內容不清]; 藍筆手寫內容難以辨識
\fieldvalue{\handwritten{[內容不清]}} &
承辦人 & 
% #VALUE_ID: LEX-P0049-V0009
% #FIELD_VALUE: 承辦人
% #TODO #HANDWRITTEN: [姓名不清]; 藍筆手寫姓名難以辨識
\fieldvalue{\handwritten{[姓名不清]}} \\
\hline
電話 & 
% #VALUE_ID: LEX-P0049-V0010
% #FIELD_VALUE: 電話
% #TODO #HANDWRITTEN: [號碼不清]; 藍筆手寫電話號碼難以辨識
\fieldvalue{\handwritten{[號碼不清]}} &
傳真 & \\
\hline
\end{tabularx}

\vspace{0.5em}

\noindent
% #VALUE_ID: LEX-P0049-V0011
% #FIELD_VALUE: 頁面下方日期欄
% #TODO #HANDWRITTEN: 81/09/02（疑似）; 藍筆手寫日期辨識不確定
\fieldvalue{\handwritten{81/09/02}} \hfill
% #VALUE_ID: LEX-P0049-V0012
% #FIELD_VALUE: 頁面下方簽名欄
% #TODO #HANDWRITTEN: [簽名不清]; 藍筆手寫簽名難以辨識
\fieldvalue{\handwritten{[簽名不清]}} \hfill
% #VALUE_ID: LEX-P0049-V0013
% #FIELD_VALUE: 頁面下方日期／編號欄
% #TODO #HANDWRITTEN: [內容不清]; 藍筆手寫內容難以辨識
\fieldvalue{\handwritten{[內容不清]}} 

\vspace{1em}

\noindent\textbf{附註：}本頁含多處藍筆手寫及勾選內容；因掃描影像方向及解析度限制，無法可靠辨識其完整文字。
% TODO: The visible form text and several handwritten entries are partially illegible in the supplied scan.
% LEXOID_PAGE_COMPLETED: 49/78

\newpage

\begin{center}
\textbf{別紙様式}
\end{center}

\begin{tabularx}{\textwidth}{|l|X|l|}
\hline
文書 & 様式 & 記載例 \\
\hline
別紙 & 手引き & 200--480 \\
\hline
\end{tabularx}

\vspace{0.5em}

\begin{tabularx}{\textwidth}{|p{0.22\textwidth}|X|}
\hline
項目 & 内容 \\
\hline
評価 & 
% #VALUE_ID: LEX-P0050-V0001
% #FIELD_VALUE: 記載番号
% #TODO #HANDWRITTEN: 0009050?; 文字列が不鮮明
\fieldvalue{\handwritten{0009050?}} \\
\hline
\end{tabularx}

\vspace{0.5em}

\noindent
\textbf{点検項目}

\begin{tabularx}{\textwidth}{|p{0.22\textwidth}|X|p{0.20\textwidth}|}
\hline
項目 & 点検内容 & 評価 \\
\hline
総合評価 &
点検結果を記載する。 &
% #VALUE_ID: LEX-P0050-V0002
% #FIELD_VALUE: 総合評価
% #HANDWRITTEN: 35
\fieldvalue{\handwritten{35}} \\
\hline
総合評価 &
評価者記入欄。 &
% #VALUE_ID: LEX-P0050-V0003
% #FIELD_VALUE: 評価者記入欄
% #HANDWRITTEN: 1
\fieldvalue{\handwritten{1}} \\
\hline
\multicolumn{2}{|l|}{点検項目} &
% #VALUE_ID: LEX-P0050-V0004
% #FIELD_VALUE: 点検項目1
% #HANDWRITTEN: 0
\fieldvalue{\handwritten{0}} \\
\hline
\multicolumn{2}{|l|}{点検項目} &
% #VALUE_ID: LEX-P0050-V0005
% #FIELD_VALUE: 点検項目2
% #HANDWRITTEN: 0
\fieldvalue{\handwritten{0}} \\
\hline
\multicolumn{2}{|l|}{点検項目} &
% #VALUE_ID: LEX-P0050-V0006
% #FIELD_VALUE: 点検項目3
% #HANDWRITTEN: 0
\fieldvalue{\handwritten{0}} \\
\hline
\multicolumn{2}{|l|}{点検項目} &
% #VALUE_ID: LEX-P0050-V0007
% #FIELD_VALUE: 点検項目4
% #HANDWRITTEN: 0
\fieldvalue{\handwritten{0}} \\
\hline
\multicolumn{2}{|l|}{点検項目} &
% #VALUE_ID: LEX-P0050-V0008
% #FIELD_VALUE: 点検項目5
% #HANDWRITTEN: 0
\fieldvalue{\handwritten{0}} \\
\hline
\multicolumn{2}{|l|}{点検項目} &
% #VALUE_ID: LEX-P0050-V0009
% #FIELD_VALUE: 点検項目6
% #HANDWRITTEN: 1
\fieldvalue{\handwritten{1}} \\
\hline
\multicolumn{2}{|l|}{点検項目} &
% #VALUE_ID: LEX-P0050-V0010
% #FIELD_VALUE: 点検項目7
% #HANDWRITTEN: 0
\fieldvalue{\handwritten{0}} \\
\hline
\multicolumn{2}{|l|}{点検項目} &
% #VALUE_ID: LEX-P0050-V0011
% #FIELD_VALUE: 点検項目8
% #HANDWRITTEN: 0
\fieldvalue{\handwritten{0}} \\
\hline
\multicolumn{2}{|l|}{点検項目} &
% #VALUE_ID: LEX-P0050-V0012
% #FIELD_VALUE: 点検項目9
% #HANDWRITTEN: 0
\fieldvalue{\handwritten{0}} \\
\hline
\multicolumn{2}{|l|}{点検項目} &
% #VALUE_ID: LEX-P0050-V0013
% #FIELD_VALUE: 点検項目10
% #HANDWRITTEN: 0
\fieldvalue{\handwritten{0}} \\
\hline
\multicolumn{2}{|l|}{点検項目} &
% #VALUE_ID: LEX-P0050-V0014
% #FIELD_VALUE: 点検項目11
% #TODO #HANDWRITTEN: 4?; 手書き文字が不鮮明
\fieldvalue{\handwritten{4?}} \\
\hline
\multicolumn{2}{|l|}{点検項目} &
% #VALUE_ID: LEX-P0050-V0015
% #FIELD_VALUE: 点検項目12
% #HANDWRITTEN: 0
\fieldvalue{\handwritten{0}} \\
\hline
\multicolumn{2}{|l|}{点検項目} &
% #VALUE_ID: LEX-P0050-V0016
% #FIELD_VALUE: 点検項目13
% #HANDWRITTEN: 0
\fieldvalue{\handwritten{0}} \\
\hline
\multicolumn{2}{|l|}{点検項目} &
% #VALUE_ID: LEX-P0050-V0017
% #FIELD_VALUE: 点検項目14
% #HANDWRITTEN: 0
\fieldvalue{\handwritten{0}} \\
\hline
\multicolumn{2}{|l|}{点検項目} &
% #VALUE_ID: LEX-P0050-V0018
% #FIELD_VALUE: 点検項目15
% #HANDWRITTEN: 0
\fieldvalue{\handwritten{0}} \\
\hline
\multicolumn{2}{|l|}{点検項目} &
% #VALUE_ID: LEX-P0050-V0019
% #FIELD_VALUE: 点検項目16
% #HANDWRITTEN: 0
\fieldvalue{\handwritten{0}} \\
\hline
\multicolumn{2}{|l|}{点検項目} &
% #VALUE_ID: LEX-P0050-V0020
% #FIELD_VALUE: 点検項目17
% #HANDWRITTEN: 0
\fieldvalue{\handwritten{0}} \\
\hline
\multicolumn{2}{|l|}{点検項目} &
% #VALUE_ID: LEX-P0050-V0021
% #FIELD_VALUE: 点検項目18
% #HANDWRITTEN: 0
\fieldvalue{\handwritten{0}} \\
\hline
\multicolumn{2}{|l|}{点検項目} &
% #VALUE_ID: LEX-P0050-V0022
% #FIELD_VALUE: 点検項目19
% #TODO #HANDWRITTEN: 2?; 手書き文字が不鮮明
\fieldvalue{\handwritten{2?}} \\
\hline
\multicolumn{2}{|l|}{点検項目} &
% #VALUE_ID: LEX-P0050-V0023
% #FIELD_VALUE: 点検項目20
% #HANDWRITTEN: 0
\fieldvalue{\handwritten{0}} \\
\hline
\end{tabularx}

% TODO: The source page is rotated and several printed Japanese labels are too
% low-resolution to transcribe reliably; the visible handwritten evaluation
% values have been preserved in visual reading order.
% LEXOID_PAGE_COMPLETED: 50/78

\newpage

% TODO: The page is a Hebrew form; several fixed labels are difficult to read at the
% available resolution. The checklist below contains only rows visible on this page.

\begin{center}
\textbf{טופס בדיקה}
\end{center}

\noindent
\begin{tabularx}{\textwidth}{|X|c|}
\hline
\multicolumn{2}{|c|}{\textbf{רשימת בדיקה}}\\
\cline{1-2}
\text{סעיף בדיקה} & \text{ערך}\\
\hline
% #VALUE_ID: LEX-P0051-V0001
% #FIELD_VALUE: ערך בשורת הבדיקה 1
% #HANDWRITTEN: 0
\text{סעיף בדיקה} & \fieldvalue{\handwritten{0}}\\
\hline
% #VALUE_ID: LEX-P0051-V0002
% #FIELD_VALUE: ערך בשורת הבדיקה 2
% #HANDWRITTEN: 0
\text{סעיף בדיקה} & \fieldvalue{\handwritten{0}}\\
\hline
% #VALUE_ID: LEX-P0051-V0003
% #FIELD_VALUE: ערך בשורת הבדיקה 3
% #HANDWRITTEN: 0
\text{סעיף בדיקה} & \fieldvalue{\handwritten{0}}\\
\hline
% #VALUE_ID: LEX-P0051-V0004
% #FIELD_VALUE: ערך בשורת הבדיקה 4
% #HANDWRITTEN: 0
\text{סעיף בדיקה} & \fieldvalue{\handwritten{0}}\\
\hline
% #VALUE_ID: LEX-P0051-V0005
% #FIELD_VALUE: ערך בשורת הבדיקה 5
% #HANDWRITTEN: 0
\text{סעיף בדיקה} & \fieldvalue{\handwritten{0}}\\
\hline
% #VALUE_ID: LEX-P0051-V0006
% #FIELD_VALUE: ערך בשורת הבדיקה 6
% #HANDWRITTEN: 0
\text{סעיף בדיקה} & \fieldvalue{\handwritten{0}}\\
\hline
% #VALUE_ID: LEX-P0051-V0007
% #FIELD_VALUE: ערך בשורת הבדיקה 7
% #HANDWRITTEN: 0
\text{סעיף בדיקה} & \fieldvalue{\handwritten{0}}\\
\hline
% #VALUE_ID: LEX-P0051-V0008
% #FIELD_VALUE: ערך בשורת הבדיקה 8
% #HANDWRITTEN: 0
\text{סעיף בדיקה} & \fieldvalue{\handwritten{0}}\\
\hline
% #VALUE_ID: LEX-P0051-V0009
% #FIELD_VALUE: ערך בשורת הבדיקה 9
% #HANDWRITTEN: 0
\text{סעיף בדיקה} & \fieldvalue{\handwritten{0}}\\
\hline
% #VALUE_ID: LEX-P0051-V0010
% #FIELD_VALUE: ערך בשורת הבדיקה 10
% #HANDWRITTEN: 0
\text{סעיף בדיקה} & \fieldvalue{\handwritten{0}}\\
\hline
% #VALUE_ID: LEX-P0051-V0011
% #FIELD_VALUE: ערך בשורת הבדיקה 11
% #HANDWRITTEN: 0
\text{סעיף בדיקה} & \fieldvalue{\handwritten{0}}\\
\hline
% #VALUE_ID: LEX-P0051-V0012
% #FIELD_VALUE: ערך בשורת הבדיקה 12
% #HANDWRITTEN: 0
\text{סעיף בדיקה} & \fieldvalue{\handwritten{0}}\\
\hline
% #VALUE_ID: LEX-P0051-V0013
% #FIELD_VALUE: ערך בשורת הבדיקה 13
% #HANDWRITTEN: 0
\text{סעיף בדיקה} & \fieldvalue{\handwritten{0}}\\
\hline
% #VALUE_ID: LEX-P0051-V0014
% #FIELD_VALUE: ערך בשורת הבדיקה 14
% #HANDWRITTEN: 0
\text{סעיף בדיקה} & \fieldvalue{\handwritten{0}}\\
\hline
\end{tabularx}

\vspace{1em}

\noindent
\begin{tabularx}{\textwidth}{|X|}
\hline
\textbf{הערות:}\\[17em]
\hline
\end{tabularx}

\vspace{1em}

\noindent
\begin{tabularx}{\textwidth}{|X|X|}
\hline
\text{חתימה} & \text{תאריך}\\
\hline
% #VALUE_ID: LEX-P0051-V0015
% #FIELD_VALUE: חתימה
% TODO #HANDWRITTEN: illegible signature; handwritten blue signature
\fieldvalue{\handwritten{[חתימה לא קריאה]}}
&
% #VALUE_ID: LEX-P0051-V0016
% #FIELD_VALUE: תאריך
% TODO #HANDWRITTEN: 18/01/50?; handwritten date is unclear
\fieldvalue{\handwritten{18/01/50?}}\\
\hline
% #VALUE_ID: LEX-P0051-V0017
% #FIELD_VALUE: חתימה
% TODO #HANDWRITTEN: illegible signature; handwritten blue signature
\fieldvalue{\handwritten{[חתימה לא קריאה]}}
&
% #VALUE_ID: LEX-P0051-V0018
% #FIELD_VALUE: תאריך
% TODO #HANDWRITTEN: 18/09/??; handwritten date is unclear
\fieldvalue{\handwritten{18/09/??}}\\
\hline
\end{tabularx}

\vspace{1em}

\noindent
\begin{tabularx}{\textwidth}{|X|X|X|}
\hline
\text{מספר טופס} & \text{שם הטופס} & \text{קוד}\\
\hline
\text{[טקסט מודפס]} & \text{[טקסט מודפס]} & \text{[טקסט מודפס]}\\
\hline
\end{tabularx}

% TODO: A red stamp and additional printed administrative text are visible at the
% bottom of the page; the stamp text is not reliably legible.
% LEXOID_PAGE_COMPLETED: 51/78

\newpage

% TODO: The page is a Japanese administrative form. Several fixed labels are faint or partially illegible.
\section*{実績報告書}

\begin{tabularx}{\textwidth}{|p{0.24\textwidth}|X|}
\hline
\textbf{項目} & \textbf{内容} \\
\hline
事業の概要 &
\begin{minipage}[t]{\linewidth}
① 事業の概要について記載する欄。\\[2pt]
② 事業の実施状況について記載する欄。\\[2pt]
③ 事業の成果について記載する欄。\\[2pt]
④ 事業の対象者について記載する欄。\\[2pt]
⑤ 事業の経費について記載する欄。\\[2pt]
⑥ その他について記載する欄。
\end{minipage}
\\
\hline
\end{tabularx}

\vspace{0.5em}

\begin{tabularx}{\textwidth}{|p{0.28\textwidth}|X|}
\hline
\textbf{記載欄} & \textbf{記載内容} \\
\hline
% #VALUE_ID: LEX-P0052-V0001
% #FIELD_VALUE: 日付
% #HANDWRITTEN: 令和202?年?月?日
\fieldvalue{\handwritten{令和202?年?月?日}} &
\\
\hline
% #VALUE_ID: LEX-P0052-V0002
% #FIELD_VALUE: 項目1
% #HANDWRITTEN: 2
\fieldvalue{\handwritten{2}} &
% #VALUE_ID: LEX-P0052-V0003
% #FIELD_VALUE: 項目2
% #HANDWRITTEN: 5
\fieldvalue{\handwritten{5}} \\
\hline
% #VALUE_ID: LEX-P0052-V0004
% #FIELD_VALUE: 項目3
% #HANDWRITTEN: 0
\fieldvalue{\handwritten{0}} &
\\
\hline
% #VALUE_ID: LEX-P0052-V0005
% #FIELD_VALUE: 項目4
% #HANDWRITTEN: 0
\fieldvalue{\handwritten{0}} &
\\
\hline
% #VALUE_ID: LEX-P0052-V0006
% #FIELD_VALUE: 項目5
% #HANDWRITTEN: 0
\fieldvalue{\handwritten{0}} &
\\
\hline
% #VALUE_ID: LEX-P0052-V0007
% #FIELD_VALUE: 項目6
% #HANDWRITTEN: 0
\fieldvalue{\handwritten{0}} &
\\
\hline
% #VALUE_ID: LEX-P0052-V0008
% #FIELD_VALUE: 項目7
% #HANDWRITTEN: 0
\fieldvalue{\handwritten{0}} &
\\
\hline
% #VALUE_ID: LEX-P0052-V0009
% #FIELD_VALUE: 項目8
% #HANDWRITTEN: 0
\fieldvalue{\handwritten{0}} &
\\
\hline
% #VALUE_ID: LEX-P0052-V0010
% #FIELD_VALUE: 項目9
% #HANDWRITTEN: 0
\fieldvalue{\handwritten{0}} &
\\
\hline
% #VALUE_ID: LEX-P0052-V0011
% #FIELD_VALUE: 項目10
% #HANDWRITTEN: 0
\fieldvalue{\handwritten{0}} &
\\
\hline
% #VALUE_ID: LEX-P0052-V0012
% #FIELD_VALUE: 項目11
% #HANDWRITTEN: 0
\fieldvalue{\handwritten{0}} &
\\
\hline
% #VALUE_ID: LEX-P0052-V0013
% #FIELD_VALUE: 項目12
% #HANDWRITTEN: 0
\fieldvalue{\handwritten{0}} &
\\
\hline
% #VALUE_ID: LEX-P0052-V0014
% #FIELD_VALUE: 項目13
% #HANDWRITTEN: 0
\fieldvalue{\handwritten{0}} &
\\
\hline
% #VALUE_ID: LEX-P0052-V0015
% #FIELD_VALUE: 項目14
% #HANDWRITTEN: 0
\fieldvalue{\handwritten{0}} &
\\
\hline
% #VALUE_ID: LEX-P0052-V0016
% #FIELD_VALUE: 項目15
% #HANDWRITTEN: 0
\fieldvalue{\handwritten{0}} &
\\
\hline
% #VALUE_ID: LEX-P0052-V0017
% #FIELD_VALUE: 項目16
% #HANDWRITTEN: 0
\fieldvalue{\handwritten{0}} &
\\
\hline
% #VALUE_ID: LEX-P0052-V0018
% #FIELD_VALUE: 項目17
% #HANDWRITTEN: 0
\fieldvalue{\handwritten{0}} &
\\
\hline
% #VALUE_ID: LEX-P0052-V0019
% #FIELD_VALUE: 項目18
% #HANDWRITTEN: 0
\fieldvalue{\handwritten{0}} &
\\
\hline
% #VALUE_ID: LEX-P0052-V0020
% #FIELD_VALUE: 項目19
% #HANDWRITTEN: 0
\fieldvalue{\handwritten{0}} &
\\
\hline
% #VALUE_ID: LEX-P0052-V0021
% #FIELD_VALUE: 項目20
% #HANDWRITTEN: 0
\fieldvalue{\handwritten{0}} &
\\
\hline
% #VALUE_ID: LEX-P0052-V0022
% #FIELD_VALUE: 項目21
% #HANDWRITTEN: 0
\fieldvalue{\handwritten{0}} &
\\
\hline
% #VALUE_ID: LEX-P0052-V0023
% #FIELD_VALUE: 項目22
% #HANDWRITTEN: 0
\fieldvalue{\handwritten{0}} &
\\
\hline
% #VALUE_ID: LEX-P0052-V0024
% #FIELD_VALUE: 項目23
% #HANDWRITTEN: 0
\fieldvalue{\handwritten{0}} &
\\
\hline
% #VALUE_ID: LEX-P0052-V0025
% #FIELD_VALUE: 項目24
% #HANDWRITTEN: 0
\fieldvalue{\handwritten{0}} &
\\
\hline
\end{tabularx}

% TODO: The form continues beyond the visible page boundary.
% LEXOID_PAGE_COMPLETED: 52/78

\newpage

\begin{center}
\textbf{医疗废物管理检查表（2025年）}
\end{center}

\noindent
\begin{tabularx}{\textwidth}{|l|X|l|X|}
\hline
被检查单位 & 
% #VALUE_ID: LEX-P0053-V0001
% #FIELD_VALUE: 被检查单位
% #TODO #HANDWRITTEN: 难辨；蓝色手写内容模糊
\fieldvalue{\handwritten{难辨}} &
检查日期 &
% #VALUE_ID: LEX-P0053-V0002
% #FIELD_VALUE: 检查日期
% #HANDWRITTEN: 2025.6.18
\fieldvalue{\handwritten{2025.6.18}}
\\
\hline
检查人 &
% #VALUE_ID: LEX-P0053-V0003
% #FIELD_VALUE: 检查人
% #HANDWRITTEN: 难辨
\fieldvalue{\handwritten{难辨}} &
编号 &
% #VALUE_ID: LEX-P0053-V0004
% #FIELD_VALUE: 编号
% #TODO #HANDWRITTEN: /；仅能辨认出斜线
\fieldvalue{\handwritten{/}}
\\
\hline
\end{tabularx}

\vspace{0.5em}

\noindent
\begin{tabularx}{\textwidth}{|c|X|c|}
\hline
序号 & 检查项目 & 扣分/得分 \\
\hline
1 &
医疗废物分类收集及标识情况 &
% #VALUE_ID: LEX-P0053-V0005
% #FIELD_VALUE: 第1项扣分/得分
% #HANDWRITTEN: 0
\fieldvalue{\handwritten{0}}
\\
\hline
2 &
医疗废物分类包装情况 &
% #VALUE_ID: LEX-P0053-V0006
% #FIELD_VALUE: 第2项扣分/得分
% #HANDWRITTEN: 0
\fieldvalue{\handwritten{0}}
\\
\hline
3 &
医疗废物包装物及容器使用情况 &
% #VALUE_ID: LEX-P0053-V0007
% #FIELD_VALUE: 第3项扣分/得分
% #HANDWRITTEN: 0
\fieldvalue{\handwritten{0}}
\\
\hline
4 &
医疗废物交接登记情况 &
% #VALUE_ID: LEX-P0053-V0008
% #FIELD_VALUE: 第4项扣分/得分
% #HANDWRITTEN: 0
\fieldvalue{\handwritten{0}}
\\
\hline
5 &
医疗废物暂存设施及管理情况 &
% #VALUE_ID: LEX-P0053-V0009
% #FIELD_VALUE: 第5项扣分/得分
% #HANDWRITTEN: 0
\fieldvalue{\handwritten{0}}
\\
\hline
6 &
医疗废物运送管理情况 &
% #VALUE_ID: LEX-P0053-V0010
% #FIELD_VALUE: 第6项扣分/得分
% #HANDWRITTEN: 0
\fieldvalue{\handwritten{0}}
\\
\hline
7 &
医疗废物处置情况 &
% #VALUE_ID: LEX-P0053-V0011
% #FIELD_VALUE: 第7项扣分/得分
% #HANDWRITTEN: 0
\fieldvalue{\handwritten{0}}
\\
\hline
8 &
医疗废物管理制度落实情况 &
% #VALUE_ID: LEX-P0053-V0012
% #FIELD_VALUE: 第8项扣分/得分
% #HANDWRITTEN: 0
\fieldvalue{\handwritten{0}}
\\
\hline
9 &
相关人员培训及防护情况 &
% #VALUE_ID: LEX-P0053-V0013
% #FIELD_VALUE: 第9项扣分/得分
% #HANDWRITTEN: 0
\fieldvalue{\handwritten{0}}
\\
\hline
10 &
医疗废物资料及记录情况 &
% #VALUE_ID: LEX-P0053-V0014
% #FIELD_VALUE: 第10项扣分/得分
% #HANDWRITTEN: 0
\fieldvalue{\handwritten{0}}
\\
\hline
11 &
医疗废物分类目录及警示标识情况 &
% #VALUE_ID: LEX-P0053-V0015
% #FIELD_VALUE: 第11项扣分/得分
% #HANDWRITTEN: 0
\fieldvalue{\handwritten{0}}
\\
\hline
12 &
医疗废物收集、贮存过程管理情况 &
% #VALUE_ID: LEX-P0053-V0016
% #FIELD_VALUE: 第12项扣分/得分
% #HANDWRITTEN: 0
\fieldvalue{\handwritten{0}}
\\
\hline
13 &
医疗废物专用包装物管理情况 &
% #VALUE_ID: LEX-P0053-V0017
% #FIELD_VALUE: 第13项扣分/得分
% #HANDWRITTEN: 0
\fieldvalue{\handwritten{0}}
\\
\hline
14 &
其他检查项目 &
% #VALUE_ID: LEX-P0053-V0018
% #FIELD_VALUE: 第14项扣分/得分
% #TODO #HANDWRITTEN: 6；蓝色手写数字辨认为6
\fieldvalue{\handwritten{6}}
\\
\hline
\end{tabularx}

\vspace{0.5em}

\noindent
备注：\rule{0.8\textwidth}{0.4pt}

\vfill

\noindent
% TODO：页面底部存在印章及表格文字，部分内容因扫描方向及分辨率难以辨认。
% LEXOID_PAGE_COMPLETED: 53/78

\newpage

\section*{粉じん障害防止規則様式第2号}
\textit{（第22条関係）}

\begin{tabularx}{\textwidth}{|p{0.22\textwidth}|X|p{0.22\textwidth}|}
\hline
事業場の名称 & \multicolumn{2}{l|}{\rule{0pt}{2.2ex}}\\
\hline
所在地 & \multicolumn{2}{l|}{\rule{0pt}{2.2ex}}\\
\hline
\end{tabularx}

\vspace{1em}

\begin{tabularx}{\textwidth}{|p{0.47\textwidth}|X|}
\hline
\multicolumn{2}{|c|}{\textbf{粉じん作業に係る測定結果等}}\\
\hline
{\begin{minipage}[t]{\linewidth}
\textcircled{1}\quad 測定年月日

\vspace{0.4em}
\begin{tabularx}{\linewidth}{@{}lXlX@{}}
月&
% #VALUE_ID: LEX-P0054-V0001
% #FIELD_VALUE: 測定年月日の月
\fieldvalue{\handwritten{12}}&
日&
% #VALUE_ID: LEX-P0054-V0002
% #FIELD_VALUE: 測定年月日の日
\fieldvalue{\handwritten{25}}\\
\end{tabularx}

\vspace{0.7em}
\textcircled{2}\quad 測定結果

\begin{tabularx}{\linewidth}{@{}lX@{}}
第1項目&
% #VALUE_ID: LEX-P0054-V0003
% #FIELD_VALUE: ② 第1項目
\fieldvalue{\handwritten{01}}\\
第2項目&
% #VALUE_ID: LEX-P0054-V0004
% #FIELD_VALUE: ② 第2項目
\fieldvalue{\handwritten{0}}\\
第3項目&
% #VALUE_ID: LEX-P0054-V0005
% #FIELD_VALUE: ② 第3項目
\fieldvalue{\handwritten{0}}\\
第4項目&
% #VALUE_ID: LEX-P0054-V0006
% #FIELD_VALUE: ② 第4項目
\fieldvalue{\handwritten{0}}\\
\end{tabularx}

\vspace{0.7em}
\textcircled{3}\quad 測定結果

\begin{tabularx}{\linewidth}{@{}lX@{}}
第1項目&
% #VALUE_ID: LEX-P0054-V0007
% #FIELD_VALUE: ③ 第1項目
\fieldvalue{\handwritten{01}}\\
第2項目&
% #VALUE_ID: LEX-P0054-V0008
% #FIELD_VALUE: ③ 第2項目
\fieldvalue{\handwritten{0}}\\
第3項目&
% #VALUE_ID: LEX-P0054-V0009
% #FIELD_VALUE: ③ 第3項目
\fieldvalue{\handwritten{0}}\\
第4項目&
% #VALUE_ID: LEX-P0054-V0010
% #FIELD_VALUE: ③ 第4項目
\fieldvalue{\handwritten{0}}\\
\end{tabularx}

\vspace{0.7em}
\textcircled{4}\quad 測定結果

\begin{tabularx}{\linewidth}{@{}lX@{}}
第1項目&
% #VALUE_ID: LEX-P0054-V0011
% #FIELD_VALUE: ④ 第1項目
\fieldvalue{\handwritten{01}}\\
第2項目&
% #VALUE_ID: LEX-P0054-V0012
% #FIELD_VALUE: ④ 第2項目
\fieldvalue{\handwritten{0}}\\
第3項目&
% #VALUE_ID: LEX-P0054-V0013
% #FIELD_VALUE: ④ 第3項目
\fieldvalue{\handwritten{0}}\\
第4項目&
% #VALUE_ID: LEX-P0054-V0014
% #FIELD_VALUE: ④ 第4項目
\fieldvalue{\handwritten{0}}\\
\end{tabularx}

\vspace{0.7em}
\textcircled{5}\quad 測定結果

\begin{tabularx}{\linewidth}{@{}lX@{}}
第1項目&
% #VALUE_ID: LEX-P0054-V0015
% #FIELD_VALUE: ⑤ 第1項目
\fieldvalue{\handwritten{01}}\\
第2項目&
% #VALUE_ID: LEX-P0054-V0016
% #FIELD_VALUE: ⑤ 第2項目
\fieldvalue{\handwritten{0}}\\
第3項目&
% #VALUE_ID: LEX-P0054-V0017
% #FIELD_VALUE: ⑤ 第3項目
\fieldvalue{\handwritten{0}}\\
第4項目&
% #VALUE_ID: LEX-P0054-V0018
% #FIELD_VALUE: ⑤ 第4項目
\fieldvalue{\handwritten{0}}\\
\end{tabularx}

\vspace{0.7em}
\textcircled{6}\quad 測定結果

\begin{tabularx}{\linewidth}{@{}lX@{}}
第1項目&
% #VALUE_ID: LEX-P0054-V0019
% #FIELD_VALUE: ⑥ 第1項目
\fieldvalue{\handwritten{01}}\\
第2項目&
% #VALUE_ID: LEX-P0054-V0020
% #FIELD_VALUE: ⑥ 第2項目
\fieldvalue{\handwritten{0}}\\
第3項目&
% #VALUE_ID: LEX-P0054-V0021
% #FIELD_VALUE: ⑥ 第3項目
\fieldvalue{\handwritten{0}}\\
第4項目&
% #VALUE_ID: LEX-P0054-V0022
% #FIELD_VALUE: ⑥ 第4項目
\fieldvalue{\handwritten{0}}\\
\end{tabularx}
\end{minipage}}
&
\begin{minipage}[t]{\linewidth}
\textbf{記載上の注意}

\begin{enumerate}
\item 測定を行った作業場所、作業内容および測定結果を記載すること。
\item 測定結果は、測定年月日ごとに記載すること。
\item その他必要な事項を記載すること。
\end{enumerate}

\vspace{1em}
% TODO: The explanatory text in this column is partially illegible in the scan; only its visible structure is represented.
\end{minipage}
\\
\hline
\end{tabularx}

\vspace{1em}

\begin{tabularx}{\textwidth}{|p{0.24\textwidth}|X|p{0.24\textwidth}|}
\hline
提出年月日 & \multicolumn{2}{l|}{\rule{0pt}{2.2ex}}\\
\hline
所轄労働基準監督署長殿 & \multicolumn{2}{l|}{\rule{0pt}{2.2ex}}\\
\hline
\end{tabularx}

\noindent
% TODO: The lower administrative stamp and printed office information are visible on the page but are too small to transcribe reliably.
% LEXOID_PAGE_COMPLETED: 54/78

\newpage

\noindent
\begin{center}
\textbf{事業実績報告書}
\end{center}

\noindent
\begin{tabularx}{\textwidth}{|l|X|l|l|}
\hline
年度 & 事業名 & 所管課 & 受付印 \\
\hline
令和元年度 & 国際交流事業 &  &  \\
\hline
\end{tabularx}

\vspace{2mm}

\noindent
\begin{tabularx}{\textwidth}{|l|X|X|}
\hline
団体名 &  & 備考 \\
\hline
\multicolumn{2}{|l|}{実績内容} & \\
\cline{1-2}
\multicolumn{2}{|l|}{\begin{minipage}[t]{0.62\textwidth}
\small
\begin{tabularx}{\linewidth}{@{}lXr@{}}
 & 項目 & 実績 \\
\cline{2-3}
 & 事業参加者数 & 
% #VALUE_ID: LEX-P0055-V0001
% #FIELD_VALUE: 事業参加者数
% #HANDWRITTEN: 0
\fieldvalue{\handwritten{0}} \\
 & 事業実施回数 &
% #VALUE_ID: LEX-P0055-V0002
% #FIELD_VALUE: 事業実施回数
% #HANDWRITTEN: 0
\fieldvalue{\handwritten{0}} \\
 & 参加者数（一般） &
% #VALUE_ID: LEX-P0055-V0003
% #FIELD_VALUE: 参加者数（一般）
% #HANDWRITTEN: 0
\fieldvalue{\handwritten{0}} \\
 & 参加者数（学生） &
% #VALUE_ID: LEX-P0055-V0004
% #FIELD_VALUE: 参加者数（学生）
% #HANDWRITTEN: 0
\fieldvalue{\handwritten{0}} \\
 & 参加者数（その他） &
% #VALUE_ID: LEX-P0055-V0005
% #FIELD_VALUE: 参加者数（その他）
% #HANDWRITTEN: 0
\fieldvalue{\handwritten{0}} \\
 & 事業費総額 &
% #VALUE_ID: LEX-P0055-V0006
% #FIELD_VALUE: 事業費総額
% #HANDWRITTEN: 0
\fieldvalue{\handwritten{0}} \\
 & 補助対象経費 &
% #VALUE_ID: LEX-P0055-V0007
% #FIELD_VALUE: 補助対象経費
% #HANDWRITTEN: 0
\fieldvalue{\handwritten{0}} \\
 & 自己負担額 &
% #VALUE_ID: LEX-P0055-V0008
% #FIELD_VALUE: 自己負担額
% #HANDWRITTEN: 0
\fieldvalue{\handwritten{0}} \\
 & その他収入 &
% #VALUE_ID: LEX-P0055-V0009
% #FIELD_VALUE: その他収入
% #HANDWRITTEN: 0
\fieldvalue{\handwritten{0}} \\
 & 支出合計 &
% #VALUE_ID: LEX-P0055-V0010
% #FIELD_VALUE: 支出合計
% #HANDWRITTEN: 0
\fieldvalue{\handwritten{0}} \\
 & 事業費残額 &
% #VALUE_ID: LEX-P0055-V0011
% #FIELD_VALUE: 事業費残額
% #HANDWRITTEN: 0
\fieldvalue{\handwritten{0}} \\
 & 執行率 &
% #VALUE_ID: LEX-P0055-V0012
% #FIELD_VALUE: 執行率
% #HANDWRITTEN: 0
\fieldvalue{\handwritten{0}} \\
 & 補助金額 &
% #VALUE_ID: LEX-P0055-V0013
% #FIELD_VALUE: 補助金額
% #HANDWRITTEN: 0
\fieldvalue{\handwritten{0}} \\
 & 交付済額 &
% #VALUE_ID: LEX-P0055-V0014
% #FIELD_VALUE: 交付済額
% #HANDWRITTEN: 0
\fieldvalue{\handwritten{0}} \\
\end{tabularx}
\end{minipage}} & \begin{minipage}[t]{0.30\textwidth}\vspace{35mm}\end{minipage} \\
\hline
\end{tabularx}

\vspace{3mm}

\noindent
\begin{tabularx}{\textwidth}{|l|X|l|X|}
\hline
事業実施期間 & 
% #VALUE_ID: LEX-P0055-V0015
% #FIELD_VALUE: 事業実施期間
% #HANDWRITTEN: 令和元年
\fieldvalue{\handwritten{令和元年}} & ～ &
% #VALUE_ID: LEX-P0055-V0016
% #FIELD_VALUE: 事業実施期間（終了）
% #HANDWRITTEN: 令和元年
\fieldvalue{\handwritten{令和元年}} \\
\hline
実績報告年月日 &
% #VALUE_ID: LEX-P0055-V0017
% #FIELD_VALUE: 実績報告年月日
% #HANDWRITTEN: 2019.10.01
\fieldvalue{\handwritten{2019.10.01}} &  & \\
\hline
\end{tabularx}

\vspace{3mm}

\noindent
\textbf{備考}

\noindent
% #VALUE_ID: LEX-P0055-V0018
% #FIELD_VALUE: 備考欄
% #HANDWRITTEN: （判読困難）
\fieldvalue{\handwritten{（判読困難）}}

% TODO: The source page is rotated and several Japanese labels and handwritten entries are partially illegible. Verify the transcriptions against the original scan.
% LEXOID_PAGE_COMPLETED: 55/78

\newpage

\section*{检查表}

% TODO：页面扫描方向及部分中文文字较模糊，以下表格保留可辨认的版式与字段。

\begin{tabularx}{\textwidth}{|l|X|l|}
\hline
\multicolumn{3}{|c|}{\textbf{检查表}}\\
\hline
项目编号 &
\multicolumn{2}{l|}{%
% #VALUE_ID: LEX-P0056-V0001
% #FIELD_VALUE: 项目编号
% #HANDWRITTEN: 24
\fieldvalue{\handwritten{24}}%
}\\
\hline
检查日期 &
\multicolumn{2}{l|}{%
% #VALUE_ID: LEX-P0056-V0002
% #FIELD_VALUE: 检查日期
% TODO #HANDWRITTEN: 200?; 日期字迹及扫描质量不清
\fieldvalue{\handwritten{200?}}%
}\\
\hline
\textbf{项目} & \textbf{检查内容} & \textbf{检查结果}\\
\hline
1 & \textit{检查内容（扫描文字不清）} &
% #VALUE_ID: LEX-P0056-V0003
% #FIELD_VALUE: 第1项检查结果
% #HANDWRITTEN: 0
\fieldvalue{\handwritten{0}}\\
\hline
2 & \textit{检查内容（扫描文字不清）} &
% #VALUE_ID: LEX-P0056-V0004
% #FIELD_VALUE: 第2项检查结果
% #HANDWRITTEN: 0
\fieldvalue{\handwritten{0}}\\
\hline
3 & \textit{检查内容（扫描文字不清）} &
% #VALUE_ID: LEX-P0056-V0005
% #FIELD_VALUE: 第3项检查结果
% #HANDWRITTEN: 0
\fieldvalue{\handwritten{0}}\\
\hline
4 & \textit{检查内容（扫描文字不清）} &
% #VALUE_ID: LEX-P0056-V0006
% #FIELD_VALUE: 第4项检查结果
% #HANDWRITTEN: 0
\fieldvalue{\handwritten{0}}\\
\hline
5 & \textit{检查内容（扫描文字不清）} &
% #VALUE_ID: LEX-P0056-V0007
% #FIELD_VALUE: 第5项检查结果
% #HANDWRITTEN: 0
\fieldvalue{\handwritten{0}}\\
\hline
6 & \textit{检查内容（扫描文字不清）} &
% #VALUE_ID: LEX-P0056-V0008
% #FIELD_VALUE: 第6项检查结果
% #HANDWRITTEN: 0
\fieldvalue{\handwritten{0}}\\
\hline
7 & \textit{检查内容（扫描文字不清）} &
% #VALUE_ID: LEX-P0056-V0009
% #FIELD_VALUE: 第7项检查结果
% #HANDWRITTEN: 0
\fieldvalue{\handwritten{0}}\\
\hline
8 & \textit{检查内容（扫描文字不清）} &
% #VALUE_ID: LEX-P0056-V0010
% #FIELD_VALUE: 第8项检查结果
% #HANDWRITTEN: 0
\fieldvalue{\handwritten{0}}\\
\hline
9 & \textit{检查内容（扫描文字不清）} &
% #VALUE_ID: LEX-P0056-V0011
% #FIELD_VALUE: 第9项检查结果
% #HANDWRITTEN: 0
\fieldvalue{\handwritten{0}}\\
\hline
10 & \textit{检查内容（扫描文字不清）} &
% #VALUE_ID: LEX-P0056-V0012
% #FIELD_VALUE: 第10项检查结果
% #HANDWRITTEN: 0
\fieldvalue{\handwritten{0}}\\
\hline
11 & \textit{检查内容（扫描文字不清）} &
% #VALUE_ID: LEX-P0056-V0013
% #FIELD_VALUE: 第11项检查结果
% #HANDWRITTEN: 0
\fieldvalue{\handwritten{0}}\\
\hline
12 & \textit{检查内容（扫描文字不清）} &
% #VALUE_ID: LEX-P0056-V0014
% #FIELD_VALUE: 第12项检查结果
% #HANDWRITTEN: 0
\fieldvalue{\handwritten{0}}\\
\hline
13 & \textit{检查内容（扫描文字不清）} &
% #VALUE_ID: LEX-P0056-V0015
% #FIELD_VALUE: 第13项检查结果
% #HANDWRITTEN: 0
\fieldvalue{\handwritten{0}}\\
\hline
14 & \textit{检查内容（扫描文字不清）} &
% #VALUE_ID: LEX-P0056-V0016
% #FIELD_VALUE: 第14项检查结果
% #HANDWRITTEN: 0
\fieldvalue{\handwritten{0}}\\
\hline
15 & \textit{检查内容（扫描文字不清）} &
% #VALUE_ID: LEX-P0056-V0017
% #FIELD_VALUE: 第15项检查结果
% #HANDWRITTEN: 0
\fieldvalue{\handwritten{0}}\\
\hline
16 & \textit{检查内容（扫描文字不清）} &
% #VALUE_ID: LEX-P0056-V0018
% #FIELD_VALUE: 第16项检查结果
% #HANDWRITTEN: 0
\fieldvalue{\handwritten{0}}\\
\hline
17 & \textit{检查内容（扫描文字不清）} &
% #VALUE_ID: LEX-P0056-V0019
% #FIELD_VALUE: 第17项检查结果
% #HANDWRITTEN: 0
\fieldvalue{\handwritten{0}}\\
\hline
18 & \textit{检查内容（扫描文字不清）} &
% #VALUE_ID: LEX-P0056-V0020
% #FIELD_VALUE: 第18项检查结果
% #HANDWRITTEN: 0
\fieldvalue{\handwritten{0}}\\
\hline
19 & \textit{检查内容（扫描文字不清）} &
% #VALUE_ID: LEX-P0056-V0021
% #FIELD_VALUE: 第19项检查结果
% #HANDWRITTEN: 0
\fieldvalue{\handwritten{0}}\\
\hline
20 & \textit{检查内容（扫描文字不清）} &
% #VALUE_ID: LEX-P0056-V0022
% #FIELD_VALUE: 第20项检查结果
% #HANDWRITTEN: 0
\fieldvalue{\handwritten{0}}\\
\hline
21 & \textit{检查内容（扫描文字不清）} &
% #VALUE_ID: LEX-P0056-V0023
% #FIELD_VALUE: 第21项检查结果
% #HANDWRITTEN: 0
\fieldvalue{\handwritten{0}}\\
\hline
22 & \textit{检查内容（扫描文字不清）} &
% #VALUE_ID: LEX-P0056-V0024
% #FIELD_VALUE: 第22项检查结果
% #HANDWRITTEN: 0
\fieldvalue{\handwritten{0}}\\
\hline
23 & \textit{检查内容（扫描文字不清）} &
% #VALUE_ID: LEX-P0056-V0025
% #FIELD_VALUE: 第23项检查结果
% #HANDWRITTEN: 0
\fieldvalue{\handwritten{0}}\\
\hline
24 & \textit{检查内容（扫描文字不清）} &
% #VALUE_ID: LEX-P0056-V0026
% #FIELD_VALUE: 第24项检查结果
% #HANDWRITTEN: 25
\fieldvalue{\handwritten{25}}\\
\hline
25 & \textit{检查内容（扫描文字不清）} &
% #VALUE_ID: LEX-P0056-V0027
% #FIELD_VALUE: 第25项检查结果
% #HANDWRITTEN: 12
\fieldvalue{\handwritten{12}}\\
\hline
\end{tabularx}

% TODO：检查表及其文字内容在下一页继续。
% LEXOID_PAGE_COMPLETED: 56/78

\newpage

% TODO: The source page is rotated and several fixed Japanese labels are difficult to read.
\begin{center}
\textbf{（判読困難な日本語様式）}
\end{center}

\noindent
\begin{tabularx}{\textwidth}{|>{\RaggedRight\arraybackslash}X|>{\RaggedRight\arraybackslash}X|}
\hline
% #VALUE_ID: LEX-P0057-V0001
% #FIELD_VALUE: 日付
% #HANDWRITTEN: 2025.02.01
\fieldvalue{\handwritten{2025.02.01}}
&
% #VALUE_ID: LEX-P0057-V0002
% #FIELD_VALUE: 記入者名・署名
% #HANDWRITTEN: 判読困難
\fieldvalue{\handwritten{判読困難}}
\\
\hline
\multicolumn{2}{|l|}{\textbf{記録欄}}\\
\hline
% #VALUE_ID: LEX-P0057-V0003
% #FIELD_VALUE: 第1項目の数値
% #HANDWRITTEN: 0
\fieldvalue{\handwritten{0}}
&
\text{（印刷された項目名）}\\
\hline
% #VALUE_ID: LEX-P0057-V0004
% #FIELD_VALUE: 第2項目の数値
% #HANDWRITTEN: 0
\fieldvalue{\handwritten{0}}
&
\text{（印刷された項目名）}\\
\hline
% #VALUE_ID: LEX-P0057-V0005
% #FIELD_VALUE: 第3項目の数値
% #HANDWRITTEN: 0
\fieldvalue{\handwritten{0}}
&
\text{（印刷された項目名）}\\
\hline
% #VALUE_ID: LEX-P0057-V0006
% #FIELD_VALUE: 第4項目の数値
% #HANDWRITTEN: 0
\fieldvalue{\handwritten{0}}
&
\text{（印刷された項目名）}\\
\hline
% #VALUE_ID: LEX-P0057-V0007
% #FIELD_VALUE: 第5項目の数値
% #HANDWRITTEN: 0
\fieldvalue{\handwritten{0}}
&
\text{（印刷された項目名）}\\
\hline
% #VALUE_ID: LEX-P0057-V0008
% #FIELD_VALUE: 第6項目の数値
% #HANDWRITTEN: 0
\fieldvalue{\handwritten{0}}
&
\text{（印刷された項目名）}\\
\hline
% #VALUE_ID: LEX-P0057-V0009
% #FIELD_VALUE: 第7項目の数値
% #HANDWRITTEN: 0
\fieldvalue{\handwritten{0}}
&
\text{（印刷された項目名）}\\
\hline
% #VALUE_ID: LEX-P0057-V0010
% #FIELD_VALUE: 第8項目の数値
% #HANDWRITTEN: 0
\fieldvalue{\handwritten{0}}
&
\text{（印刷された項目名）}\\
\hline
% #VALUE_ID: LEX-P0057-V0011
% #FIELD_VALUE: 第9項目の数値
% #HANDWRITTEN: 0
\fieldvalue{\handwritten{0}}
&
\text{（印刷された項目名）}\\
\hline
% #VALUE_ID: LEX-P0057-V0012
% #FIELD_VALUE: 第10項目の数値
% #HANDWRITTEN: 0
\fieldvalue{\handwritten{0}}
&
\text{（印刷された項目名）}\\
\hline
% #VALUE_ID: LEX-P0057-V0013
% #FIELD_VALUE: 第11項目の数値
% #HANDWRITTEN: 01
\fieldvalue{\handwritten{01}}
&
\text{（印刷された項目名）}\\
\hline
% #VALUE_ID: LEX-P0057-V0014
% #FIELD_VALUE: 第12項目の数値
% #HANDWRITTEN: 0
\fieldvalue{\handwritten{0}}
&
\text{（印刷された項目名）}\\
\hline
% #VALUE_ID: LEX-P0057-V0015
% #FIELD_VALUE: 第13項目の数値
% #HANDWRITTEN: 0
\fieldvalue{\handwritten{0}}
&
\text{（印刷された項目名）}\\
\hline
% #VALUE_ID: LEX-P0057-V0016
% #FIELD_VALUE: 第14項目の数値
% #HANDWRITTEN: 0
\fieldvalue{\handwritten{0}}
&
\text{（印刷された項目名）}\\
\hline
% #VALUE_ID: LEX-P0057-V0017
% #FIELD_VALUE: 第15項目の数値
% #HANDWRITTEN: 0
\fieldvalue{\handwritten{0}}
&
\text{（印刷された項目名）}\\
\hline
% #VALUE_ID: LEX-P0057-V0018
% #FIELD_VALUE: 第16項目の数値
% #HANDWRITTEN: 2
\fieldvalue{\handwritten{2}}
&
\text{（印刷された項目名）}\\
\hline
% #VALUE_ID: LEX-P0057-V0019
% #FIELD_VALUE: 第17項目の数値
% #HANDWRITTEN: 80
\fieldvalue{\handwritten{80}}
&
\text{（印刷された項目名）}\\
\hline
\end{tabularx}

\vspace{1em}

\noindent
\text{（下部に印刷された様式情報・注意書き）}

% TODO: The page contains additional fixed Japanese text, stamps, and form headings that are
% rotated and too low-resolution for reliable transcription.
% LEXOID_PAGE_COMPLETED: 57/78

\newpage

% TODO: The source image is rotated 180 degrees and several characters are too low-resolution to transcribe confidently.

\begin{center}
\textbf{检验报告}
\end{center}

\begin{tabularx}{\textwidth}{|p{3.2cm}|X|p{3.2cm}|X|}
\hline
报告编号 &
% #VALUE_ID: LEX-P0058-V0001
% #FIELD_VALUE: 报告编号
% #TODO #HANDWRITTEN: [编号不清]; 蓝色手写内容分辨率不足
\fieldvalue{\handwritten{[编号不清]}} &
检验日期 &
% #VALUE_ID: LEX-P0058-V0002
% #FIELD_VALUE: 检验日期
% #HANDWRITTEN: 8/9/2024
\fieldvalue{\handwritten{8/9/2024}} \\
\hline
受检单位名称 &
% #VALUE_ID: LEX-P0058-V0003
% #FIELD_VALUE: 受检单位名称
% #TODO #HANDWRITTEN: [单位名称不清]; 蓝色手写字迹无法辨认
\fieldvalue{\handwritten{[单位名称不清]}} &
检验类别 &
% #VALUE_ID: LEX-P0058-V0004
% #FIELD_VALUE: 检验类别
% #TODO #HANDWRITTEN: [类别不清]; 蓝色手写内容无法辨认
\fieldvalue{\handwritten{[类别不清]}} \\
\hline
设备名称 &
% #VALUE_ID: LEX-P0058-V0005
% #FIELD_VALUE: 设备名称
% #TODO #HANDWRITTEN: [设备名称不清]; 蓝色手写内容无法辨认
\fieldvalue{\handwritten{[设备名称不清]}} &
规格型号 &
% #VALUE_ID: LEX-P0058-V0006
% #FIELD_VALUE: 规格型号
% #TODO #HANDWRITTEN: 200--480（末位不清）
\fieldvalue{\handwritten{200--480}} \\
\hline
\end{tabularx}

\vspace{0.5em}

\begin{tabularx}{\textwidth}{|p{4.2cm}|X|p{4.2cm}|X|}
\hline
\multicolumn{4}{|c|}{\textbf{检验项目及检验结果}}\\
\hline
\multicolumn{2}{|c|}{检验项目} &
\multicolumn{2}{c|}{检验结果}\\
\hline
资料审查 &
\begin{itemize}
\item 资料齐全
\item 资料不齐全
\end{itemize}
&
检验结论 &
% #VALUE_ID: LEX-P0058-V0007
% #FIELD_VALUE: 检验结论
\fieldvalue{合格} \\
\hline
外观检查 &
% #VALUE_ID: LEX-P0058-V0008
% #FIELD_VALUE: 外观检查
\fieldvalue{符合要求} &
安全装置检查 &
% #VALUE_ID: LEX-P0058-V0009
% #FIELD_VALUE: 安全装置检查
\fieldvalue{符合要求} \\
\hline
运行试验 &
% #VALUE_ID: LEX-P0058-V0010
% #FIELD_VALUE: 运行试验
\fieldvalue{符合要求} &
性能试验 &
% #VALUE_ID: LEX-P0058-V0011
% #FIELD_VALUE: 性能试验
\fieldvalue{符合要求} \\
\hline
\end{tabularx}

\vspace{0.5em}

\begin{tabularx}{\textwidth}{|p{3.0cm}|X|p{3.0cm}|X|}
\hline
\multicolumn{4}{|c|}{检验记录}\\
\hline
检验项目 &
检验内容 &
检验项目 &
检验内容\\
\hline
\textbf{一、资料审查} &
% #VALUE_ID: LEX-P0058-V0012
% #FIELD_VALUE: 资料审查结果
\fieldvalue{资料齐全} &
\textbf{二、外观检查} &
% #VALUE_ID: LEX-P0058-V0013
% #FIELD_VALUE: 外观检查结果
\fieldvalue{符合要求} \\
\hline
\textbf{三、安全装置} &
% #VALUE_ID: LEX-P0058-V0014
% #FIELD_VALUE: 安全装置检查结果
\fieldvalue{符合要求} &
\textbf{四、运行试验} &
% #VALUE_ID: LEX-P0058-V0015
% #FIELD_VALUE: 运行试验结果
\fieldvalue{符合要求} \\
\hline
\textbf{五、性能试验} &
% #VALUE_ID: LEX-P0058-V0016
% #FIELD_VALUE: 性能试验结果
\fieldvalue{符合要求} &
\textbf{六、检验结论} &
% #VALUE_ID: LEX-P0058-V0017
% #FIELD_VALUE: 检验结论
\fieldvalue{合格} \\
\hline
\end{tabularx}

\vspace{0.5em}

\noindent
备注：检验过程中发现的问题及处理意见：
% #VALUE_ID: LEX-P0058-V0018
% #FIELD_VALUE: 问题及处理意见
% #TODO #HANDWRITTEN: [手写内容不清]; 表格内蓝色手写记录无法辨认
\fieldvalue{\handwritten{[手写内容不清]}}。

\vspace{1em}

\begin{tabularx}{\textwidth}{|p{3cm}|X|p{3cm}|X|}
\hline
检验人员 &
% #VALUE_ID: LEX-P0058-V0019
% #FIELD_VALUE: 检验人员
% #TODO #HANDWRITTEN: [姓名不清]; 蓝色签名无法辨认
\fieldvalue{\handwritten{[姓名不清]}} &
审核人员 &
% #VALUE_ID: LEX-P0058-V0020
% #FIELD_VALUE: 审核人员
% #TODO #HANDWRITTEN: [签名不清]; 蓝色手写签名无法辨认
\fieldvalue{\handwritten{[签名不清]}} \\
\hline
检验日期 &
% #VALUE_ID: LEX-P0058-V0021
% #FIELD_VALUE: 检验日期
% #TODO #HANDWRITTEN: [日期不清]; 页眉蓝色手写日期与表内日期均难以辨认
\fieldvalue{\handwritten{[日期不清]}} &
批准日期 &
% #VALUE_ID: LEX-P0058-V0022
% #FIELD_VALUE: 批准日期
% #TODO #HANDWRITTEN: [日期不清]; 蓝色手写日期无法辨认
\fieldvalue{\handwritten{[日期不清]}} \\
\hline
\end{tabularx}

\noindent
\textbf{注：} 本页内容为表格续页。
% TODO: The lower portion of the source page contains additional table rows and handwritten entries that are partly obscured and illegible.
% LEXOID_PAGE_COMPLETED: 58/78

\newpage

\begin{center}
\textbf{（別紙様式22）}\\[2pt]
\textbf{実績報告書}
\end{center}

\begin{tabularx}{\textwidth}{|l|X|r|}
\hline
文書番号 & \centering 令和\underline{\hspace{1.2cm}}年\underline{\hspace{1.2cm}}月\underline{\hspace{1.2cm}}日 &
\begin{minipage}{0.18\textwidth}
% #VALUE_ID: LEX-P0059-V0001
% #FIELD_VALUE: 文書番号・記入欄
% #TODO #HANDWRITTEN: 判読困難; 青色の手書き文字
\fieldvalue{\handwritten{判読困難}}
\end{minipage}
\\
\hline
\multicolumn{3}{|c|}{事業名・団体名等}\\
\hline
\end{tabularx}

\vspace{0.5em}

\begin{tabularx}{\textwidth}{|X|p{0.29\textwidth}|}
\hline
\centering 項目 & \centering 記入欄\arraybackslash\\
\hline
\multicolumn{2}{|l|}{\textbf{①}　事業の実施状況}\\
\hline
実施した事業の内容 & 
% #VALUE_ID: LEX-P0059-V0002
% #FIELD_VALUE: 実施した事業の内容
% #TODO #HANDWRITTEN: 5; 青色の手書き文字
\fieldvalue{\handwritten{5}}\\
\hline
参加人数等 &
% #VALUE_ID: LEX-P0059-V0003
% #FIELD_VALUE: 参加人数等
% #HANDWRITTEN: 17
\fieldvalue{\handwritten{17}}\\
\hline
事業の成果 &
% #VALUE_ID: LEX-P0059-V0004
% #FIELD_VALUE: 事業の成果
% #HANDWRITTEN: 0
\fieldvalue{\handwritten{0}}\\
\hline
今後の課題 &
% #VALUE_ID: LEX-P0059-V0005
% #FIELD_VALUE: 今後の課題
% #HANDWRITTEN: 0
\fieldvalue{\handwritten{0}}\\
\hline
その他 &
% #VALUE_ID: LEX-P0059-V0006
% #FIELD_VALUE: その他
% #HANDWRITTEN: 0
\fieldvalue{\handwritten{0}}\\
\hline
\multicolumn{2}{|l|}{\textbf{②}　事業費の執行状況}\\
\hline
支出項目 1 &
% #VALUE_ID: LEX-P0059-V0007
% #FIELD_VALUE: 支出項目 1
% #HANDWRITTEN: 0
\fieldvalue{\handwritten{0}}\\
\hline
支出項目 2 &
% #VALUE_ID: LEX-P0059-V0008
% #FIELD_VALUE: 支出項目 2
% #HANDWRITTEN: 0
\fieldvalue{\handwritten{0}}\\
\hline
支出項目 3 &
% #VALUE_ID: LEX-P0059-V0009
% #FIELD_VALUE: 支出項目 3
% #HANDWRITTEN: 0
\fieldvalue{\handwritten{0}}\\
\hline
支出項目 4 &
% #VALUE_ID: LEX-P0059-V0010
% #FIELD_VALUE: 支出項目 4
% #HANDWRITTEN: 0
\fieldvalue{\handwritten{0}}\\
\hline
支出項目 5 &
% #VALUE_ID: LEX-P0059-V0011
% #FIELD_VALUE: 支出項目 5
% #HANDWRITTEN: 0
\fieldvalue{\handwritten{0}}\\
\hline
\multicolumn{2}{|l|}{\textbf{③}　収支の状況}\\
\hline
収入の部 1 &
% #VALUE_ID: LEX-P0059-V0012
% #FIELD_VALUE: 収入の部 1
% #HANDWRITTEN: 0
\fieldvalue{\handwritten{0}}\\
\hline
収入の部 2 &
% #VALUE_ID: LEX-P0059-V0013
% #FIELD_VALUE: 収入の部 2
% #HANDWRITTEN: 0
\fieldvalue{\handwritten{0}}\\
\hline
支出の部 1 &
% #VALUE_ID: LEX-P0059-V0014
% #FIELD_VALUE: 支出の部 1
% #HANDWRITTEN: 0
\fieldvalue{\handwritten{0}}\\
\hline
支出の部 2 &
% #VALUE_ID: LEX-P0059-V0015
% #FIELD_VALUE: 支出の部 2
% #HANDWRITTEN: 0
\fieldvalue{\handwritten{0}}\\
\hline
差引残額 &
% #VALUE_ID: LEX-P0059-V0016
% #FIELD_VALUE: 差引残額
% #TODO #HANDWRITTEN: 01; 数字の境界が不鮮明
\fieldvalue{\handwritten{01}}\\
\hline
\multicolumn{2}{|l|}{\textbf{④}　添付資料等}\\
\hline
添付資料 1 &
% #VALUE_ID: LEX-P0059-V0017
% #FIELD_VALUE: 添付資料 1
% #HANDWRITTEN: 0
\fieldvalue{\handwritten{0}}\\
\hline
添付資料 2 &
% #VALUE_ID: LEX-P0059-V0018
% #FIELD_VALUE: 添付資料 2
% #HANDWRITTEN: 0
\fieldvalue{\handwritten{0}}\\
\hline
添付資料 3 &
% #VALUE_ID: LEX-P0059-V0019
% #FIELD_VALUE: 添付資料 3
% #HANDWRITTEN: 0
\fieldvalue{\handwritten{0}}\\
\hline
添付資料 4 &
% #VALUE_ID: LEX-P0059-V0020
% #FIELD_VALUE: 添付資料 4
% #HANDWRITTEN: 0
\fieldvalue{\handwritten{0}}\\
\hline
添付資料 5 &
% #VALUE_ID: LEX-P0059-V0021
% #FIELD_VALUE: 添付資料 5
% #HANDWRITTEN: 0
\fieldvalue{\handwritten{0}}\\
\hline
\multicolumn{2}{|l|}{\textbf{⑤}　確認事項}\\
\hline
確認事項 1 &
% #VALUE_ID: LEX-P0059-V0022
% #FIELD_VALUE: 確認事項 1
% #HANDWRITTEN: 0
\fieldvalue{\handwritten{0}}\\
\hline
確認事項 2 &
% #VALUE_ID: LEX-P0059-V0023
% #FIELD_VALUE: 確認事項 2
% #HANDWRITTEN: 0
\fieldvalue{\handwritten{0}}\\
\hline
\end{tabularx}

% TODO: The page contains additional Japanese checklist text and form labels that are partially illegible in the supplied scan; reproduce the exact wording after higher-resolution inspection.
% LEXOID_PAGE_COMPLETED: 59/78

\newpage

\begin{center}
\textbf{检验批质量验收记录表（续）}
\end{center}

\noindent
\begin{tabularx}{\textwidth}{@{}lXlX@{}}
工程名称：&
% #VALUE_ID: LEX-P0060-V0001
% #FIELD_VALUE: 工程名称
% #TODO #HANDWRITTEN: （难以辨认）; 蓝色手写内容不清晰
\fieldvalue{\handwritten{（难以辨认）}}&
施工单位：&
% #VALUE_ID: LEX-P0060-V0002
% #FIELD_VALUE: 施工单位
% #TODO #HANDWRITTEN: （难以辨认）; 蓝色手写内容不清晰
\fieldvalue{\handwritten{（难以辨认）}}\\
验收部位：&
% #VALUE_ID: LEX-P0060-V0003
% #FIELD_VALUE: 验收部位
% #TODO #HANDWRITTEN: （难以辨认）; 蓝色手写内容不清晰
\fieldvalue{\handwritten{（难以辨认）}}&
日期：&
% #VALUE_ID: LEX-P0060-V0004
% #FIELD_VALUE: 日期
% #HANDWRITTEN: 2015.01.18
\fieldvalue{\handwritten{2015.01.18}}
\end{tabularx}

\vspace{0.5em}

\small
\begin{tabularx}{\textwidth}{|c|X|c|X|}
\hline
序号 & 检验项目 & 验收规范规定 & 施工单位检查评定记录 \\
\hline
\multicolumn{4}{|c|}{主控项目} \\
\hline
1 & 检验项目（本页上部所列） &
规定值 &
% #VALUE_ID: LEX-P0060-V0005
% #FIELD_VALUE: 第1项施工单位检查评定记录
% #HANDWRITTEN: 23--21
\fieldvalue{\handwritten{23--21}} \\
\hline
2 & 检验项目 &
规定值 &
% #VALUE_ID: LEX-P0060-V0006
% #FIELD_VALUE: 第2项施工单位检查评定记录
% #HANDWRITTEN: 2
\fieldvalue{\handwritten{2}} \\
\hline
3 & 检验项目 &
规定值 &
% #VALUE_ID: LEX-P0060-V0007
% #FIELD_VALUE: 第3项施工单位检查评定记录
% #HANDWRITTEN: 0
\fieldvalue{\handwritten{0}} \\
\hline
4 & 检验项目 &
规定值 &
% #VALUE_ID: LEX-P0060-V0008
% #FIELD_VALUE: 第4项施工单位检查评定记录
% #HANDWRITTEN: 0
\fieldvalue{\handwritten{0}} \\
\hline
5 & 检验项目 &
规定值 &
% #VALUE_ID: LEX-P0060-V0009
% #FIELD_VALUE: 第5项施工单位检查评定记录
% #HANDWRITTEN: 0
\fieldvalue{\handwritten{0}} \\
\hline
6 & 检验项目 &
规定值 &
% #VALUE_ID: LEX-P0060-V0010
% #FIELD_VALUE: 第6项施工单位检查评定记录
% #HANDWRITTEN: 0
\fieldvalue{\handwritten{0}} \\
\hline
\multicolumn{4}{|c|}{一般项目} \\
\hline
7 & 检验项目 &
规定值 &
% #VALUE_ID: LEX-P0060-V0011
% #FIELD_VALUE: 第7项施工单位检查评定记录
% #HANDWRITTEN: 0
\fieldvalue{\handwritten{0}} \\
\hline
8 & 检验项目 &
规定值 &
% #VALUE_ID: LEX-P0060-V0012
% #FIELD_VALUE: 第8项施工单位检查评定记录
% #HANDWRITTEN: 0
\fieldvalue{\handwritten{0}} \\
\hline
9 & 检验项目 &
规定值 &
% #VALUE_ID: LEX-P0060-V0013
% #FIELD_VALUE: 第9项施工单位检查评定记录
% #HANDWRITTEN: 0
\fieldvalue{\handwritten{0}} \\
\hline
10 & 检验项目 &
规定值 &
% #VALUE_ID: LEX-P0060-V0014
% #FIELD_VALUE: 第10项施工单位检查评定记录
% #HANDWRITTEN: 0
\fieldvalue{\handwritten{0}} \\
\hline
11 & 检验项目 &
规定值 &
% #VALUE_ID: LEX-P0060-V0015
% #FIELD_VALUE: 第11项施工单位检查评定记录
% #HANDWRITTEN: 0
\fieldvalue{\handwritten{0}} \\
\hline
12 & 检验项目 &
规定值 &
% #VALUE_ID: LEX-P0060-V0016
% #FIELD_VALUE: 第12项施工单位检查评定记录
% #HANDWRITTEN: 0
\fieldvalue{\handwritten{0}} \\
\hline
13 & 检验项目 &
规定值 &
% #VALUE_ID: LEX-P0060-V0017
% #FIELD_VALUE: 第13项施工单位检查评定记录
% #HANDWRITTEN: 0
\fieldvalue{\handwritten{0}} \\
\hline
14 & 检验项目 &
规定值 &
% #VALUE_ID: LEX-P0060-V0018
% #FIELD_VALUE: 第14项施工单位检查评定记录
% #HANDWRITTEN: 0
\fieldvalue{\handwritten{0}} \\
\hline
15 & 检验项目 &
规定值 &
% #VALUE_ID: LEX-P0060-V0019
% #FIELD_VALUE: 第15项施工单位检查评定记录
% #HANDWRITTEN: 0
\fieldvalue{\handwritten{0}} \\
\hline
16 & 检验项目 &
规定值 &
% #VALUE_ID: LEX-P0060-V0020
% #FIELD_VALUE: 第16项施工单位检查评定记录
% #HANDWRITTEN: 0
\fieldvalue{\handwritten{0}} \\
\hline
\end{tabularx}
\normalsize

% TODO: The source page is rotated and several fixed labels and handwritten header entries are illegible. The table above represents only the visible rows on this page; continuation and remaining form fields are on other pages.
% LEXOID_PAGE_COMPLETED: 60/78

\newpage

% TODO: The page is a Japanese administrative form; several printed labels are
% difficult to read because of the scan quality and have been retained as
% approximate transcriptions/placeholders.

\begin{center}
\textbf{令和7年度　（22）調査票}
\end{center}

\noindent
\begin{tabularx}{\textwidth}{|l|X|l|}
\hline
\textbf{区分} & \textbf{項目} & \textbf{記入欄}\\
\hline
令和7年度 & 調査年月日等 & 
% #VALUE_ID: LEX-P0061-V0001
% #FIELD_VALUE: 調査年月日等
% #TODO #HANDWRITTEN: R7.??/??; handwritten date is unclear
\fieldvalue{\handwritten{R7.??/??}}\\
\hline
\end{tabularx}

\vspace{1em}

\noindent
\begin{tabularx}{\textwidth}{|X|p{0.18\textwidth}|}
\hline
\multicolumn{2}{|c|}{\textbf{調査項目}}\\
\cline{1-2}
\textbf{各項目について該当する数値を記入} & \textbf{記入欄}\\
\hline
\multicolumn{2}{|l|}{\textcircled{\scriptsize 1}\quad 調査対象に関する事項}\\
\hline
\text{項目} & 
% #VALUE_ID: LEX-P0061-V0002
% #FIELD_VALUE: ① 記入欄
% #HANDWRITTEN: 有：1
\fieldvalue{\handwritten{有：1}}\\
\hline
\multicolumn{2}{|l|}{\textcircled{\scriptsize 2}\quad 調査項目}\\
\hline
\text{項目1} &
% #VALUE_ID: LEX-P0061-V0003
% #FIELD_VALUE: ② 項目1
% #HANDWRITTEN: 0
\fieldvalue{\handwritten{0}}\\
\text{項目2} &
% #VALUE_ID: LEX-P0061-V0004
% #FIELD_VALUE: ② 項目2
% #HANDWRITTEN: 0
\fieldvalue{\handwritten{0}}\\
\text{項目3} &
% #VALUE_ID: LEX-P0061-V0005
% #FIELD_VALUE: ② 項目3
% #HANDWRITTEN: 0
\fieldvalue{\handwritten{0}}\\
\text{項目4} &
% #VALUE_ID: LEX-P0061-V0006
% #FIELD_VALUE: ② 項目4
% #HANDWRITTEN: 0
\fieldvalue{\handwritten{0}}\\
\text{項目5} &
% #VALUE_ID: LEX-P0061-V0007
% #FIELD_VALUE: ② 項目5
% #HANDWRITTEN: 01
\fieldvalue{\handwritten{01}}\\
\hline
\multicolumn{2}{|l|}{\textcircled{\scriptsize 3}\quad 調査項目}\\
\hline
\text{項目1} &
% #VALUE_ID: LEX-P0061-V0008
% #FIELD_VALUE: ③ 項目1
% #HANDWRITTEN: 0
\fieldvalue{\handwritten{0}}\\
\text{項目2} &
% #VALUE_ID: LEX-P0061-V0009
% #FIELD_VALUE: ③ 項目2
% #HANDWRITTEN: 0
\fieldvalue{\handwritten{0}}\\
\text{項目3} &
% #VALUE_ID: LEX-P0061-V0010
% #FIELD_VALUE: ③ 項目3
% #HANDWRITTEN: 0
\fieldvalue{\handwritten{0}}\\
\text{項目4} &
% #VALUE_ID: LEX-P0061-V0011
% #FIELD_VALUE: ③ 項目4
% #HANDWRITTEN: 0
\fieldvalue{\handwritten{0}}\\
\text{項目5} &
% #VALUE_ID: LEX-P0061-V0012
% #FIELD_VALUE: ③ 項目5
% #HANDWRITTEN: 01
\fieldvalue{\handwritten{01}}\\
\hline
\multicolumn{2}{|l|}{\textcircled{\scriptsize 4}\quad 調査項目}\\
\hline
\text{項目1} &
% #VALUE_ID: LEX-P0061-V0013
% #FIELD_VALUE: ④ 項目1
% #HANDWRITTEN: 0
\fieldvalue{\handwritten{0}}\\
\text{項目2} &
% #VALUE_ID: LEX-P0061-V0014
% #FIELD_VALUE: ④ 項目2
% #HANDWRITTEN: 0
\fieldvalue{\handwritten{0}}\\
\text{項目3} &
% #VALUE_ID: LEX-P0061-V0015
% #FIELD_VALUE: ④ 項目3
% #HANDWRITTEN: 0
\fieldvalue{\handwritten{0}}\\
\text{項目4} &
% #VALUE_ID: LEX-P0061-V0016
% #FIELD_VALUE: ④ 項目4
% #HANDWRITTEN: 0
\fieldvalue{\handwritten{0}}\\
\text{項目5} &
% #VALUE_ID: LEX-P0061-V0017
% #FIELD_VALUE: ④ 項目5
% #HANDWRITTEN: 01
\fieldvalue{\handwritten{01}}\\
\hline
\multicolumn{2}{|l|}{\textcircled{\scriptsize 5}\quad 調査項目}\\
\hline
\text{項目1} &
% #VALUE_ID: LEX-P0061-V0018
% #FIELD_VALUE: ⑤ 項目1
% #HANDWRITTEN: 0
\fieldvalue{\handwritten{0}}\\
\text{項目2} &
% #VALUE_ID: LEX-P0061-V0019
% #FIELD_VALUE: ⑤ 項目2
% #HANDWRITTEN: 0
\fieldvalue{\handwritten{0}}\\
\text{項目3} &
% #VALUE_ID: LEX-P0061-V0020
% #FIELD_VALUE: ⑤ 項目3
% #HANDWRITTEN: 0
\fieldvalue{\handwritten{0}}\\
\text{項目4} &
% #VALUE_ID: LEX-P0061-V0021
% #FIELD_VALUE: ⑤ 項目4
% #HANDWRITTEN: 0
\fieldvalue{\handwritten{0}}\\
\text{項目5} &
% #VALUE_ID: LEX-P0061-V0022
% #FIELD_VALUE: ⑤ 項目5
% #HANDWRITTEN: 01
\fieldvalue{\handwritten{01}}\\
\hline
\multicolumn{2}{|l|}{\textcircled{\scriptsize 6}\quad 調査項目}\\
\hline
\text{項目1} &
% #VALUE_ID: LEX-P0061-V0023
% #FIELD_VALUE: ⑥ 項目1
% #HANDWRITTEN: 0
\fieldvalue{\handwritten{0}}\\
\text{項目2} &
% #VALUE_ID: LEX-P0061-V0024
% #FIELD_VALUE: ⑥ 項目2
% #HANDWRITTEN: 0
\fieldvalue{\handwritten{0}}\\
\text{項目3} &
% #VALUE_ID: LEX-P0061-V0025
% #FIELD_VALUE: ⑥ 項目3
% #HANDWRITTEN: 0
\fieldvalue{\handwritten{0}}\\
\text{項目4} &
% #VALUE_ID: LEX-P0061-V0026
% #FIELD_VALUE: ⑥ 項目4
% #HANDWRITTEN: 0
\fieldvalue{\handwritten{0}}\\
\text{項目5} &
% #VALUE_ID: LEX-P0061-V0027
% #FIELD_VALUE: ⑥ 項目5
% #HANDWRITTEN: 01
\fieldvalue{\handwritten{01}}\\
\hline
\end{tabularx}

% TODO: The original page contains additional Japanese explanatory text in a
% large adjacent panel; the text is not reliably legible in the supplied scan.
% LEXOID_PAGE_COMPLETED: 61/78

\newpage

\section*{医院感染管理检查记录表}

\begin{tabularx}{\textwidth}{|p{0.18\textwidth}|X|}
\hline
备注 &
% #VALUE_ID: LEX-P0062-V0001
% #FIELD_VALUE: 备注
% #HANDWRITTEN: /
\fieldvalue{\handwritten{/}} \\ \hline
\end{tabularx}

\vspace{0.5em}

% TODO: The checklist text is partially unclear in the scan; only the visible rows and handwritten entries are reproduced.
\begin{tabularx}{\textwidth}{|X|p{0.16\textwidth}|}
\hline
\multicolumn{2}{|c|}{\textbf{检查项目及结果}} \\ \hline

检查项目（可见行1） &
% #VALUE_ID: LEX-P0062-V0002
% #FIELD_VALUE: 检查结果
% #HANDWRITTEN: 0
\fieldvalue{\handwritten{0}} \\ \hline

检查项目（可见行2） &
% #VALUE_ID: LEX-P0062-V0003
% #FIELD_VALUE: 检查结果
% #HANDWRITTEN: 0
\fieldvalue{\handwritten{0}} \\ \hline

检查项目（可见行3） &
% #VALUE_ID: LEX-P0062-V0004
% #FIELD_VALUE: 检查结果
% #HANDWRITTEN: 0
\fieldvalue{\handwritten{0}} \\ \hline

检查项目（可见行4） &
% #VALUE_ID: LEX-P0062-V0005
% #FIELD_VALUE: 检查结果
% #HANDWRITTEN: 0
\fieldvalue{\handwritten{0}} \\ \hline

检查项目（可见行5） &
% #VALUE_ID: LEX-P0062-V0006
% #FIELD_VALUE: 检查结果
% #HANDWRITTEN: 0
\fieldvalue{\handwritten{0}} \\ \hline

检查项目（可见行6） &
% #VALUE_ID: LEX-P0062-V0007
% #FIELD_VALUE: 检查结果
% #HANDWRITTEN: 0
\fieldvalue{\handwritten{0}} \\ \hline

检查项目（可见行7） &
% #VALUE_ID: LEX-P0062-V0008
% #FIELD_VALUE: 检查结果
% #HANDWRITTEN: 0
\fieldvalue{\handwritten{0}} \\ \hline

检查项目（可见行8） &
% #VALUE_ID: LEX-P0062-V0009
% #FIELD_VALUE: 检查结果
% #HANDWRITTEN: 0
\fieldvalue{\handwritten{0}} \\ \hline

检查项目（可见行9） &
% #VALUE_ID: LEX-P0062-V0010
% #FIELD_VALUE: 检查结果
% #HANDWRITTEN: 0
\fieldvalue{\handwritten{0}} \\ \hline

检查项目（可见行10） &
% #VALUE_ID: LEX-P0062-V0011
% #FIELD_VALUE: 检查结果
% #HANDWRITTEN: 0
\fieldvalue{\handwritten{0}} \\ \hline

检查项目（可见行11） &
% #VALUE_ID: LEX-P0062-V0012
% #FIELD_VALUE: 检查结果
% #HANDWRITTEN: 0
\fieldvalue{\handwritten{0}} \\ \hline

检查项目（可见行12） &
% #VALUE_ID: LEX-P0062-V0013
% #FIELD_VALUE: 检查结果
% #HANDWRITTEN: 10
\fieldvalue{\handwritten{10}} \\ \hline

检查项目（可见行13） &
% #VALUE_ID: LEX-P0062-V0014
% #FIELD_VALUE: 检查结果
% #HANDWRITTEN: 0
\fieldvalue{\handwritten{0}} \\ \hline

检查项目（可见行14） &
% #VALUE_ID: LEX-P0062-V0015
% #FIELD_VALUE: 检查结果
% #HANDWRITTEN: 0
\fieldvalue{\handwritten{0}} \\ \hline

检查项目（可见行15） &
% #VALUE_ID: LEX-P0062-V0016
% #FIELD_VALUE: 检查结果
% #HANDWRITTEN: 0
\fieldvalue{\handwritten{0}} \\ \hline

检查项目（可见行16） &
% #VALUE_ID: LEX-P0062-V0017
% #FIELD_VALUE: 检查结果
% #HANDWRITTEN: 0
\fieldvalue{\handwritten{0}} \\ \hline

检查项目（可见行17） &
% #VALUE_ID: LEX-P0062-V0018
% #FIELD_VALUE: 检查结果
% #HANDWRITTEN: 2
\fieldvalue{\handwritten{2}} \\ \hline
\end{tabularx}

\vspace{0.75em}

\begin{tabularx}{\textwidth}{|p{0.22\textwidth}|X|p{0.22\textwidth}|X|}
\hline
检查日期 &
% #VALUE_ID: LEX-P0062-V0019
% #FIELD_VALUE: 检查日期
% #HANDWRITTEN: 2025.06.18
\fieldvalue{\handwritten{2025.06.18}} &
复核日期 &
% #VALUE_ID: LEX-P0062-V0020
% #FIELD_VALUE: 复核日期
% #HANDWRITTEN: 2025.06.18
\fieldvalue{\handwritten{2025.06.18}} \\ \hline
检查人 &
% #VALUE_ID: LEX-P0062-V0021
% #FIELD_VALUE: 检查人
% TODO: handwritten name is unclear in the scan
% #TODO #HANDWRITTEN: 王??; handwritten name is partially illegible
\fieldvalue{\handwritten{王??}} &
\multicolumn{2}{c|}{} \\ \hline
\end{tabularx}

% TODO: The form continues beyond the visible page boundary.
% LEXOID_PAGE_COMPLETED: 62/78

\newpage

% TODO: The page is rendered upside down in the source image; the following structure reflects the corrected reading orientation.
\begin{center}
\textbf{审查项目及审查结论}
\end{center}

\begin{tabularx}{\textwidth}{|p{0.18\textwidth}|X|p{0.18\textwidth}|}
\hline
\textbf{序号} & \textbf{审查项目} & \textbf{审查结论} \\
\hline
① & 审查结论 & 
% #VALUE_ID: LEX-P0063-V0001
% #FIELD_VALUE: 审查结论（左侧数值）
% #TODO #HANDWRITTEN: 1; handwritten mark is unclear
\fieldvalue{\handwritten{1}}
\quad
% #VALUE_ID: LEX-P0063-V0002
% #FIELD_VALUE: 审查结论（右侧数值）
% #TODO #HANDWRITTEN: 45; handwritten mark is unclear
\fieldvalue{\handwritten{45}}
\\
\hline
\multicolumn{3}{|l|}{② \quad 审查项目} \\
\hline
 & 项目 1 &
% #VALUE_ID: LEX-P0063-V0003
% #FIELD_VALUE: ②项目1审查值
% #HANDWRITTEN: 0
\fieldvalue{\handwritten{0}}
\\
\cline{2-3}
 & 项目 2 &
% #VALUE_ID: LEX-P0063-V0004
% #FIELD_VALUE: ②项目2审查值
% #HANDWRITTEN: 0
\fieldvalue{\handwritten{0}}
\\
\cline{2-3}
 & 项目 3 &
% #VALUE_ID: LEX-P0063-V0005
% #FIELD_VALUE: ②项目3审查值
% #HANDWRITTEN: 0
\fieldvalue{\handwritten{0}}
\\
\cline{2-3}
 & 项目 4 &
% #VALUE_ID: LEX-P0063-V0006
% #FIELD_VALUE: ②项目4审查值
% #HANDWRITTEN: 0
\fieldvalue{\handwritten{0}}
\\
\cline{2-3}
 & 项目 5 &
% #VALUE_ID: LEX-P0063-V0007
% #FIELD_VALUE: ②项目5审查值
% #HANDWRITTEN: 0
\fieldvalue{\handwritten{0}}
\\
\cline{2-3}
 & 项目 6 &
% #VALUE_ID: LEX-P0063-V0008
% #FIELD_VALUE: ②项目6审查值
% #HANDWRITTEN: 0
\fieldvalue{\handwritten{0}}
\\
\hline
\multicolumn{3}{|l|}{③ \quad 审查项目} \\
\hline
 & 项目 1 &
% #VALUE_ID: LEX-P0063-V0009
% #FIELD_VALUE: ③项目1审查值
% #HANDWRITTEN: 0
\fieldvalue{\handwritten{0}}
\\
\cline{2-3}
 & 项目 2 &
% #VALUE_ID: LEX-P0063-V0010
% #FIELD_VALUE: ③项目2审查值
% #HANDWRITTEN: 0
\fieldvalue{\handwritten{0}}
\\
\cline{2-3}
 & 项目 3 &
% #VALUE_ID: LEX-P0063-V0011
% #FIELD_VALUE: ③项目3审查值
% #HANDWRITTEN: 0
\fieldvalue{\handwritten{0}}
\\
\cline{2-3}
 & 项目 4 &
% #VALUE_ID: LEX-P0063-V0012
% #FIELD_VALUE: ③项目4审查值
% #HANDWRITTEN: 0
\fieldvalue{\handwritten{0}}
\\
\cline{2-3}
 & 项目 5 &
% #VALUE_ID: LEX-P0063-V0013
% #FIELD_VALUE: ③项目5审查值
% #HANDWRITTEN: 0
\fieldvalue{\handwritten{0}}
\\
\cline{2-3}
 & 项目 6 &
% #VALUE_ID: LEX-P0063-V0014
% #FIELD_VALUE: ③项目6审查值
% #HANDWRITTEN: 0
\fieldvalue{\handwritten{0}}
\\
\hline
\multicolumn{3}{|l|}{④ \quad 审查项目} \\
\hline
 & 项目 1 &
% #VALUE_ID: LEX-P0063-V0015
% #FIELD_VALUE: ④项目1审查值
% #HANDWRITTEN: 0
\fieldvalue{\handwritten{0}}
\\
\cline{2-3}
 & 项目 2 &
% #VALUE_ID: LEX-P0063-V0016
% #FIELD_VALUE: ④项目2审查值
% #HANDWRITTEN: 0
\fieldvalue{\handwritten{0}}
\\
\cline{2-3}
 & 项目 3 &
% #VALUE_ID: LEX-P0063-V0017
% #FIELD_VALUE: ④项目3审查值
% #HANDWRITTEN: 0
\fieldvalue{\handwritten{0}}
\\
\cline{2-3}
 & 项目 4 &
% #VALUE_ID: LEX-P0063-V0018
% #FIELD_VALUE: ④项目4审查值
% #HANDWRITTEN: 0
\fieldvalue{\handwritten{0}}
\\
\cline{2-3}
 & 项目 5 &
% #VALUE_ID: LEX-P0063-V0019
% #FIELD_VALUE: ④项目5审查值
% #HANDWRITTEN: 0
\fieldvalue{\handwritten{0}}
\\
\hline
\multicolumn{3}{|l|}{⑤ \quad 审查项目} \\
\hline
 & 项目 1 &
% #VALUE_ID: LEX-P0063-V0020
% #FIELD_VALUE: ⑤项目1审查值
% #HANDWRITTEN: 0
\fieldvalue{\handwritten{0}}
\\
\cline{2-3}
 & 项目 2 &
% #VALUE_ID: LEX-P0063-V0021
% #FIELD_VALUE: ⑤项目2审查值
% #HANDWRITTEN: 0
\fieldvalue{\handwritten{0}}
\\
\cline{2-3}
 & 项目 3 &
% #VALUE_ID: LEX-P0063-V0022
% #FIELD_VALUE: ⑤项目3审查值
% #HANDWRITTEN: 0
\fieldvalue{\handwritten{0}}
\\
\cline{2-3}
 & 项目 4 &
% #VALUE_ID: LEX-P0063-V0023
% #FIELD_VALUE: ⑤项目4审查值
% #HANDWRITTEN: 0
\fieldvalue{\handwritten{0}}
\\
\cline{2-3}
 & 项目 5 &
% #VALUE_ID: LEX-P0063-V0024
% #FIELD_VALUE: ⑤项目5审查值
% #HANDWRITTEN: 0
\fieldvalue{\handwritten{0}}
\\
\hline
\multicolumn{3}{|l|}{⑥ \quad 审查项目} \\
\hline
 & 项目 1 &
% #VALUE_ID: LEX-P0063-V0025
% #FIELD_VALUE: ⑥项目1审查值
% #HANDWRITTEN: 0
\fieldvalue{\handwritten{0}}
\\
\cline{2-3}
 & 项目 2 &
% #VALUE_ID: LEX-P0063-V0026
% #FIELD_VALUE: ⑥项目2审查值
% #HANDWRITTEN: 0
\fieldvalue{\handwritten{0}}
\\
\cline{2-3}
 & 项目 3 &
% #VALUE_ID: LEX-P0063-V0027
% #FIELD_VALUE: ⑥项目3审查值
% #HANDWRITTEN: 0
\fieldvalue{\handwritten{0}}
\\
\cline{2-3}
 & 项目 4 &
% #VALUE_ID: LEX-P0063-V0028
% #FIELD_VALUE: ⑥项目4审查值
% #HANDWRITTEN: 0
\fieldvalue{\handwritten{0}}
\\
\cline{2-3}
 & 项目 5 &
% #VALUE_ID: LEX-P0063-V0029
% #FIELD_VALUE: ⑥项目5审查值
% #HANDWRITTEN: 0
\fieldvalue{\handwritten{0}}
\\
\hline
\end{tabularx}

\vspace{1em}

% TODO: Additional fixed form-header text, document identifiers, seals, and explanatory paragraphs are present on the page but are not legible at the supplied resolution.
% LEXOID_PAGE_COMPLETED: 63/78

\newpage

\begin{center}
\textbf{安全衛生点検表（22）}
\end{center}

\noindent
\begin{tabularx}{\textwidth}{|l|X|l|}
\hline
工事名 & \multicolumn{2}{l|}{\textit{（記入欄）}}\\
\hline
点検日 &
% #VALUE_ID: LEX-P0064-V0001
% #FIELD_VALUE: 点検日
% #HANDWRITTEN: 8/19
\fieldvalue{\handwritten{8/19}} &
\multicolumn{1}{l|}{\textit{年　月　日}}\\
\hline
点検者 &
% #VALUE_ID: LEX-P0064-V0002
% #FIELD_VALUE: 点検者
% #HANDWRITTEN: （判読困難）
\fieldvalue{\handwritten{（判読困難）}} &
\\
\hline
\end{tabularx}

\vspace{0.5em}

\noindent
\begin{tabularx}{\textwidth}{|c|X|c|}
\hline
番号 & 点検項目 & 判定\\
\hline
\multicolumn{3}{|l|}{\textbf{安全設備等}}\\
\hline
1 & 墜落防止設備等は適切に設置されているか。 &
% #VALUE_ID: LEX-P0064-V0003
% #FIELD_VALUE: 1番判定
% #HANDWRITTEN: 1
\fieldvalue{\handwritten{1}}\\
\hline
2 & 作業場所の整理整頓は行われているか。 &
% #VALUE_ID: LEX-P0064-V0004
% #FIELD_VALUE: 2番判定
% #HANDWRITTEN: 0
\fieldvalue{\handwritten{0}}\\
\hline
3 & 通路及び作業床は安全に確保されているか。 &
% #VALUE_ID: LEX-P0064-V0005
% #FIELD_VALUE: 3番判定
% #HANDWRITTEN: 0
\fieldvalue{\handwritten{0}}\\
\hline
4 & 開口部等には必要な措置が講じられているか。 &
% #VALUE_ID: LEX-P0064-V0006
% #FIELD_VALUE: 4番判定
% #HANDWRITTEN: 0
\fieldvalue{\handwritten{0}}\\
\hline
5 & 足場及び作業台は安全な状態か。 &
% #VALUE_ID: LEX-P0064-V0007
% #FIELD_VALUE: 5番判定
% #HANDWRITTEN: 0
\fieldvalue{\handwritten{0}}\\
\hline
6 & はしご等は適正に使用されているか。 &
% #VALUE_ID: LEX-P0064-V0008
% #FIELD_VALUE: 6番判定
% #HANDWRITTEN: 0
\fieldvalue{\handwritten{0}}\\
\hline
7 & 作業員は保護具を適切に着用しているか。 &
% #VALUE_ID: LEX-P0064-V0009
% #FIELD_VALUE: 7番判定
% #HANDWRITTEN: 0
\fieldvalue{\handwritten{0}}\\
\hline
\multicolumn{3}{|l|}{\textbf{作業状況}}\\
\hline
8 & 作業手順及び作業方法は適切か。 &
% #VALUE_ID: LEX-P0064-V0010
% #FIELD_VALUE: 8番判定
% #HANDWRITTEN: 0
\fieldvalue{\handwritten{0}}\\
\hline
9 & 機械・器具等の点検は実施されているか。 &
% #VALUE_ID: LEX-P0064-V0011
% #FIELD_VALUE: 9番判定
% #HANDWRITTEN: 0
\fieldvalue{\handwritten{0}}\\
\hline
10 & 電気設備及び配線は安全な状態か。 &
% #VALUE_ID: LEX-P0064-V0012
% #FIELD_VALUE: 10番判定
% #HANDWRITTEN: 0
\fieldvalue{\handwritten{0}}\\
\hline
11 & 火気の取扱い及び防火措置は適切か。 &
% #VALUE_ID: LEX-P0064-V0013
% #FIELD_VALUE: 11番判定
% #HANDWRITTEN: 0
\fieldvalue{\handwritten{0}}\\
\hline
12 & 危険物等は適切に保管されているか。 &
% #VALUE_ID: LEX-P0064-V0014
% #FIELD_VALUE: 12番判定
% #HANDWRITTEN: 0
\fieldvalue{\handwritten{0}}\\
\hline
13 & 作業区域の立入禁止措置は適切か。 &
% #VALUE_ID: LEX-P0064-V0015
% #FIELD_VALUE: 13番判定
% #HANDWRITTEN: 0
\fieldvalue{\handwritten{0}}\\
\hline
\multicolumn{3}{|l|}{\textbf{その他}}\\
\hline
14 & 標識及び表示は適切に設置されているか。 &
% #VALUE_ID: LEX-P0064-V0016
% #FIELD_VALUE: 14番判定
% #HANDWRITTEN: 0
\fieldvalue{\handwritten{0}}\\
\hline
15 & 作業員への安全指示は行われているか。 &
% #VALUE_ID: LEX-P0064-V0017
% #FIELD_VALUE: 15番判定
% #HANDWRITTEN: 0
\fieldvalue{\handwritten{0}}\\
\hline
16 & その他の安全衛生上の問題はないか。 &
% #VALUE_ID: LEX-P0064-V0018
% #FIELD_VALUE: 16番判定
% #HANDWRITTEN: 0
\fieldvalue{\handwritten{0}}\\
\hline
17 & 総合的に安全が確保されているか。 &
% #VALUE_ID: LEX-P0064-V0019
% #FIELD_VALUE: 17番判定
% #HANDWRITTEN: 0
\fieldvalue{\handwritten{0}}\\
\hline
\end{tabularx}

% TODO: The source page is rotated and several Japanese labels and handwritten entries are low-resolution; verify the title, field labels, date, name, and individual checklist wording against the original PDF.
% LEXOID_PAGE_COMPLETED: 64/78

\newpage

\section*{附件二}

\begin{center}
\textbf{来访登记表（2）}
\end{center}

% TODO: The header fields and several printed labels are low-resolution and partially illegible.

\noindent
日期：
% #VALUE_ID: LEX-P0065-V0001
% #FIELD_VALUE: 日期
% #TODO #HANDWRITTEN: 2024年……；字迹及日期数字不清
\fieldvalue{\handwritten{2024年……}}

\vspace{0.5em}

\begin{tabularx}{\textwidth}{|p{0.23\textwidth}|X|}
\hline
\textbf{项目} & \textbf{内容} \\
\hline
\multicolumn{2}{|l|}{\textbf{检查记录}} \\
\hline
\multicolumn{2}{|l|}{\textcircled{9}\ 检查项目} \\
\hline
项目9--1 &
% #VALUE_ID: LEX-P0065-V0002
% #FIELD_VALUE: 项目9--1记录值
% #HANDWRITTEN: 0
\fieldvalue{\handwritten{0}} \\
\hline
项目9--2 &
% #VALUE_ID: LEX-P0065-V0003
% #FIELD_VALUE: 项目9--2记录值
% #HANDWRITTEN: 0
\fieldvalue{\handwritten{0}} \\
\hline
项目9--3 &
% #VALUE_ID: LEX-P0065-V0004
% #FIELD_VALUE: 项目9--3记录值
% #HANDWRITTEN: 0
\fieldvalue{\handwritten{0}} \\
\hline
项目9--4 &
% #VALUE_ID: LEX-P0065-V0005
% #FIELD_VALUE: 项目9--4记录值
% #HANDWRITTEN: 0
\fieldvalue{\handwritten{0}} \\
\hline
项目9--5 &
% #VALUE_ID: LEX-P0065-V0006
% #FIELD_VALUE: 项目9--5记录值
% #HANDWRITTEN: 1
\fieldvalue{\handwritten{1}} \\
\hline
\multicolumn{2}{|l|}{\textcircled{8}\ 检查项目} \\
\hline
项目8--1 &
% #VALUE_ID: LEX-P0065-V0007
% #FIELD_VALUE: 项目8--1记录值
% #HANDWRITTEN: 0
\fieldvalue{\handwritten{0}} \\
\hline
项目8--2 &
% #VALUE_ID: LEX-P0065-V0008
% #FIELD_VALUE: 项目8--2记录值
% #HANDWRITTEN: 0
\fieldvalue{\handwritten{0}} \\
\hline
项目8--3 &
% #VALUE_ID: LEX-P0065-V0009
% #FIELD_VALUE: 项目8--3记录值
% #HANDWRITTEN: 0
\fieldvalue{\handwritten{0}} \\
\hline
项目8--4 &
% #VALUE_ID: LEX-P0065-V0010
% #FIELD_VALUE: 项目8--4记录值
% #HANDWRITTEN: 0
\fieldvalue{\handwritten{0}} \\
\hline
项目8--5 &
% #VALUE_ID: LEX-P0065-V0011
% #FIELD_VALUE: 项目8--5记录值
% #HANDWRITTEN: 0
\fieldvalue{\handwritten{0}} \\
\hline
\multicolumn{2}{|l|}{\textcircled{7}\ 检查项目} \\
\hline
项目7--1 &
% #VALUE_ID: LEX-P0065-V0012
% #FIELD_VALUE: 项目7--1记录值
% #HANDWRITTEN: 0
\fieldvalue{\handwritten{0}} \\
\hline
项目7--2 &
% #VALUE_ID: LEX-P0065-V0013
% #FIELD_VALUE: 项目7--2记录值
% #HANDWRITTEN: 0
\fieldvalue{\handwritten{0}} \\
\hline
项目7--3 &
% #VALUE_ID: LEX-P0065-V0014
% #FIELD_VALUE: 项目7--3记录值
% #HANDWRITTEN: 0
\fieldvalue{\handwritten{0}} \\
\hline
项目7--4 &
% #VALUE_ID: LEX-P0065-V0015
% #FIELD_VALUE: 项目7--4记录值
% #HANDWRITTEN: 0
\fieldvalue{\handwritten{0}} \\
\hline
项目7--5 &
% #VALUE_ID: LEX-P0065-V0016
% #FIELD_VALUE: 项目7--5记录值
% #HANDWRITTEN: 0
\fieldvalue{\handwritten{0}} \\
\hline
\multicolumn{2}{|l|}{\textcircled{6}\ 检查项目} \\
\hline
项目6--1 &
% #VALUE_ID: LEX-P0065-V0017
% #FIELD_VALUE: 项目6--1记录值
% #HANDWRITTEN: 0
\fieldvalue{\handwritten{0}} \\
\hline
项目6--2 &
% #VALUE_ID: LEX-P0065-V0018
% #FIELD_VALUE: 项目6--2记录值
% #HANDWRITTEN: 0
\fieldvalue{\handwritten{0}} \\
\hline
项目6--3 &
% #VALUE_ID: LEX-P0065-V0019
% #FIELD_VALUE: 项目6--3记录值
% #HANDWRITTEN: 0
\fieldvalue{\handwritten{0}} \\
\hline
项目6--4 &
% #VALUE_ID: LEX-P0065-V0020
% #FIELD_VALUE: 项目6--4记录值
% #HANDWRITTEN: 0
\fieldvalue{\handwritten{0}} \\
\hline
项目6--5 &
% #VALUE_ID: LEX-P0065-V0021
% #FIELD_VALUE: 项目6--5记录值
% #HANDWRITTEN: 0
\fieldvalue{\handwritten{0}} \\
\hline
项目6--6 &
% #VALUE_ID: LEX-P0065-V0022
% #FIELD_VALUE: 项目6--6记录值
% #HANDWRITTEN: 0
\fieldvalue{\handwritten{0}} \\
\hline
项目6--7 &
% #VALUE_ID: LEX-P0065-V0023
% #FIELD_VALUE: 项目6--7记录值
% #HANDWRITTEN: 0
\fieldvalue{\handwritten{0}} \\
\hline
\textbf{合计} &
% #VALUE_ID: LEX-P0065-V0024
% #FIELD_VALUE: 合计
% #HANDWRITTEN: 15
\fieldvalue{\handwritten{15}} \\
\hline
\textbf{复核} &
% #VALUE_ID: LEX-P0065-V0025
% #FIELD_VALUE: 复核
% #HANDWRITTEN: 12
\fieldvalue{\handwritten{12}} \\
\hline
\end{tabularx}

% TODO: The table continues or contains additional printed explanatory text outside the legible crop.
% LEXOID_PAGE_COMPLETED: 65/78

\newpage

\section*{（2）校外実習実施報告書}

\noindent
\begin{tabularx}{\textwidth}{|l|X|l|X|}
\hline
学籍番号 &
% #VALUE_ID: LEX-P0066-V0001
% #FIELD_VALUE: 学籍番号
\fieldvalue{200-0480} &
実施日 &
% #VALUE_ID: LEX-P0066-V0002
% #FIELD_VALUE: 実施日
\fieldvalue{2025/10/01} \\
\hline
氏名 & & 学科・学年 & \\
\hline
\end{tabularx}

\vspace{0.5em}

\begin{center}
\textbf{実習実施報告書（2）}
\end{center}

\noindent
\begin{tabularx}{\textwidth}{|l|X|}
\hline
報告日 &
% #VALUE_ID: LEX-P0066-V0003
% #FIELD_VALUE: 報告日
% #HANDWRITTEN: 8/10/25
\fieldvalue{\handwritten{8/10/25}} \\
\hline
担当者・確認者 &
% #VALUE_ID: LEX-P0066-V0004
% #FIELD_VALUE: 担当者・確認者
% #HANDWRITTEN: 判読困難な署名
\fieldvalue{\handwritten{判読困難な署名}} \\
\hline
備考 &
% #VALUE_ID: LEX-P0066-V0005
% #FIELD_VALUE: 備考
% #HANDWRITTEN: /
\fieldvalue{\handwritten{/}} \\
\hline
\end{tabularx}

\vspace{0.5em}

\noindent
\begin{tabularx}{\textwidth}{|c|X|}
\hline
\multicolumn{2}{|c|}{実施内容・確認事項} \\
\hline
① &
% #VALUE_ID: LEX-P0066-V0006
% #FIELD_VALUE: 実施内容・確認事項①
% #HANDWRITTEN: 判読困難な記入
\fieldvalue{\handwritten{判読困難な記入}} \\
\hline
② &
% #VALUE_ID: LEX-P0066-V0007
% #FIELD_VALUE: 実施内容・確認事項②
% #HANDWRITTEN: ✓
\fieldvalue{\handwritten{✓}} \\
\hline
③ &
% #VALUE_ID: LEX-P0066-V0008
% #FIELD_VALUE: 実施内容・確認事項③
% #HANDWRITTEN: /
\fieldvalue{\handwritten{/}} \\
\hline
④ &
% #VALUE_ID: LEX-P0066-V0009
% #FIELD_VALUE: 実施内容・確認事項④
% #HANDWRITTEN: ✓
\fieldvalue{\handwritten{✓}} \\
\hline
⑤ &
% #VALUE_ID: LEX-P0066-V0010
% #FIELD_VALUE: 実施内容・確認事項⑤
% #HANDWRITTEN: /
\fieldvalue{\handwritten{/}} \\
\hline
⑥ &
% #VALUE_ID: LEX-P0066-V0011
% #FIELD_VALUE: 実施内容・確認事項⑥
% #HANDWRITTEN: /
\fieldvalue{\handwritten{/}} \\
\hline
⑦ &
% #VALUE_ID: LEX-P0066-V0012
% #FIELD_VALUE: 実施内容・確認事項⑦
% #HANDWRITTEN: ✓
\fieldvalue{\handwritten{✓}} \\
\hline
⑧ &
% #VALUE_ID: LEX-P0066-V0013
% #FIELD_VALUE: 実施内容・確認事項⑧
% #HANDWRITTEN: 判読困難な文字列
\fieldvalue{\handwritten{判読困難な文字列}} \\
\hline
⑨ &
% #VALUE_ID: LEX-P0066-V0014
% #FIELD_VALUE: 実施内容・確認事項⑨
% #HANDWRITTEN: 0
\fieldvalue{\handwritten{0}} \\
\hline
⑩ &
% #VALUE_ID: LEX-P0066-V0015
% #FIELD_VALUE: 実施内容・確認事項⑩
% #HANDWRITTEN: 0
\fieldvalue{\handwritten{0}} \\
\hline
⑪ &
% #VALUE_ID: LEX-P0066-V0016
% #FIELD_VALUE: 実施内容・確認事項⑪
% #HANDWRITTEN: 0
\fieldvalue{\handwritten{0}} \\
\hline
⑫ &
% #VALUE_ID: LEX-P0066-V0017
% #FIELD_VALUE: 実施内容・確認事項⑫
% #HANDWRITTEN: 0
\fieldvalue{\handwritten{0}} \\
\hline
⑬ &
% #VALUE_ID: LEX-P0066-V0018
% #FIELD_VALUE: 実施内容・確認事項⑬
% #HANDWRITTEN: 01
\fieldvalue{\handwritten{01}} \\
\hline
⑭ &
% #VALUE_ID: LEX-P0066-V0019
% #FIELD_VALUE: 実施内容・確認事項⑭
% #HANDWRITTEN: 0
\fieldvalue{\handwritten{0}} \\
\hline
⑮ &
% #VALUE_ID: LEX-P0066-V0020
% #FIELD_VALUE: 実施内容・確認事項⑮
% #HANDWRITTEN: 0
\fieldvalue{\handwritten{0}} \\
\hline
⑯ &
% #VALUE_ID: LEX-P0066-V0021
% #FIELD_VALUE: 実施内容・確認事項⑯
% #HANDWRITTEN: 0
\fieldvalue{\handwritten{0}} \\
\hline
⑰ &
% #VALUE_ID: LEX-P0066-V0022
% #FIELD_VALUE: 実施内容・確認事項⑰
% #HANDWRITTEN: 0
\fieldvalue{\handwritten{0}} \\
\hline
\end{tabularx}

\vspace{0.5em}

\noindent
% TODO: A red stamped or handwritten entry is visible near the lower edge of the page; its text is not legible.
\begin{flushright}
% #VALUE_ID: LEX-P0066-V0023
% #HANDWRITTEN: 判読困難
\handwritten{判読困難}
\end{flushright}

% TODO: The form continues beyond the visible page boundary.
% LEXOID_PAGE_COMPLETED: 66/78

\newpage

\begin{center}
\textbf{检验报告（22）}
\end{center}

\begin{tabularx}{\textwidth}{|X|X|X|X|}
\hline
名称 & 规格 & 数量 & 送检单位 \\ \hline
\TODO{表格中的具体名称、规格及数量文字较模糊。}
%
% #VALUE_ID: LEX-P0067-V0001
% #FIELD_VALUE: 表格灰色栏内手写值
% #TODO #HANDWRITTEN: （字迹不清）; 蓝色手写内容难以辨认
\fieldvalue{\handwritten{（字迹不清）}} &
\TODO{印刷值不清} &
\TODO{印刷值不清} &
\TODO{印刷值不清} \\ \hline
\end{tabularx}

\vspace{0.5em}

检验报告编号：\TODO{编号文字模糊}

\vspace{0.5em}

\begin{tabularx}{\textwidth}{@{}lX@{}}
检验依据： & ZK-SMP-PP-03-001 \\
           & ZK-SMP-PP-03-02 \\
检验结论： & \TODO{结论文字模糊} \\
\end{tabularx}

\vspace{0.5em}

\begin{tabularx}{\textwidth}{|l|X|}
\hline
产品名称： &
% #VALUE_ID: LEX-P0067-V0002
% #FIELD_VALUE: 产品名称
% #TODO #HANDWRITTEN: （字迹不清）; 蓝色手写产品名称无法辨认
\fieldvalue{\handwritten{（字迹不清）}} \\ \hline
样品编号： &
% #VALUE_ID: LEX-P0067-V0003
% #FIELD_VALUE: 样品编号
% #TODO #HANDWRITTEN: （字迹不清）; 蓝色手写编号无法辨认
\fieldvalue{\handwritten{（字迹不清）}} \\ \hline
检验日期： &
% #VALUE_ID: LEX-P0067-V0004
% #FIELD_VALUE: 检验日期
% #TODO #HANDWRITTEN: （字迹不清）; 蓝色手写日期无法辨认
\fieldvalue{\handwritten{（字迹不清）}} \\ \hline
\end{tabularx}

\vspace{0.5em}

\begin{tabularx}{\textwidth}{|X|l|}
\hline
\multicolumn{1}{|c|}{检验项目} & \multicolumn{1}{c|}{检验结果} \\ \hline
\TODO{检验项目文字模糊} &
% #VALUE_ID: LEX-P0067-V0005
% #FIELD_VALUE: 第一项检验结果
% #HANDWRITTEN: 勾选“合格”
\fieldvalue{\handwritten{合格}} \\ \hline
\TODO{检验项目文字模糊} &
% #VALUE_ID: LEX-P0067-V0006
% #FIELD_VALUE: 第二项检验结果
% #HANDWRITTEN: 勾选“合格”
\fieldvalue{\handwritten{合格}} \\ \hline
\TODO{检验项目文字模糊} &
% #VALUE_ID: LEX-P0067-V0007
% #FIELD_VALUE: 第三项检验结果
% #HANDWRITTEN: 勾选“合格”
\fieldvalue{\handwritten{合格}} \\ \hline
\TODO{检验项目文字模糊} &
% #VALUE_ID: LEX-P0067-V0008
% #FIELD_VALUE: 第四项检验结果
% #HANDWRITTEN: 勾选“合格”
\fieldvalue{\handwritten{合格}} \\ \hline
\TODO{检验项目文字模糊} &
% #VALUE_ID: LEX-P0067-V0009
% #FIELD_VALUE: 第五项检验结果
% #HANDWRITTEN: 勾选“合格”
\fieldvalue{\handwritten{合格}} \\ \hline
\TODO{检验项目文字模糊} &
% #VALUE_ID: LEX-P0067-V0010
% #FIELD_VALUE: 第六项检验结果
% #HANDWRITTEN: 勾选“合格”
\fieldvalue{\handwritten{合格}} \\ \hline
\end{tabularx}

\vspace{0.5em}

\begin{tabularx}{\textwidth}{|X|}
\hline
检验结果及结论： \\
\TODO{本页下方表格中的印刷内容及手写内容均较模糊。} \\ \hline
备注： \\
\TODO{备注栏为空或内容无法辨认。} \\ \hline
\end{tabularx}

% TODO：本页包含倒置扫描及多个难以辨认的蓝色手写字段；已按可见表格结构整理，无法确认的文字保留待核对标记。
% LEXOID_PAGE_COMPLETED: 67/78

\newpage

\section*{医疗器械注册申报资料}

% TODO: 页面顶部的机构标识、文号及印章文字较小且部分倒置，需结合高清原图复核。

\noindent
\textbf{医疗器械注册申请表（续）}

\begin{tabularx}{\textwidth}{|p{0.22\textwidth}|X|}
\hline
项目 & 内容 \\
\hline
申请事项 &
% #VALUE_ID: LEX-P0068-V0001
% #FIELD_VALUE: 申请事项
% #HANDWRITTEN: [手写内容不清]
\fieldvalue{\handwritten{[手写内容不清]}} \\
\hline
申请日期 &
% #VALUE_ID: LEX-P0068-V0002
% #FIELD_VALUE: 申请日期
% #HANDWRITTEN: 2025.06.18
\fieldvalue{\handwritten{2025.06.18}} \\
\hline
联系人及联系方式 &
% #VALUE_ID: LEX-P0068-V0003
% #FIELD_VALUE: 联系人及联系方式
% #HANDWRITTEN: [手写内容不清]
\fieldvalue{\handwritten{[手写内容不清]}} \\
\hline
\end{tabularx}

\subsection*{申请资料及相关情况}

\begin{tabularx}{\textwidth}{|p{0.25\textwidth}|X|p{0.15\textwidth}|}
\hline
项目 & 相关内容 & 是否 \\
\hline
产品名称 &
医疗器械产品名称及相关信息。 &
% #VALUE_ID: LEX-P0068-V0004
% #FIELD_VALUE: 产品名称相关资料是否齐全
% #HANDWRITTEN: 是
\fieldvalue{\handwritten{是}} \\
\hline
产品描述 &
产品的结构组成、工作原理、作用机理及预期用途等。 &
% #VALUE_ID: LEX-P0068-V0005
% #FIELD_VALUE: 产品描述相关资料是否齐全
% #HANDWRITTEN: 是
\fieldvalue{\handwritten{是}} \\
\hline
适用范围 &
产品的适用范围、适用人群及禁忌症等。 &
% #VALUE_ID: LEX-P0068-V0006
% #FIELD_VALUE: 适用范围相关资料是否齐全
% #HANDWRITTEN: 是
\fieldvalue{\handwritten{是}} \\
\hline
产品技术要求 &
产品技术要求及检验报告等相关材料。 &
% #VALUE_ID: LEX-P0068-V0007
% #FIELD_VALUE: 产品技术要求相关资料是否齐全
% #HANDWRITTEN: 是
\fieldvalue{\handwritten{是}} \\
\hline
说明书和标签样稿 &
产品说明书、标签样稿及其符合性说明。 &
% #VALUE_ID: LEX-P0068-V0008
% #FIELD_VALUE: 说明书和标签样稿相关资料是否齐全
% #HANDWRITTEN: 是
\fieldvalue{\handwritten{是}} \\
\hline
\end{tabularx}

\subsection*{其他说明}

\begin{tabularx}{\textwidth}{|p{0.25\textwidth}|X|}
\hline
项目 & 填写内容 \\
\hline
备注 &
% #VALUE_ID: LEX-P0068-V0009
% #FIELD_VALUE: 备注
% #HANDWRITTEN: [手写内容不清]
\fieldvalue{\handwritten{[手写内容不清]}} \\
\hline
注册证号 &
% #VALUE_ID: LEX-P0068-V0010
% #FIELD_VALUE: 注册证号
% #HANDWRITTEN: [数字不清]
\fieldvalue{\handwritten{[数字不清]}} \\
\hline
\end{tabularx}

\subsection*{相关产品信息}

\begin{tabularx}{\textwidth}{|p{0.08\textwidth}|X|p{0.16\textwidth}|p{0.16\textwidth}|p{0.18\textwidth}|p{0.12\textwidth}|}
\hline
序号 & 产品名称 & 型号 & 规格 & 注册证号 & 是否有效 \\
\hline
1 &
相关医疗器械产品 &
% #VALUE_ID: LEX-P0068-V0011
% #FIELD_VALUE: 型号
% #HANDWRITTEN: 110906
\fieldvalue{\handwritten{110906}} &
% #VALUE_ID: LEX-P0068-V0012
% #FIELD_VALUE: 规格
% #HANDWRITTEN: [数字不清]
\fieldvalue{\handwritten{[数字不清]}} &
医疗器械注册证相关信息 &
是 \\
\hline
2 &
相关医疗器械产品 &
% #VALUE_ID: LEX-P0068-V0013
% #FIELD_VALUE: 型号
% #HANDWRITTEN: 120006
\fieldvalue{\handwritten{120006}} &
% #VALUE_ID: LEX-P0068-V0014
% #FIELD_VALUE: 规格
% #HANDWRITTEN: [数字不清]
\fieldvalue{\handwritten{[数字不清]}} &
医疗器械注册证相关信息 &
是 \\
\hline
\end{tabularx}

% TODO: 本页内容在表格边界处疑似与相邻页面连续；仅保留本页可见的表格行。
% LEXOID_PAGE_COMPLETED: 68/78

\newpage

% TODO: The page is scanned upside down and several Chinese labels are low-resolution; the following reproduces the visible form structure and legible entries.

\section*{特种设备相关记录表}

\noindent
\textbf{填表日期：}
% #VALUE_ID: LEX-P0069-V0001
% #FIELD_VALUE: 填表日期
% #HANDWRITTEN: 2017-06-18
\fieldvalue{\handwritten{2017-06-18}}
\hfill
\textbf{验收日期：}
% #VALUE_ID: LEX-P0069-V0002
% #FIELD_VALUE: 验收日期
% #TODO #HANDWRITTEN: 2017-06-?; final digit is unclear
\fieldvalue{\handwritten{2017-06-?}}

\vspace{0.5em}

\begin{tabularx}{\textwidth}{|c|X|X|X|X|}
\hline
\textbf{序号} & \textbf{检查项目} & \textbf{标准值} & \textbf{实测值} & \textbf{备注} \\
\hline
1 & 项目一 & / & / & \\
\hline
2 & 项目二 & / & / & \\
\hline
3 & 项目三 & / & / &
% #VALUE_ID: LEX-P0069-V0003
% #FIELD_VALUE: 项目三实测值
% #TODO #HANDWRITTEN: 30; handwritten numeral is partially unclear
\fieldvalue{\handwritten{30}} \\
\hline
4 & 项目四 & / & / & \\
\hline
5 & 项目五 & / & / &
% #VALUE_ID: LEX-P0069-V0004
% #FIELD_VALUE: 项目五实测值
% #TODO #HANDWRITTEN: 20; handwritten numeral is partially unclear
\fieldvalue{\handwritten{20}} \\
\hline
6 & 项目六 & / & / & \\
\hline
\end{tabularx}

\vspace{1em}

\noindent
\textbf{填表日期：}
% #VALUE_ID: LEX-P0069-V0005
% #FIELD_VALUE: 填表日期
% #TODO #HANDWRITTEN: 2017-06-18; date is faint
\fieldvalue{\handwritten{2017-06-18}}
\hfill
\textbf{验收日期：}
% #VALUE_ID: LEX-P0069-V0006
% #FIELD_VALUE: 验收日期
% #TODO #HANDWRITTEN: 2019-?; date is partly obscured
\fieldvalue{\handwritten{2019-?}}

\begin{tabularx}{\textwidth}{|c|X|X|X|X|}
\hline
\textbf{序号} & \textbf{检查项目} & \textbf{检查结果} & \textbf{备注} & \textbf{签字} \\
\hline
1 & 项目一 & / & / & \\
\hline
2 & 项目二 & / & / &
% #VALUE_ID: LEX-P0069-V0007
% #FIELD_VALUE: 项目二检查结果
% #TODO #HANDWRITTEN: 35; handwritten value is unclear
\fieldvalue{\handwritten{35}} \\
\hline
3 & 项目三 & / & / &
% #VALUE_ID: LEX-P0069-V0008
% #FIELD_VALUE: 项目三检查结果
% #TODO #HANDWRITTEN: 30; handwritten value is unclear
\fieldvalue{\handwritten{30}} \\
\hline
4 & 项目四 & / & / &
% #VALUE_ID: LEX-P0069-V0009
% #FIELD_VALUE: 项目四检查结果
% #TODO #HANDWRITTEN: 20; handwritten value is unclear
\fieldvalue{\handwritten{20}} \\
\hline
5 & 项目五 & / & / &
% #VALUE_ID: LEX-P0069-V0010
% #FIELD_VALUE: 项目五检查结果
% #TODO #HANDWRITTEN: 40; handwritten value is unclear
\fieldvalue{\handwritten{40}} \\
\hline
6 & 项目六 & / & / & \\
\hline
7 & 项目七 & / & / & \\
\hline
8 & 项目八 & / & / &
% #VALUE_ID: LEX-P0069-V0011
% #FIELD_VALUE: 项目八检查结果
% #TODO #HANDWRITTEN: 12000; handwritten number is partly obscured
\fieldvalue{\handwritten{12000}} \\
\hline
9 & 项目九 & / & / & \\
\hline
\end{tabularx}

\vspace{0.5em}

% #VALUE_ID: LEX-P0069-V0012
% #HANDWRITTEN: 蓝色斜线
\handwritten{\text{蓝色斜线标记}}

% TODO: A blue diagonal handwritten line crosses both table regions; its exact intended meaning is not legible.
% LEXOID_PAGE_COMPLETED: 69/78

\newpage

\begin{center}
\textbf{附件1-2}\\
% TODO: 页面顶部标题及表单编号因图像倒置且分辨率有限，部分文字无法可靠辨认。
\end{center}

\begin{tabularx}{\textwidth}{|l|X|l|X|}
\hline
名称 & % TODO: 表头文字部分无法辨认
编号 &
% #VALUE_ID: LEX-P0070-V0001
% #FIELD_VALUE: 编号
% #TODO #HANDWRITTEN: （无法辨认）；蓝色手写内容分辨率不足
\fieldvalue{\handwritten{（无法辨认）}} &
% TODO: 右侧表头及固定内容无法辨认
\\
\hline
\end{tabularx}

\vspace{0.5em}

\noindent
\begin{tabularx}{\textwidth}{|p{0.48\textwidth}|X|}
\hline
\begin{minipage}[t]{0.46\textwidth}
\textbf{调查内容}

\begin{enumerate}
\item
% #VALUE_ID: LEX-P0070-V0002
% #FIELD_VALUE: 第1项选择
% #HANDWRITTEN: ✓
\fieldvalue{\handwritten{\checkmark}}
\quad 相关选项
% TODO: 选项文字无法辨认。

\item
% #VALUE_ID: LEX-P0070-V0003
% #FIELD_VALUE: 第2项选择
% #HANDWRITTEN: ✓
\fieldvalue{\handwritten{\checkmark}}
\quad 相关选项
% TODO: 选项文字无法辨认。

\item
% #VALUE_ID: LEX-P0070-V0004
% #FIELD_VALUE: 第3项选择
% #HANDWRITTEN: ✓
\fieldvalue{\handwritten{\checkmark}}
\quad 相关选项
% TODO: 选项文字无法辨认。

\item
% #VALUE_ID: LEX-P0070-V0005
% #FIELD_VALUE: 第4项选择
% #HANDWRITTEN: ✓
\fieldvalue{\handwritten{\checkmark}}
\quad 相关选项
% TODO: 选项文字无法辨认。

\item
% #VALUE_ID: LEX-P0070-V0006
% #FIELD_VALUE: 第5项选择
% #HANDWRITTEN: ✓
\fieldvalue{\handwritten{\checkmark}}
\quad 相关选项
% TODO: 选项文字无法辨认。
\end{enumerate}

\vspace{1em}

% #VALUE_ID: LEX-P0070-V0007
% #FIELD_VALUE: 选项
% #HANDWRITTEN: ✓
\fieldvalue{\handwritten{\checkmark}}\quad
% TODO: 原表中的选项文字无法辨认。

% #VALUE_ID: LEX-P0070-V0008
% #FIELD_VALUE: 选项
% #HANDWRITTEN: ✓
\fieldvalue{\handwritten{\checkmark}}\quad
% TODO: 原表中的选项文字无法辨认。

% #VALUE_ID: LEX-P0070-V0009
% #FIELD_VALUE: 选项
% #HANDWRITTEN: ✓
\fieldvalue{\handwritten{\checkmark}}\quad
% TODO: 原表中的选项文字无法辨认。

% #VALUE_ID: LEX-P0070-V0010
% #FIELD_VALUE: 选项
% #HANDWRITTEN: ✓
\fieldvalue{\handwritten{\checkmark}}\quad
% TODO: 原表中的选项文字无法辨认。

% #VALUE_ID: LEX-P0070-V0011
% #FIELD_VALUE: 选项
% #HANDWRITTEN: ✓
\fieldvalue{\handwritten{\checkmark}}\quad
% TODO: 原表中的选项文字无法辨认。

% #VALUE_ID: LEX-P0070-V0012
% #FIELD_VALUE: 选项
% #HANDWRITTEN: ✓
\fieldvalue{\handwritten{\checkmark}}\quad
% TODO: 原表中的选项文字无法辨认。

% #VALUE_ID: LEX-P0070-V0013
% #FIELD_VALUE: 选项
% #HANDWRITTEN: ✓
\fieldvalue{\handwritten{\checkmark}}\quad
% TODO: 原表中的选项文字无法辨认。

% #VALUE_ID: LEX-P0070-V0014
% #FIELD_VALUE: 选项
% #HANDWRITTEN: ✓
\fieldvalue{\handwritten{\checkmark}}\quad
% TODO: 原表中的选项文字无法辨认。

% #VALUE_ID: LEX-P0070-V0015
% #FIELD_VALUE: 选项
% #HANDWRITTEN: ✓
\fieldvalue{\handwritten{\checkmark}}\quad
% TODO: 原表中的选项文字无法辨认。

% #VALUE_ID: LEX-P0070-V0016
% #FIELD_VALUE: 选项
% #HANDWRITTEN: ✓
\fieldvalue{\handwritten{\checkmark}}\quad
% TODO: 原表中的选项文字无法辨认。

% #VALUE_ID: LEX-P0070-V0017
% #FIELD_VALUE: 选项
% #HANDWRITTEN: ✓
\fieldvalue{\handwritten{\checkmark}}\quad
% TODO: 原表中的选项文字无法辨认。

% #VALUE_ID: LEX-P0070-V0018
% #FIELD_VALUE: 选项
% #HANDWRITTEN: ✓
\fieldvalue{\handwritten{\checkmark}}\quad
% TODO: 原表中的选项文字无法辨认。

% #VALUE_ID: LEX-P0070-V0019
% #FIELD_VALUE: 选项
% #HANDWRITTEN: ✓
\fieldvalue{\handwritten{\checkmark}}\quad
% TODO: 原表中的选项文字无法辨认。

% #VALUE_ID: LEX-P0070-V0020
% #FIELD_VALUE: 选项
% #HANDWRITTEN: ✓
\fieldvalue{\handwritten{\checkmark}}\quad
% TODO: 原表中的选项文字无法辨认。

% #VALUE_ID: LEX-P0070-V0021
% #FIELD_VALUE: 选项
% #HANDWRITTEN: ✓
\fieldvalue{\handwritten{\checkmark}}\quad
% TODO: 原表中的选项文字无法辨认。

% #VALUE_ID: LEX-P0070-V0022
% #FIELD_VALUE: 选项
% #HANDWRITTEN: ✓
\fieldvalue{\handwritten{\checkmark}}\quad
% TODO: 原表中的选项文字无法辨认。

% #VALUE_ID: LEX-P0070-V0023
% #FIELD_VALUE: 选项
% #HANDWRITTEN: ✓
\fieldvalue{\handwritten{\checkmark}}\quad
% TODO: 原表中的选项文字无法辨认。
\end{minipage}
&
\begin{minipage}[t]{0.48\textwidth}
\textbf{调查说明}

\begin{enumerate}
\item
% TODO: 本项说明文字因图像倒置及分辨率不足无法完整辨认。
调查说明文字。

\item
% TODO: 本项说明文字因图像倒置及分辨率不足无法完整辨认。
调查说明文字。

\item
% TODO: 本项说明文字因图像倒置及分辨率不足无法完整辨认。
调查说明文字。

\item
% TODO: 本项说明文字因图像倒置及分辨率不足无法完整辨认。
调查说明文字。

\item
% TODO: 本项说明文字因图像倒置及分辨率不足无法完整辨认。
调查说明文字。
\end{enumerate}
\end{minipage}
\\
\hline
\end{tabularx}

% TODO: 本页内容在扫描件中呈倒置方向；表格下方及页脚的固定文字仅部分可辨认，未作推断。
% LEXOID_PAGE_COMPLETED: 70/78

\newpage

\begin{center}
\textbf{記入用紙}
\end{center}

\begin{tabularx}{\textwidth}{|l|X|l|}
\hline
番号 & 職種 & 氏名・記入日 \\
\hline
 &  & 
% #VALUE_ID: LEX-P0071-V0001
% #FIELD_VALUE: 記入日
\fieldvalue{令和6年10月1日} \\
\hline
\end{tabularx}

\vspace{0.5em}

\noindent
\begin{tabularx}{\textwidth}{|p{0.18\textwidth}|X|}
\hline
\multicolumn{2}{|c|}{確認事項} \\
\hline
\multicolumn{2}{|l|}{以下の項目について該当する欄に記入すること。} \\
\hline
\end{tabularx}

\vspace{0.5em}

\noindent
\small
\begin{tabularx}{\textwidth}{|c|X|c|}
\hline
\textbf{項目} & \textbf{確認内容} & \textbf{記入欄} \\
\hline
1 & 確認事項 & 
% #VALUE_ID: LEX-P0071-V0002
% #FIELD_VALUE: 記入欄1
% #HANDWRITTEN: C
\fieldvalue{\handwritten{C}} \\
\hline
2 & 確認事項 & 
% #VALUE_ID: LEX-P0071-V0003
% #FIELD_VALUE: 記入欄2
% #HANDWRITTEN: ✓
\fieldvalue{\handwritten{\checkmark}} \\
\hline
3 & 確認事項 & 
% #VALUE_ID: LEX-P0071-V0004
% #FIELD_VALUE: 記入欄3
% #HANDWRITTEN: C
\fieldvalue{\handwritten{C}} \\
\hline
4 & 確認事項 & 
% #VALUE_ID: LEX-P0071-V0005
% #FIELD_VALUE: 記入欄4
% #HANDWRITTEN: ○
\fieldvalue{\handwritten{○}} \\
\hline
5 & 確認事項 & 
% #VALUE_ID: LEX-P0071-V0006
% #FIELD_VALUE: 記入欄5
% #HANDWRITTEN: 2
\fieldvalue{\handwritten{2}} \\
\hline
6 & 確認事項 & 
% #VALUE_ID: LEX-P0071-V0007
% #FIELD_VALUE: 記入欄6
% #HANDWRITTEN: 0
\fieldvalue{\handwritten{0}} \\
\hline
7 & 確認事項 & 
% #VALUE_ID: LEX-P0071-V0008
% #FIELD_VALUE: 記入欄7
% #HANDWRITTEN: ^
\fieldvalue{\handwritten{\textasciicircum}} \\
\hline
8 & 確認事項 & 
% #VALUE_ID: LEX-P0071-V0009
% #FIELD_VALUE: 記入欄8
% #HANDWRITTEN: △
\fieldvalue{\handwritten{△}} \\
\hline
9 & 確認事項 & 
% #VALUE_ID: LEX-P0071-V0010
% #FIELD_VALUE: 記入欄9
% #HANDWRITTEN: C
\fieldvalue{\handwritten{C}} \\
\hline
10 & 確認事項 & 
% #VALUE_ID: LEX-P0071-V0011
% #FIELD_VALUE: 記入欄10
% #HANDWRITTEN: ✓
\fieldvalue{\handwritten{\checkmark}} \\
\hline
11 & 確認事項 & 
% #VALUE_ID: LEX-P0071-V0012
% #FIELD_VALUE: 記入欄11
% #HANDWRITTEN: C
\fieldvalue{\handwritten{C}} \\
\hline
12 & 確認事項 & 
% #VALUE_ID: LEX-P0071-V0013
% #FIELD_VALUE: 記入欄12
% #HANDWRITTEN: ○
\fieldvalue{\handwritten{○}} \\
\hline
13 & 確認事項 & 
% #VALUE_ID: LEX-P0071-V0014
% #FIELD_VALUE: 記入欄13
% #HANDWRITTEN: 2
\fieldvalue{\handwritten{2}} \\
\hline
14 & 確認事項 & 
% #VALUE_ID: LEX-P0071-V0015
% #FIELD_VALUE: 記入欄14
% #HANDWRITTEN: 0
\fieldvalue{\handwritten{0}} \\
\hline
15 & 確認事項 & 
% #VALUE_ID: LEX-P0071-V0016
% #FIELD_VALUE: 記入欄15
% #HANDWRITTEN: ^
\fieldvalue{\handwritten{\textasciicircum}} \\
\hline
16 & 確認事項 & 
% #VALUE_ID: LEX-P0071-V0017
% #FIELD_VALUE: 記入欄16
% #HANDWRITTEN: △
\fieldvalue{\handwritten{△}} \\
\hline
17 & 確認事項 & 
% #VALUE_ID: LEX-P0071-V0018
% #FIELD_VALUE: 記入欄17
% #HANDWRITTEN: C
\fieldvalue{\handwritten{C}} \\
\hline
18 & 確認事項 & 
% #VALUE_ID: LEX-P0071-V0019
% #FIELD_VALUE: 記入欄18
% #HANDWRITTEN: ✓
\fieldvalue{\handwritten{\checkmark}} \\
\hline
19 & 確認事項 & 
% #VALUE_ID: LEX-P0071-V0020
% #FIELD_VALUE: 記入欄19
% #HANDWRITTEN: C
\fieldvalue{\handwritten{C}} \\
\hline
20 & 確認事項 & 
% #VALUE_ID: LEX-P0071-V0021
% #FIELD_VALUE: 記入欄20
% #HANDWRITTEN: ○
\fieldvalue{\handwritten{○}} \\
\hline
21 & 確認事項 & 
% #VALUE_ID: LEX-P0071-V0022
% #FIELD_VALUE: 記入欄21
% #HANDWRITTEN: 2
\fieldvalue{\handwritten{2}} \\
\hline
22 & 確認事項 & 
% #VALUE_ID: LEX-P0071-V0023
% #FIELD_VALUE: 記入欄22
% #HANDWRITTEN: 0
\fieldvalue{\handwritten{0}} \\
\hline
23 & 確認事項 & 
% #VALUE_ID: LEX-P0071-V0024
% #FIELD_VALUE: 記入欄23
% #HANDWRITTEN: C
\fieldvalue{\handwritten{C}} \\
\hline
24 & 確認事項 & 
% #VALUE_ID: LEX-P0071-V0025
% #FIELD_VALUE: 記入欄24
% #HANDWRITTEN: C
\fieldvalue{\handwritten{C}} \\
\hline
25 & 確認事項 & 
% #VALUE_ID: LEX-P0071-V0026
% #FIELD_VALUE: 記入欄25
% #HANDWRITTEN: C
\fieldvalue{\handwritten{C}} \\
\hline
26 & 確認事項 & 
% #VALUE_ID: LEX-P0071-V0027
% #FIELD_VALUE: 記入欄26
% #HANDWRITTEN: C
\fieldvalue{\handwritten{C}} \\
\hline
27 & 確認事項 & 
% #VALUE_ID: LEX-P0071-V0028
% #FIELD_VALUE: 記入欄27
% #HANDWRITTEN: 2
\fieldvalue{\handwritten{2}} \\
\hline
\end{tabularx}
\normalsize

% TODO: The source page is rotated and contains a continuation of the form; some printed Japanese labels are too low-resolution to transcribe reliably.
% LEXOID_PAGE_COMPLETED: 71/78

\newpage

% TODO: The source page is rotated and several Japanese labels are difficult to read; layout and visible entries are represented below.
\begin{center}
\textbf{歯科医師臨床研修（?）実施記録}\\
\end{center}

\noindent
\begin{tabularx}{\textwidth}{|p{0.18\textwidth}|X|p{0.18\textwidth}|}
\hline
施設名 & \multicolumn{2}{l|}{\rule{0.65\textwidth}{0.4pt}}\\
\hline
日付 & \rule{0.25\textwidth}{0.4pt} & 研修期間 \rule{0.20\textwidth}{0.4pt}\\
\hline
\end{tabularx}

\vspace{1em}

\begin{tabularx}{\textwidth}{|p{0.16\textwidth}|X|p{0.20\textwidth}|}
\hline
項目 & 評価内容 & 評価欄\\
\hline
\multicolumn{3}{|l|}{\textbf{研修目標達成度}}\\
\hline
\multicolumn{2}{|l|}{基本的診療能力に関する評価} &
% #VALUE_ID: LEX-P0072-V0001
% #FIELD_VALUE: 基本的診療能力に関する評価
% #HANDWRITTEN: 0
\fieldvalue{\handwritten{0}}\\
\hline
\multicolumn{2}{|l|}{患者への対応に関する評価} &
% #VALUE_ID: LEX-P0072-V0002
% #FIELD_VALUE: 患者への対応に関する評価
% #HANDWRITTEN: 0
\fieldvalue{\handwritten{0}}\\
\hline
\multicolumn{2}{|l|}{診療内容の記録に関する評価} &
% #VALUE_ID: LEX-P0072-V0003
% #FIELD_VALUE: 診療内容の記録に関する評価
% #HANDWRITTEN: 0
\fieldvalue{\handwritten{0}}\\
\hline
\multicolumn{2}{|l|}{医療安全に関する評価} &
% #VALUE_ID: LEX-P0072-V0004
% #FIELD_VALUE: 医療安全に関する評価
% #HANDWRITTEN: 0
\fieldvalue{\handwritten{0}}\\
\hline
\multicolumn{2}{|l|}{感染対策に関する評価} &
% #VALUE_ID: LEX-P0072-V0005
% #FIELD_VALUE: 感染対策に関する評価
% #HANDWRITTEN: 0
\fieldvalue{\handwritten{0}}\\
\hline
\multicolumn{2}{|l|}{診療技術に関する評価} &
% #VALUE_ID: LEX-P0072-V0006
% #FIELD_VALUE: 診療技術に関する評価
% #HANDWRITTEN: 0
\fieldvalue{\handwritten{0}}\\
\hline
\multicolumn{2}{|l|}{コミュニケーション能力に関する評価} &
% #VALUE_ID: LEX-P0072-V0007
% #FIELD_VALUE: コミュニケーション能力に関する評価
% #HANDWRITTEN: 0
\fieldvalue{\handwritten{0}}\\
\hline
\multicolumn{2}{|l|}{チーム医療に関する評価} &
% #VALUE_ID: LEX-P0072-V0008
% #FIELD_VALUE: チーム医療に関する評価
% #HANDWRITTEN: 0
\fieldvalue{\handwritten{0}}\\
\hline
\multicolumn{2}{|l|}{診療管理に関する評価} &
% #VALUE_ID: LEX-P0072-V0009
% #FIELD_VALUE: 診療管理に関する評価
% #HANDWRITTEN: 0
\fieldvalue{\handwritten{0}}\\
\hline
\multicolumn{2}{|l|}{研修態度に関する評価} &
% #VALUE_ID: LEX-P0072-V0010
% #FIELD_VALUE: 研修態度に関する評価
% #HANDWRITTEN: 0
\fieldvalue{\handwritten{0}}\\
\hline
\multicolumn{2}{|l|}{総合評価} &
% #VALUE_ID: LEX-P0072-V0011
% #FIELD_VALUE: 総合評価
% #HANDWRITTEN: 0
\fieldvalue{\handwritten{0}}\\
\hline
\multicolumn{2}{|l|}{特記事項} &
% #VALUE_ID: LEX-P0072-V0012
% #FIELD_VALUE: 特記事項
% #HANDWRITTEN: 0
\fieldvalue{\handwritten{0}}\\
\hline
\multicolumn{2}{|l|}{評価者記入欄} &
% #VALUE_ID: LEX-P0072-V0013
% #FIELD_VALUE: 評価者記入欄
% #HANDWRITTEN: 0
\fieldvalue{\handwritten{0}}\\
\hline
\multicolumn{2}{|l|}{確認欄} &
% #VALUE_ID: LEX-P0072-V0014
% #FIELD_VALUE: 確認欄
% #HANDWRITTEN: 0
\fieldvalue{\handwritten{0}}\\
\hline
\end{tabularx}

\vfill

\noindent
\begin{tabularx}{\textwidth}{|p{0.22\textwidth}|X|p{0.22\textwidth}|}
\hline
確認印 & \rule{0.30\textwidth}{0.4pt} & 日付 \rule{0.14\textwidth}{0.4pt}\\
\hline
記入日 &
% #VALUE_ID: LEX-P0072-V0015
% #FIELD_VALUE: 記入日
% #HANDWRITTEN: 令和7年9月25日（判読困難）
\fieldvalue{\handwritten{令和7年9月25日}} &
\\
\hline
評価日 &
% #VALUE_ID: LEX-P0072-V0016
% #FIELD_VALUE: 評価日
% #HANDWRITTEN: 令和7年9月25日（判読困難）
\fieldvalue{\handwritten{令和7年9月25日}} &
\\
\hline
\end{tabularx}

% TODO: Additional printed labels and the precise wording of the evaluation rows are partially illegible in the supplied scan.
% LEXOID_PAGE_COMPLETED: 72/78

\newpage

\begin{center}
\textbf{表單資料及檢核表}
\end{center}

% TODO: The source scan is rotated and several fixed printed labels are not legible.
\begin{tabularx}{\textwidth}{|p{0.22\textwidth}|X|p{0.20\textwidth}|}
\hline
\textbf{日期} & \textbf{表單名稱／案件資料} & \textbf{編號} \\
\hline
% #VALUE_ID: LEX-P0073-V0001
% #FIELD_VALUE: 日期欄
民國\fieldvalue{112年12月13日} & 
% #VALUE_ID: LEX-P0073-V0002
% #FIELD_VALUE: 案件日期或期間
% #TODO #HANDWRITTEN: 2?--2?; handwritten entry is unclear in the scan
% #HANDWRITTEN: 2?--2?
\fieldvalue{\handwritten{2?--2?}} &
\text{（文字不清）}
\\
\hline
\end{tabularx}

\vspace{1em}

\begin{tabularx}{\textwidth}{|p{0.12\textwidth}|X|p{0.18\textwidth}|}
\hline
\textbf{項次} & \textbf{檢核內容} & \textbf{檢核結果} \\
\hline
1 &
\text{（固定印刷內容辨識不清）} &
% #VALUE_ID: LEX-P0073-V0003
% #FIELD_VALUE: 第1項檢核結果
% #HANDWRITTEN: check mark
\fieldvalue{\handwritten{\checkmark}}
\\
\hline
2 &
\text{（固定印刷內容辨識不清）} &
% #VALUE_ID: LEX-P0073-V0004
% #FIELD_VALUE: 第2項檢核結果
% #HANDWRITTEN: check mark
\fieldvalue{\handwritten{\checkmark}}
\\
\hline
3 &
\text{（固定印刷內容辨識不清）} &
% #VALUE_ID: LEX-P0073-V0005
% #FIELD_VALUE: 第3項檢核結果
% #HANDWRITTEN: check mark
\fieldvalue{\handwritten{\checkmark}}
\\
\hline
4 &
\text{（固定印刷內容辨識不清）} &
% #VALUE_ID: LEX-P0073-V0006
% #FIELD_VALUE: 第4項檢核結果
% #HANDWRITTEN: check mark
\fieldvalue{\handwritten{\checkmark}}
\\
\hline
5 &
\text{（固定印刷內容辨識不清）} &
% #VALUE_ID: LEX-P0073-V0007
% #FIELD_VALUE: 第5項檢核結果
% #HANDWRITTEN: check mark
\fieldvalue{\handwritten{\checkmark}}
\\
\hline
6 &
\text{（固定印刷內容辨識不清）} &
% #VALUE_ID: LEX-P0073-V0008
% #FIELD_VALUE: 第6項檢核結果
% #HANDWRITTEN: check mark
\fieldvalue{\handwritten{\checkmark}}
\\
\hline
7 &
\text{（固定印刷內容辨識不清）} &
% #VALUE_ID: LEX-P0073-V0009
% #FIELD_VALUE: 第7項檢核結果
% #HANDWRITTEN: check mark
\fieldvalue{\handwritten{\checkmark}}
\\
\hline
\end{tabularx}

\vspace{1em}

\begin{tabularx}{\textwidth}{|p{0.15\textwidth}|X|}
\hline
\textbf{備註} & \text{（本頁可見之固定文字及備註內容辨識不清）} \\
\hline
\end{tabularx}

% TODO: This page appears to be a continuation of a larger form; only the visible table and entries on this page are represented.
% LEXOID_PAGE_COMPLETED: 73/78

\newpage

\section*{学校給食実施状況調査票}

\begin{tabularx}{\textwidth}{|p{0.16\textwidth}|X|p{0.16\textwidth}|X|}
\hline
都道府県 & & 市区町村 & \\
\hline
学校名 & & 記入者 & \\
\hline
記入日 &
% #VALUE_ID: LEX-P0074-V0001
% #FIELD_VALUE: 記入日
% #TODO #HANDWRITTEN: R6.10.1; 青字の記入内容は判読が不確実
\fieldvalue{\handwritten{R6.10.1}} &
 & \\
\hline
\end{tabularx}

\vspace{1em}

\noindent
\begin{tabularx}{\textwidth}{|p{0.43\textwidth}|X|}
\hline
\multicolumn{1}{|c|}{調査項目} & \multicolumn{1}{c|}{記入欄} \\
\hline
% #VALUE_ID: LEX-P0074-V0002
% #FIELD_VALUE: 第1項目
% #TODO #HANDWRITTEN: 10; 青字の数字は判読が不確実
第1項目 &
\fieldvalue{\handwritten{10}} \\
\hline
% #VALUE_ID: LEX-P0074-V0003
% #FIELD_VALUE: 第2項目
% #TODO #HANDWRITTEN: 0; 青字の記入は丸印または数字の可能性あり
第2項目 &
\fieldvalue{\handwritten{0}} \\
\hline
% #VALUE_ID: LEX-P0074-V0004
% #FIELD_VALUE: 第3項目
% #TODO #HANDWRITTEN: 0; 青字の記入は丸印または数字の可能性あり
第3項目 &
\fieldvalue{\handwritten{0}} \\
\hline
% #VALUE_ID: LEX-P0074-V0005
% #FIELD_VALUE: 第4項目
% #TODO #HANDWRITTEN: 0; 青字の記入は丸印または数字の可能性あり
第4項目 &
\fieldvalue{\handwritten{0}} \\
\hline
% #VALUE_ID: LEX-P0074-V0006
% #FIELD_VALUE: 第5項目
% #TODO #HANDWRITTEN: 0; 青字の記入は丸印または数字の可能性あり
第5項目 &
\fieldvalue{\handwritten{0}} \\
\hline
% #VALUE_ID: LEX-P0074-V0007
% #FIELD_VALUE: 第6項目
% #TODO #HANDWRITTEN: 0; 青字の記入は丸印または数字の可能性あり
第6項目 &
\fieldvalue{\handwritten{0}} \\
\hline
% #VALUE_ID: LEX-P0074-V0008
% #FIELD_VALUE: 第7項目
% #TODO #HANDWRITTEN: 0; 青字の記入は丸印または数字の可能性あり
第7項目 &
\fieldvalue{\handwritten{0}} \\
\hline
% #VALUE_ID: LEX-P0074-V0009
% #FIELD_VALUE: 第8項目
% #TODO #HANDWRITTEN: 0; 青字の記入は丸印または数字の可能性あり
第8項目 &
\fieldvalue{\handwritten{0}} \\
\hline
% #VALUE_ID: LEX-P0074-V0010
% #FIELD_VALUE: 第9項目
% #TODO #HANDWRITTEN: 0; 青字の記入は丸印または数字の可能性あり
第9項目 &
\fieldvalue{\handwritten{0}} \\
\hline
% #VALUE_ID: LEX-P0074-V0011
% #FIELD_VALUE: 第10項目
% #TODO #HANDWRITTEN: 0; 青字の記入は丸印または数字の可能性あり
第10項目 &
\fieldvalue{\handwritten{0}} \\
\hline
% #VALUE_ID: LEX-P0074-V0012
% #FIELD_VALUE: 第11項目
% #TODO #HANDWRITTEN: 0; 青字の記入は丸印または数字の可能性あり
第11項目 &
\fieldvalue{\handwritten{0}} \\
\hline
% #VALUE_ID: LEX-P0074-V0013
% #FIELD_VALUE: 第12項目
% #TODO #HANDWRITTEN: 0; 青字の記入は丸印または数字の可能性あり
第12項目 &
\fieldvalue{\handwritten{0}} \\
\hline
% #VALUE_ID: LEX-P0074-V0014
% #FIELD_VALUE: 第13項目
% #TODO #HANDWRITTEN: 0; 青字の記入は丸印または数字の可能性あり
第13項目 &
\fieldvalue{\handwritten{0}} \\
\hline
% #VALUE_ID: LEX-P0074-V0015
% #FIELD_VALUE: 第14項目
% #TODO #HANDWRITTEN: 0; 青字の記入は丸印または数字の可能性あり
第14項目 &
\fieldvalue{\handwritten{0}} \\
\hline
% #VALUE_ID: LEX-P0074-V0016
% #FIELD_VALUE: 第15項目
% #TODO #HANDWRITTEN: 0; 青字の記入は丸印または数字の可能性あり
第15項目 &
\fieldvalue{\handwritten{0}} \\
\hline
% #VALUE_ID: LEX-P0074-V0017
% #FIELD_VALUE: 第16項目
% #TODO #HANDWRITTEN: 0; 青字の記入は丸印または数字の可能性あり
第16項目 &
\fieldvalue{\handwritten{0}} \\
\hline
% #VALUE_ID: LEX-P0074-V0018
% #FIELD_VALUE: 第17項目
% #TODO #HANDWRITTEN: 0; 青字の記入は丸印または数字の可能性あり
第17項目 &
\fieldvalue{\handwritten{0}} \\
\hline
% #VALUE_ID: LEX-P0074-V0019
% #FIELD_VALUE: 第18項目
% #TODO #HANDWRITTEN: 0; 青字の記入は丸印または数字の可能性あり
第18項目 &
\fieldvalue{\handwritten{0}} \\
\hline
\end{tabularx}

% TODO: The source page is scanned upside down and the detailed Japanese item labels are not reliably legible. The visible form contains three repeated groups of item rows and a large blank area to the right.
% LEXOID_PAGE_COMPLETED: 74/78

\newpage

\begin{center}
\textbf{（22）関係書類}
\end{center}

\noindent
\begin{tabularx}{\textwidth}{|p{0.22\textwidth}|X|p{0.22\textwidth}|}
\hline
記載日 &
% #VALUE_ID: LEX-P0075-V0001
% #FIELD_VALUE: 記載日
% #TODO #HANDWRITTEN: 8:10:30?; blue handwritten date/time is unclear
\fieldvalue{\handwritten{8:10:30?}} &
記載日\\
\hline
&
% #VALUE_ID: LEX-P0075-V0002
% #FIELD_VALUE: 記載日
% #TODO #HANDWRITTEN: X? 9?92; blue handwritten date is unclear
\fieldvalue{\handwritten{X? 9?92}} &
\\
\hline
\multicolumn{3}{|l|}{
% #VALUE_ID: LEX-P0075-V0003
% #FIELD_VALUE: 上欄チェック
% #HANDWRITTEN: /
\fieldvalue{\handwritten{/}}
}\\
\hline
\end{tabularx}

\vspace{0.5em}

\begin{tabularx}{\textwidth}{|p{0.20\textwidth}|X|p{0.18\textwidth}|}
\hline
\multicolumn{3}{|c|}{項目}\\
\hline
記載欄 &
内容 &
記入値\\
\hline
第1項 &
% #VALUE_ID: LEX-P0075-V0004
% #FIELD_VALUE: 第1項の記入欄
% #HANDWRITTEN: 1
\fieldvalue{\handwritten{1}} &
\\
\hline
第2項 &
% #VALUE_ID: LEX-P0075-V0005
% #FIELD_VALUE: 第2項の記入欄
% #HANDWRITTEN: /
\fieldvalue{\handwritten{/}} &
\\
\hline
第3項 &
% #VALUE_ID: LEX-P0075-V0006
% #FIELD_VALUE: 第3項の記入欄
% #HANDWRITTEN: /
\fieldvalue{\handwritten{/}} &
\\
\hline
第4項 &
% #VALUE_ID: LEX-P0075-V0007
% #FIELD_VALUE: 第4項の記入欄
% #HANDWRITTEN: /
\fieldvalue{\handwritten{/}} &
\\
\hline
第5項 &
% #VALUE_ID: LEX-P0075-V0008
% #FIELD_VALUE: 第5項の記入欄
% #HANDWRITTEN: /
\fieldvalue{\handwritten{/}} &
\\
\hline
第6項 &
% #VALUE_ID: LEX-P0075-V0009
% #FIELD_VALUE: 第6項の記入欄
% #HANDWRITTEN: /
\fieldvalue{\handwritten{/}} &
\\
\hline
第7項 &
% #VALUE_ID: LEX-P0075-V0010
% #FIELD_VALUE: 第7項の記入欄
% #HANDWRITTEN: /
\fieldvalue{\handwritten{/}} &
\\
\hline
第8項 &
% #VALUE_ID: LEX-P0075-V0011
% #FIELD_VALUE: 第8項の記入欄
% #HANDWRITTEN: /
\fieldvalue{\handwritten{/}} &
\\
\hline
第9項 &
% #VALUE_ID: LEX-P0075-V0012
% #FIELD_VALUE: 第9項の記入欄
% #HANDWRITTEN: 2
\fieldvalue{\handwritten{2}} &
\\
\hline
第10項 &
% #VALUE_ID: LEX-P0075-V0013
% #FIELD_VALUE: 第10項の記入欄
% #HANDWRITTEN: 0
\fieldvalue{\handwritten{0}} &
\\
\hline
第11項 &
% #VALUE_ID: LEX-P0075-V0014
% #FIELD_VALUE: 第11項の記入欄
% #HANDWRITTEN: 0
\fieldvalue{\handwritten{0}} &
\\
\hline
第12項 &
% #VALUE_ID: LEX-P0075-V0015
% #FIELD_VALUE: 第12項の記入欄
% #HANDWRITTEN: 0
\fieldvalue{\handwritten{0}} &
\\
\hline
第13項 &
% #VALUE_ID: LEX-P0075-V0016
% #FIELD_VALUE: 第13項の記入欄
% #HANDWRITTEN: 0
\fieldvalue{\handwritten{0}} &
\\
\hline
第14項 &
% #VALUE_ID: LEX-P0075-V0017
% #FIELD_VALUE: 第14項の記入欄
% #HANDWRITTEN: 3
\fieldvalue{\handwritten{3}} &
\\
\hline
第15項 &
% #VALUE_ID: LEX-P0075-V0018
% #FIELD_VALUE: 第15項の記入欄
% #HANDWRITTEN: /
\fieldvalue{\handwritten{/}} &
\\
\hline
第16項 &
% #VALUE_ID: LEX-P0075-V0019
% #FIELD_VALUE: 第16項の記入欄
% #HANDWRITTEN: /
\fieldvalue{\handwritten{/}} &
\\
\hline
第17項 &
% #VALUE_ID: LEX-P0075-V0020
% #FIELD_VALUE: 第17項の記入欄
% #HANDWRITTEN: 0
\fieldvalue{\handwritten{0}} &
\\
\hline
\end{tabularx}

\vspace{0.5em}

\noindent
\textbf{備考}\\
\noindent
\begin{tabularx}{\textwidth}{|X|}
\hline
記載事項および確認欄\\[5em]
\hline
\end{tabularx}

% TODO: The source image is rotated and low-resolution; several Japanese labels and the exact field grouping are illegible.
% LEXOID_PAGE_COMPLETED: 75/78

\newpage

\section*{职业健康检查结果记录表}

% TODO：本页扫描图像为倒置页面，部分标题及字段文字难以准确辨认；以下按可见版式整理。

\begin{tabularx}{\textwidth}{|l|X|l|X|l|X|}
\hline
姓名 & 
% #VALUE_ID: LEX-P0076-V0001
% #FIELD_VALUE: 姓名
% #HANDWRITTEN: （难以辨认）
\fieldvalue{\handwritten{（难以辨认）}} &
性别 & & 年龄 & \\
\hline
检查日期 &
% #VALUE_ID: LEX-P0076-V0002
% #FIELD_VALUE: 检查日期
% #HANDWRITTEN: 2023年（其余数字难以辨认）
\fieldvalue{\handwritten{2023年（其余数字难以辨认）}} &
工种 & & 接触职业病危害因素 & 200--480 \\
\hline
\end{tabularx}

\vspace{0.5em}

\begin{tabularx}{\textwidth}{|p{0.48\textwidth}|X|}
\hline
\textbf{职业健康检查项目及结果} &
\textbf{职业健康检查结果评价} \\
\hline
{\begin{tabularx}{\linewidth}{@{}lXr@{}}
项目一 &
% #VALUE_ID: LEX-P0076-V0003
% #FIELD_VALUE: 项目一选择
% #HANDWRITTEN: √
\fieldvalue{\handwritten{√}} & \\[0.25em]
项目二 &
% #VALUE_ID: LEX-P0076-V0004
% #FIELD_VALUE: 项目二选择
% #HANDWRITTEN: √
\fieldvalue{\handwritten{√}} & \\[0.25em]
项目三 &
% #VALUE_ID: LEX-P0076-V0005
% #FIELD_VALUE: 项目三选择
% #HANDWRITTEN: √
\fieldvalue{\handwritten{√}} & \\[0.25em]
项目四 &
% #VALUE_ID: LEX-P0076-V0006
% #FIELD_VALUE: 项目四选择
% #HANDWRITTEN: √
\fieldvalue{\handwritten{√}} & \\[0.25em]
项目五 &
% #VALUE_ID: LEX-P0076-V0007
% #FIELD_VALUE: 项目五选择
% #HANDWRITTEN: √
\fieldvalue{\handwritten{√}} & \\[0.25em]
项目六 &
% #VALUE_ID: LEX-P0076-V0008
% #FIELD_VALUE: 项目六选择
% #HANDWRITTEN: √
\fieldvalue{\handwritten{√}} & \\[0.25em]
项目七 &
% #VALUE_ID: LEX-P0076-V0009
% #FIELD_VALUE: 项目七选择
% #HANDWRITTEN: √
\fieldvalue{\handwritten{√}} & \\[0.25em]
项目八 &
% #VALUE_ID: LEX-P0076-V0010
% #FIELD_VALUE: 项目八选择
% #HANDWRITTEN: √
\fieldvalue{\handwritten{√}} & \\[0.25em]
项目九 &
% #VALUE_ID: LEX-P0076-V0011
% #FIELD_VALUE: 项目九选择
% #HANDWRITTEN: √
\fieldvalue{\handwritten{√}} & \\[0.25em]
项目十 &
% #VALUE_ID: LEX-P0076-V0012
% #FIELD_VALUE: 项目十选择
% #HANDWRITTEN: √
\fieldvalue{\handwritten{√}} & \\[0.25em]
项目十一 &
% #VALUE_ID: LEX-P0076-V0013
% #FIELD_VALUE: 项目十一选择
% #HANDWRITTEN: √
\fieldvalue{\handwritten{√}} & \\[0.25em]
项目十二 &
% #VALUE_ID: LEX-P0076-V0014
% #FIELD_VALUE: 项目十二选择
% #HANDWRITTEN: √
\fieldvalue{\handwritten{√}} & \\[0.25em]
项目十三 &
% #VALUE_ID: LEX-P0076-V0015
% #FIELD_VALUE: 项目十三选择
% #HANDWRITTEN: √
\fieldvalue{\handwritten{√}} & \\[0.25em]
项目十四 &
% #VALUE_ID: LEX-P0076-V0016
% #FIELD_VALUE: 项目十四选择
% #HANDWRITTEN: √
\fieldvalue{\handwritten{√}} & \\[0.25em]
项目十五 &
% #VALUE_ID: LEX-P0076-V0017
% #FIELD_VALUE: 项目十五选择
% #HANDWRITTEN: √
\fieldvalue{\handwritten{√}} & \\[0.25em]
项目十六 &
% #VALUE_ID: LEX-P0076-V0018
% #FIELD_VALUE: 项目十六选择
% #HANDWRITTEN: √
\fieldvalue{\handwritten{√}} & \\
\end{tabularx}}
&
\begin{minipage}[t]{\linewidth}
\small
\begin{enumerate}
  \item 职业健康检查结果评价。
  \item 根据职业健康检查结果，作出相应的健康评价。
  \item 对检查结果异常者，应结合职业病危害接触史进行分析。
  \item 对需要复查或进一步检查者，应提出处理意见。
  \item 对职业禁忌证及疑似职业病，应按规定进行处理。
  \item 检查结果及处理意见应如实记录并归档。
\end{enumerate}
\end{minipage}
\\
\hline
\end{tabularx}

\vspace{0.5em}

\noindent
\textbf{结论及处理意见：}

\begin{tabularx}{\textwidth}{|X|}
\hline
% TODO：结论栏内的文字及勾选项因扫描质量和页面倒置难以辨认。
\strut\\[2.5em]
\hline
\end{tabularx}

\vfill

\noindent
\footnotesize
第181页\quad 共222页
% LEXOID_PAGE_COMPLETED: 76/78

\newpage

% TODO: The page image is upside down and several printed Chinese characters are illegible; the table structure and handwritten selections are retained below.

\begin{center}
\textbf{（表格标题及表头文字难辨）}
\end{center}

\begin{tabularx}{\textwidth}{|p{0.36\textwidth}|X|}
\hline
\textbf{项目／选项} & \textbf{内容} \\
\hline
\multicolumn{2}{|l|}{\textbf{第一组}} \\
\hline
表内选项 1 &
% #VALUE_ID: LEX-P0077-V0001
% #FIELD_VALUE: 第一组表内选项1
% #TODO #HANDWRITTEN: C; 蓝色手写标记字迹不清
\fieldvalue{\handwritten{C}} \\
\hline
表内选项 2 &
% #VALUE_ID: LEX-P0077-V0002
% #FIELD_VALUE: 第一组表内选项2
% #TODO #HANDWRITTEN: C; 蓝色手写标记字迹不清
\fieldvalue{\handwritten{C}} \\
\hline
表内选项 3 &
% #VALUE_ID: LEX-P0077-V0003
% #FIELD_VALUE: 第一组表内选项3
% #TODO #HANDWRITTEN: C; 蓝色手写标记字迹不清
\fieldvalue{\handwritten{C}} \\
\hline
表内选项 4 &
% #VALUE_ID: LEX-P0077-V0004
% #FIELD_VALUE: 第一组表内选项4
% #TODO #HANDWRITTEN: C; 蓝色手写标记字迹不清
\fieldvalue{\handwritten{C}} \\
\hline
表内选项 5 &
% #VALUE_ID: LEX-P0077-V0005
% #FIELD_VALUE: 第一组表内选项5
% #TODO #HANDWRITTEN: C; 蓝色手写标记字迹不清
\fieldvalue{\handwritten{C}} \qquad
% #VALUE_ID: LEX-P0077-V0006
% #FIELD_VALUE: 第一组表内选项5旁标记
% #TODO #HANDWRITTEN: 方框标记; 蓝色手写标记内容不清
\fieldvalue{\handwritten{方框标记}} \\
\hline
表内选项 6 &
% #VALUE_ID: LEX-P0077-V0007
% #FIELD_VALUE: 第一组表内选项6旁标记
% #TODO #HANDWRITTEN: 方框标记; 蓝色手写标记内容不清
\fieldvalue{\handwritten{方框标记}} \\
\hline
表内选项 7 &
% #VALUE_ID: LEX-P0077-V0008
% #FIELD_VALUE: 第一组表内选项7
% #TODO #HANDWRITTEN: C; 蓝色手写标记字迹不清
\fieldvalue{\handwritten{C}} \\
\hline
\multicolumn{2}{|l|}{\textbf{第二组}} \\
\hline
表内选项 1 &
% #VALUE_ID: LEX-P0077-V0009
% #FIELD_VALUE: 第二组表内选项1
% #TODO #HANDWRITTEN: C; 蓝色手写标记字迹不清
\fieldvalue{\handwritten{C}} \\
\hline
表内选项 2 &
% #VALUE_ID: LEX-P0077-V0010
% #FIELD_VALUE: 第二组表内选项2
% #TODO #HANDWRITTEN: C; 蓝色手写标记字迹不清
\fieldvalue{\handwritten{C}} \\
\hline
表内选项 3 &
% #VALUE_ID: LEX-P0077-V0011
% #FIELD_VALUE: 第二组表内选项3
% #TODO #HANDWRITTEN: C; 蓝色手写标记字迹不清
\fieldvalue{\handwritten{C}} \\
\hline
表内选项 4 &
% #VALUE_ID: LEX-P0077-V0012
% #FIELD_VALUE: 第二组表内选项4
% #TODO #HANDWRITTEN: C; 蓝色手写标记字迹不清
\fieldvalue{\handwritten{C}} \\
\hline
表内选项 5 &
% #VALUE_ID: LEX-P0077-V0013
% #FIELD_VALUE: 第二组表内选项5
% #TODO #HANDWRITTEN: C; 蓝色手写标记字迹不清
\fieldvalue{\handwritten{C}} \qquad
% #VALUE_ID: LEX-P0077-V0014
% #FIELD_VALUE: 第二组表内选项5旁标记
% #TODO #HANDWRITTEN: 方框标记; 蓝色手写标记内容不清
\fieldvalue{\handwritten{方框标记}} \\
\hline
表内选项 6 &
% #VALUE_ID: LEX-P0077-V0015
% #FIELD_VALUE: 第二组表内选项6旁标记
% #TODO #HANDWRITTEN: 方框标记; 蓝色手写标记内容不清
\fieldvalue{\handwritten{方框标记}} \\
\hline
表内选项 7 &
% #VALUE_ID: LEX-P0077-V0016
% #FIELD_VALUE: 第二组表内选项7
% #TODO #HANDWRITTEN: C; 蓝色手写标记字迹不清
\fieldvalue{\handwritten{C}} \\
\hline
\multicolumn{2}{|l|}{\textbf{第三组}} \\
\hline
表内选项 1 &
% #VALUE_ID: LEX-P0077-V0017
% #FIELD_VALUE: 第三组表内选项1
% #TODO #HANDWRITTEN: C; 蓝色手写标记字迹不清
\fieldvalue{\handwritten{C}} \\
\hline
表内选项 2 &
% #VALUE_ID: LEX-P0077-V0018
% #FIELD_VALUE: 第三组表内选项2
% #TODO #HANDWRITTEN: C; 蓝色手写标记字迹不清
\fieldvalue{\handwritten{C}} \\
\hline
表内选项 3 &
% #VALUE_ID: LEX-P0077-V0019
% #FIELD_VALUE: 第三组表内选项3
% #TODO #HANDWRITTEN: C; 蓝色手写标记字迹不清
\fieldvalue{\handwritten{C}} \\
\hline
表内选项 4 &
% #VALUE_ID: LEX-P0077-V0020
% #FIELD_VALUE: 第三组表内选项4
% #TODO #HANDWRITTEN: C; 蓝色手写标记字迹不清
\fieldvalue{\handwritten{C}} \\
\hline
表内选项 5 &
% #VALUE_ID: LEX-P0077-V0021
% #FIELD_VALUE: 第三组表内选项5
% #TODO #HANDWRITTEN: C; 蓝色手写标记字迹不清
\fieldvalue{\handwritten{C}} \\
\hline
表内选项 6 &
% #VALUE_ID: LEX-P0077-V0022
% #FIELD_VALUE: 第三组表内选项6
% #TODO #HANDWRITTEN: C; 蓝色手写标记字迹不清
\fieldvalue{\handwritten{C}} \\
\hline
\end{tabularx}

% TODO: The large blank area on the right side of the visible table contains no discernible entries.
% LEXOID_PAGE_COMPLETED: 77/78

\newpage

\section*{放射性医薬品製造記録（2）}

\begin{center}
\begin{tabular}{|p{0.18\linewidth}|p{0.18\linewidth}|p{0.18\linewidth}|p{0.18\linewidth}|}
\hline
製品名 & 製造番号 & 製造年月日 & 200--280\\
\hline
\end{tabular}
\end{center}

% #VALUE_ID: LEX-P0078-V0001
% #FIELD_VALUE: 製造年月日（左欄）
% #TODO #HANDWRITTEN: 8/?? 28?; blue handwritten entry is partially illegible
\fieldvalue{\handwritten{8/?? 28?}}
\qquad
% #VALUE_ID: LEX-P0078-V0002
% #FIELD_VALUE: 製造年月日（右欄）
% #TODO #HANDWRITTEN: 24? ...; blue handwritten entry is partially illegible
\fieldvalue{\handwritten{24? ...}}

\vspace{0.5em}

% #VALUE_ID: LEX-P0078-V0003
% #FIELD_VALUE: 品名・規格等
% #TODO #HANDWRITTEN: 9? ...; blue handwritten text is partially illegible
\fieldvalue{\handwritten{9? ...}}

\begin{center}
\begin{tabular}{|p{0.62\linewidth}|p{0.26\linewidth}|}
\hline
\multicolumn{2}{|c|}{製造工程・確認事項}\\
\hline
\textbf{確認項目} & \textbf{記録}\\
\hline
原料の確認・秤量 & 
% #VALUE_ID: LEX-P0078-V0004
% #FIELD_VALUE: 原料の確認・秤量
% #HANDWRITTEN: C
\fieldvalue{\handwritten{C}}\\
\hline
調製条件の確認 & 
% #VALUE_ID: LEX-P0078-V0005
% #FIELD_VALUE: 調製条件の確認
% #HANDWRITTEN: C
\fieldvalue{\handwritten{C}}\\
\hline
外観の確認 & 
% #VALUE_ID: LEX-P0078-V0006
% #FIELD_VALUE: 外観の確認
% #HANDWRITTEN: 0
\fieldvalue{\handwritten{0}}\\
\hline
放射能の確認 & 
% #VALUE_ID: LEX-P0078-V0007
% #FIELD_VALUE: 放射能の確認
% #HANDWRITTEN: 0
\fieldvalue{\handwritten{0}}\\
\hline
容器・表示の確認 & 
% #VALUE_ID: LEX-P0078-V0008
% #FIELD_VALUE: 容器・表示の確認
% #HANDWRITTEN: 0
\fieldvalue{\handwritten{0}}\\
\hline
作業室・設備の確認 & 
% #VALUE_ID: LEX-P0078-V0009
% #FIELD_VALUE: 作業室・設備の確認
% #HANDWRITTEN: 01
\fieldvalue{\handwritten{01}}\\
\hline
記録の確認 & \\
\hline
確認者欄 & \\
\hline
\end{tabular}
\end{center}

% TODO: The printed checklist text in the following continuation is faint and partially illegible.
\begin{center}
\begin{tabular}{|p{0.62\linewidth}|p{0.26\linewidth}|}
\hline
原料の確認・秤量 & 
% #VALUE_ID: LEX-P0078-V0010
% #FIELD_VALUE: 原料の確認・秤量（第2欄）
% #HANDWRITTEN: 0
\fieldvalue{\handwritten{0}}\\
\hline
調製条件の確認 & 
% #VALUE_ID: LEX-P0078-V0011
% #FIELD_VALUE: 調製条件の確認（第2欄）
% #HANDWRITTEN: 0
\fieldvalue{\handwritten{0}}\\
\hline
外観の確認 & 
% #VALUE_ID: LEX-P0078-V0012
% #FIELD_VALUE: 外観の確認（第2欄）
% #HANDWRITTEN: 0
\fieldvalue{\handwritten{0}}\\
\hline
放射能の確認 & 
% #VALUE_ID: LEX-P0078-V0013
% #FIELD_VALUE: 放射能の確認（第2欄）
% #HANDWRITTEN: 0
\fieldvalue{\handwritten{0}}\\
\hline
容器・表示の確認 & 
% #VALUE_ID: LEX-P0078-V0014
% #FIELD_VALUE: 容器・表示の確認（第2欄）
% #HANDWRITTEN: 0
\fieldvalue{\handwritten{0}}\\
\hline
作業室・設備の確認 & 
% #VALUE_ID: LEX-P0078-V0015
% #FIELD_VALUE: 作業室・設備の確認（第2欄）
% #HANDWRITTEN: 01
\fieldvalue{\handwritten{01}}\\
\hline
記録の確認 & \\
\hline
確認者欄 & \\
\hline
最終確認 & 
% #VALUE_ID: LEX-P0078-V0016
% #FIELD_VALUE: 最終確認
% #HANDWRITTEN: 0
\fieldvalue{\handwritten{0}}\\
\hline
\end{tabular}
\end{center}

\noindent
\textbf{備考}\\
\vspace{2em}

\noindent
\textit{（以下、次ページへ続く）}

\end{document}
% LEXOID_PAGE_COMPLETED: 78/78
